\documentclass{article}

\usepackage[utf8]{inputenc}
\usepackage[T1]{fontenc}
\usepackage{helvet}
\usepackage{geometry}
\usepackage{xcolor}
\usepackage{eso-pic}
\usepackage{fancyhdr}
\usepackage{parskip}
\usepackage{microtype}
\usepackage[english, ukrainian]{babel}
\usepackage{url}
\usepackage{amsmath}
\usepackage{amssymb}
\usepackage{amsthm}
\usepackage{longtable}
\usepackage{booktabs}
\usepackage{hyperref}
\usepackage{titlesec}
\usepackage{tocloft}

\renewcommand{\cftsecfont}{\small\bfseries}
\renewcommand{\cftsubsecfont}{\footnotesize}
\renewcommand{\cftsubsubsecfont}{\scriptsize}
\renewcommand{\cftsecpagefont}{\small\bfseries}
\renewcommand{\cftsubsecpagefont}{\footnotesize}
\renewcommand{\cftsubsubsecpagefont}{\scriptsize}

\definecolor{nato}{HTML}{293757}
\definecolor{footergray}{gray}{0.65}

\geometry{a4paper,left=3cm,right=3cm,top=3cm,bottom=3cm}
\renewcommand{\familydefault}{\sfdefault}
\setcounter{secnumdepth}{3}

\titleformat{\section}
  {\color{nato}\Huge\bfseries}{\thesection.}{1em}{}
\titlespacing*{\section}{0pt}{30pt}{12pt}

\titleformat{\subsection}
  {\color{nato}\LARGE\bfseries}{\thesubsection.}{1em}{}
\titlespacing*{\subsection}{0pt}{20pt}{8pt}

\titleformat{\subsubsection}
  {\color{nato}\Large\bfseries}{\thesubsubsection.}{1em}{}
\titlespacing*{\subsubsection}{0pt}{10pt}{6pt}

\fancyhf{}
\fancyhead[R]{\small\color{footergray} 2 червня 2026}
\fancyfoot[L]{\small\color{footergray} Комплексна система захисту інформації. Профіль базових заходів із захисту інформації (155).}
\fancyfoot[R]{\small\color{footergray}\thepage}

\newcommand\HUGE[1]{\fontsize{40}{40}\selectfont #1}

\renewcommand{\headrulewidth}{0pt}
\renewcommand{\footrulewidth}{0.4pt}
\renewcommand{\footrule}{\color{footergray}\hrule width \textwidth height 0.4pt}

\begin{document}

\pagestyle{fancy}

\begin{titlepage}
  \AddToShipoutPictureBG*{\AtPageLowerLeft{\color{nato}\rule{\paperwidth}{\paperheight}}}
  \color{white}\thispagestyle{empty}\vspace*{4cm}\centering
  {\Huge  \textbf{Комплексна система захисту інформації} \par} \vspace{2cm}
  {\HUGE  \textbf{Профіль базових заходів із захисту інформації (155)} \par} \vspace{2.5cm}
  \vfill
  {\large 2026 \copyright\ Інформаційні Судові Системи \par}
\end{titlepage}

\newpage

\begin{center}
  {\Huge\color{nato}\bfseries Зміст\par}
\end{center}
\vspace{40pt}
\tableofcontents

\begin{abstract}
Базовий профіль — затверджується Адміністрацією Держспецзв’язку (наказом).
\end{abstract}

\section{AC}
Клас заходів захисту AC --- УПРАВЛІННЯ ДОСТУПОМ

Цей клас зосереджується на обмеженні доступу до інформаційних систем, ресурсів та функцій лише для авторизованих користувачів, програм і пристроїв.

\textbf{Перелік заходів захисту:} Політика та процедури управління доступом (AC-1); Управління обліковими записами (AC-2); Деактивація облікових записів (AC-2(3)); Вихід із системи за відсутності активності (AC-2(5)); Деактивація облікових записів осіб з високим рівнем ризику (AC-2(13)); Забезпечення доступу (AC-3); Управління інформаційними потоками (AC-4); Розмежування обов'язків (AC-5); Мінімізація повноважень (AC-6); Авторизований доступ до функцій безпеки (AC-6(1)); Непривілейований доступ до незахищених функцій (AC-6(2)); Привілейовані облікові записи (AC-6(5)); Перегляд повноважень користувача (AC-6(7)); Аудит використання привілейованих функцій (AC-6(9)); Заборона непривілейованим користувачам виконувати привілейовані функції (AC-6(10)); Невдалі спроби входу в систему (AC-7); Попередження про використання системи (AC-8); Блокування пристрою (AC-11); Приховані дисплеї (AC-11(1)); Припинення сеансу (AC-12); Віддалений доступ (AC-17); Керовані точки контролю доступу (AC-17(3)); Привілейовані команди та доступ (AC-17(4)); Бездротовий доступ (AC-18); Автентифікація та шифрування (AC-18(1)); Відключення бездротової мережі (AC-18(3)); Контроль доступу для мобільних пристроїв (AC-19); Повне шифрування пристроїв та сховищ інформації (AC-19(5)); Використання зовнішніх систем (AC-20); Обмеження на авторизоване використання (AC-20(1)); Переносні пристрої зберігання даних (AC-20(2)); Публічно доступний контент (AC-22).

\subsection{ПОЛІТИКА ТА ПРОЦЕДУРИ УПРАВЛІННЯ ДОСТУПОМ (AC-1)}
a. Розробити, задокументувати та поширити [Призначення: серед визначеного організацією персоналу або ролей]:\\ 
1. 2. [Вибір (один або декілька): Рівень організації; Рівень місії/бізнес-процесу; рівень системи] політики контролю доступу, яка:\\ 
(a) містить мету, сферу застосування, ролі, відповідальність, зобов'язання керівництва, координацію між організаційними підрозділами та систему контролю відповідності (compliances);\\ 
(b) відповідає чинному законодавству, нормативним документам, директивам, нормам, політикам, стандартам і керівним документам. Процедури, що сприяють реалізації політики управління доступом і відповідних заходів управління доступом.\\ 
b. Призначити на посаду [Призначення: визначену організацією посадову особу] для управління, документування і розповсюдження політики та процедур контролю доступом.\\ 
c. Переглянути та оновити:\\ 
1. поточну політику управління доступом [Призначення: з визначеною організацією частотою] та [Призначення: події, визначені організацією];\\ 
2. поточні процедури управління доступом [Призначення: з визначеною організацією частотою] та [Завдання: події, визначені організацією].

\vspace{10pt}
{\footnotesize\bfseries
No: 1\\
Name: ac\_1\_odp\_01\\
Type: list\\
Default: [\textquotedbl{}admin\textquotedbl{}, \textquotedbl{}security\_officer\textquotedbl{}]
}

{\footnotesize Визначено персонал або ролі, на які поширюється політика контролю доступу}


{\footnotesize\bfseries
No: 2\\
Name: ac\_1\_odp\_02\\
Type: list\\
Default: [\textquotedbl{}admin\textquotedbl{}, \textquotedbl{}security\_officer\textquotedbl{}]
}

{\footnotesize Визначено персонал або ролі, на які поширюються процедури контролю доступу}


{\footnotesize\bfseries
No: 3\\
Name: ac\_1\_odp\_03\\
Type: string\\
Default: nil
}

{\footnotesize Вибрано одне або декілька з наступних ЗНАЧЕНЬ ПАРАМЕТРІВ: \{рівень організації; рівень місії/бізнес-процесу; рівень системи\}}


{\footnotesize\bfseries
No: 4\\
Name: ac\_1\_odp\_04\\
Type: list\\
Default: [\textquotedbl{}admin\textquotedbl{}, \textquotedbl{}security\_officer\textquotedbl{}]
}

{\footnotesize Визначено посадову особу, яка керуватиме політикою та процедурами контролю доступу}


{\footnotesize\bfseries
No: 5\\
Name: ac\_1\_odp\_05\\
Type: list\\
Default: [\textquotedbl{}default\_deny\_rule\textquotedbl{}, \textquotedbl{}abac\_rule\_1\textquotedbl{}]
}

{\footnotesize Визначено частоту, з якою переглядається та оновлюється поточна політика контролю доступу}


{\footnotesize\bfseries
No: 6\\
Name: ac\_1\_odp\_06\\
Type: list\\
Default: [\textquotedbl{}default\_deny\_rule\textquotedbl{}, \textquotedbl{}abac\_rule\_1\textquotedbl{}]
}

{\footnotesize Визначено події, які вимагають перегляду та оновлення поточної політики контролю доступу}


{\footnotesize\bfseries
No: 7\\
Name: ac\_1\_odp\_07\\
Type: integer\\
Default: 30
}

{\footnotesize Визначено частоту, з якою переглядаються та оновлюються поточні процедури контролю доступу}




\subsection{УПРАВЛІННЯ ОБЛІКОВИМИ ЗАПИСАМИ (AC-2)}
a. Визначити та задокументувати типи облікових записів системи, дозволених для використання в ІС для підтримки цілей, завдань, функцій і процесів організації.\\ 
b. Призначити менеджерів облікових записів для управління системними обліковими записами.\\ 
c. Створити умови для групового та рольового членства.\\ 
d. Визначити авторизованих користувачів інформаційної системи, членство в групі та ролі, а також дозволи доступу (наприклад, привілеї) та інші атрибути (за потреби) для кожного облікового запису.\\ 
e. Вимагати схвалення [Призначення: визначеною організацією відповідальною особою або роллю] запитів на створення облікових записів системи.\\ 
f. Створювати, активувати, змінювати, деактивувати та видаляти системні облікові записи відповідно до [Призначення: визначених організацією політики, процедур та умов].\\ 
g. Впровадити моніторинг використання облікових записів системи.\\ 
h. Повідомляти адміністраторів облікових записів у межах [Призначення: визначеного організацією часового періоду для кожної ситуації]:\\ 
1. коли облікові записи більше не потрібні;\\ 
2. коли користувачі звільнені чи переведені;\\ 
3. коли використовуються індивідуальні системи або наявні зміни, які потребують нових знань.\\ 
i. Авторизувати доступ до системи на основі:\\ 
1. Дійсної авторизації доступу.\\ 
2. Передбачуваного використання системи.\\ 
3. Інших атрибутів, що вимагаються організацією.\\ 
j. Проводити перегляд облікових записів на відповідність вимогам управління обліковими записами з [Призначення: визначеною організацією частотою].\\ 
k. Впровадити процес повторного випуску облікових даних спільного/групового облікового запису (якщо він буде розгорнутий), коли особи виходять з групи.\\ 
l. Узгодити процеси управління обліковими записами з процесами звільнення та переводу (передачі повноважень) персоналу.

\vspace{10pt}
{\footnotesize\bfseries
No: 1\\
Name: ac\_2\_odp\_09\\
Type: integer\\
Default: 24
}

{\footnotesize Визначено передумови та критерії членства в групах і ролях; визначено атрибути (за необхідності) для кожного облікового запису; визначено персонал або ролі, необхідні для затвердження запитів на створення облікових записів; визначено політику, процедури, передумови та критерії створення, активації, зміни, деактивації та видалення облікових записів; визначено персонал або ролі, які мають бути повідомлені; визначено період часу, протягом якого адміністратори облікових записів повинні бути повідомлені про те, що облікові записи більше не потрібні; визначено термін, протягом якого необхідно повідомляти адміністраторів облікових записів про звільнення або переведення користувачів; визначено період часу, протягом якого необхідно повідомляти адміністраторів облікових записів про зміни у використанні системи або необхідність знати про зміни для окремої особи; визначено атрибути, необхідні для авторизації доступу до системи (за потреби); AC-02\_ODP[10] AC-02a.[01] AC-02a.[02] AC-02b AC-02с AC-02d.01 AC-02d.02 AC-02d.03[01] AC-02d.03[02] AC-02e AC-02f.[01] AC-02f.[02] AC-02f.[03] AC-02f.[04] AC-02f.[05] AC-02g AC-02h.01 AC-02h.02 AC-02h.03 AC-02i.01 AC-02i.02 AC-02i.03 визначено періодичність перегляду облікових записів; визначено та задокументовано типи облікових записів, дозволених для використання в системі; визначено та задокументовано типи облікових записів, які заборонено використовувати в системі; призначені менеджери облікових записів; необхідні умови та критерії для членства в групах та ролях; визначено авторизованих користувачів системи; вказано приналежність до групи або ролі; для кожного облікового запису вказуються повноваження доступу (тобто привілеї); атрибути (за необхідності) вказуються для кожного облікового запису; для запитів на створення облікових записів потрібні схвалення від персоналу або ролей; облікові записи створюються відповідно до політики, процедур, передумов та критеріїв; облікові записи активуються відповідно до політики, процедур, передумов та критеріїв; облікові записи змінюються відповідно до політики, процедур, передумов та критеріїв; облікові записи деактивуються відповідно до політики, процедур, передумов та критеріїв; облікові записи видаляються відповідно до політики, процедур, передумов та критеріїв; контролюється використання облікових записів; адміністратори облікових записів та персонал або ролі отримують повідомлення протягом періоду часу, коли облікові записи більше не потрібні; адміністратори облікових записів та персонал або ролі отримують повідомлення протягом періоду часу, коли користувачі звільнені чи переведені; адміністратори облікових записів та персонал або ролі отримують повідомлення протягом періоду часу, коли використовуються індивідуальні системи або наявні зміни, які потребують нових знань. доступ до системи здійснюється на підставі дійсної авторизації доступу; доступ до системи авторизується на основі передбачуваного використання системи; доступ до системи авторизовано на основі атрибутів (за необхідності); AC-02j AC-02k.[01] AC-02k.[02] AC-02l.[01] AC-02l.[02] облікові записи переглядаються на відповідність вимогам управління обліковими записами частота; створено процес повторного випуску облікових даних спільного доступу або групових облікових записів (якщо вони розгорнуті), коли користувачів вилучено з групи; впроваджено процес повторного випуску облікових даних спільного доступу або групових облікових записів (якщо вони розгорнуті), коли користувачів вилучено з групи; процеси управління обліковими записами узгоджуються з процесами звільнення персоналу; процеси управління обліковими записами узгоджуються з процесами переводу персоналу}




\subsubsection{ДЕАКТИВАЦІЯ ОБЛІКОВИХ ЗАПИСІВ (AC-2(3))}
Автоматично деактивувати облікові записи коли: a) їх строк дії минув; b) вони більше не пов'язані з користувачем; c) вони порушують організаційну політику; d) вони були неактивними впродовж [Призначення: визначеного організацією періоду часу].

\vspace{10pt}
{\footnotesize\bfseries
No: 1\\
Name: ac\_2\_3\_b\\
Type: list\\
Default: [\textquotedbl{}admin\textquotedbl{}, \textquotedbl{}security\_officer\textquotedbl{}]
}

{\footnotesize Облікові записи деактивуються протягом часового періоду, коли облікові записи більше не пов'язані з користувачем або фізичною особою}


{\footnotesize\bfseries
No: 2\\
Name: ac\_2\_3\_d\\
Type: list\\
Default: [\textquotedbl{}default\_deny\_rule\textquotedbl{}, \textquotedbl{}abac\_rule\_1\textquotedbl{}]
}

{\footnotesize Облікові записи відключено протягом часового періоду, коли облікові записи порушують політику організації; облікові записи деактивуються протягом , якщо вони були неактивними протягом }


{\footnotesize\bfseries
No: 3\\
Name: ac\_2\_3\_odp\_01\\
Type: integer\\
Default: 30
}

{\footnotesize Визначено період часу, протягом якого необхідно деактивувати облікові записи}


{\footnotesize\bfseries
No: 4\\
Name: ac\_2\_3\_odp\_02\\
Type: list\\
Default: [\textquotedbl{}login\textquotedbl{}, \textquotedbl{}logout\textquotedbl{}, \textquotedbl{}failed\_attempt\textquotedbl{}]
}

{\footnotesize Визначено період часу неактивності, після закінчення якого облікові записи будуть деактивовані; AC-02(03)(a) облікові записи деактивуються протягом часового періоду, коли термін дії облікових записів минув}




\subsubsection{ВИХІД ІЗ СИСТЕМИ ЗА ВІДСУТНОСТІ АКТИВНОСТІ (AC-2(5))}
Вимагати від користувачів виходити із системи, коли [Призначення: вичерпано визначений організацією періоду часу очікування або опис того, коли необхідно вийти із системи].

\vspace{10pt}
{\footnotesize\bfseries
No: 1\\
Name: ac\_2\_5\_01\\
Type: integer\\
Default: 30
}

{\footnotesize Користувачі повинні виходити з системи, коли період очікуваної бездіяльності або опис часу, коли потрібно вийти з системи}


{\footnotesize\bfseries
No: 2\\
Name: ac\_2\_5\_odp\\
Type: integer\\
Default: 30
}

{\footnotesize Визначено часовий період очікуваної бездіяльності або опис, коли потрібно вийти з системи}




\subsubsection{ДЕАКТИВАЦІЯ ОБЛІКОВИХ ЗАПИСІВ ОСІБ З ВИСОКИМ РІВНЕМ РИЗИКУ (AC-2(13))}
Деактивувати облікові записи користувачів, які становлять значний ризик, у межах [Призначення: визначеного організацією періоду часу] після виявлення ризику.

\vspace{10pt}
{\footnotesize\bfseries
No: 1\\
Name: ac\_2\_13\_01\\
Type: integer\\
Default: 30
}

{\footnotesize Облікові записи користувачів деактивуються протягом періоду часу з моменту виявлення значних ризиків}


{\footnotesize\bfseries
No: 2\\
Name: ac\_2\_13\_odp\_01\\
Type: integer\\
Default: 30
}

{\footnotesize Визначено період часу, протягом якого необхідно деактивувати облікові записи фізичних осіб, які становлять значний ризик}


{\footnotesize\bfseries
No: 3\\
Name: ac\_2\_13\_odp\_02\\
Type: string\\
Default: nil
}

{\footnotesize Визначено значні ризики, що призводять до деактивації облікових записів}




\subsection{ЗАБЕЗПЕЧЕННЯ ДОСТУПУ (AC-3)}
Застосовувати затверджені повноваження для логічного доступу до інформації та ресурсів системи відповідно до чинної політики (правил) управління доступом.

\vspace{10pt}
{\footnotesize\bfseries
No: 1\\
Name: ac\_3\_01\\
Type: list\\
Default: [\textquotedbl{}default\_deny\_rule\textquotedbl{}, \textquotedbl{}abac\_rule\_1\textquotedbl{}]
}

{\footnotesize Затверджені повноваження на логічний доступ до інформації та ресурсів системи виконуються відповідно до чинних політик(правил) управління доступом}




\subsection{УПРАВЛІННЯ ІНФОРМАЦІЙНИМИ ПОТОКАМИ (AC-4)}
Застосувати затверджені повноваження для управління потоком інформації всередині системи та між пов'язаними системами на основі [Призначення: визначеними організацією політиками управління інформаційним потоком].

\vspace{10pt}
{\footnotesize\bfseries
No: 1\\
Name: ac\_4\_01\\
Type: list\\
Default: [\textquotedbl{}default\_deny\_rule\textquotedbl{}, \textquotedbl{}abac\_rule\_1\textquotedbl{}]
}

{\footnotesize Затверджені повноваження застосовуються для контролю потоку інформації всередині системи та між підключеними системами на основі політики управління інформаційними потоками}


{\footnotesize\bfseries
No: 2\\
Name: ac\_4\_odp\\
Type: list\\
Default: [\textquotedbl{}default\_deny\_rule\textquotedbl{}, \textquotedbl{}abac\_rule\_1\textquotedbl{}]
}

{\footnotesize Визначено політики управління інформаційними потоками всередині системи та між підключеними системами}




\subsection{РОЗМЕЖУВАННЯ ОБОВ'ЯЗКІВ (AC-5)}
a. Розмежувати і документувати [Призначення: визначені організацією обов'язки окремих осіб].\\ 
b. Установити правила авторизації доступу для підтримки розмежування обов'язків.

\vspace{10pt}
{\footnotesize\bfseries
No: 1\\
Name: ac\_5\_odp\\
Type: string\\
Default: nil
}

{\footnotesize Визначено обов'язки осіб, які потребують розмежування; AC-05[a] обов'язки осіб визначені та задокументовані; AC-05[b] визначено права авторизації доступу до системи для підтримки розмежування обов'язків}




\subsection{МІНІМІЗАЦІЯ ПОВНОВАЖЕНЬ (AC-6)}
Впровадити принцип мінімізації повноважень, який дозволяє користувачам (або процесам, що діють від імені користувачів) здійснювати лише такі авторизовані звернення, які необхідні для виконання визначених завдань відповідно до цілей (призначення, місії) організації та функцій.

\vspace{10pt}
{\footnotesize\bfseries
No: 1\\
Name: ac\_6\_01\\
Type: string\\
Default: nil
}

{\footnotesize Застосовується принцип мінімалізації повноважень, який дозволяє користувачам (або процесам, що діють від імені користувачів) здійснювати лише такі авторизовані звернення, які необхідні для виконання поставлених завдань організації}




\subsubsection{АВТОРИЗОВАНИЙ ДОСТУП ДО ФУНКЦІЙ БЕЗПЕКИ (AC-6(1))}


\vspace{10pt}
{\footnotesize\bfseries
No: 1\\
Name: ac\_6\_1\_a\_01\\
Type: string\\
Default: nil
}

{\footnotesize Авторизовано доступ для осіб та ролей до функцій безпеки (розгорнуті на апаратному забезпеченні)}


{\footnotesize\bfseries
No: 2\\
Name: ac\_6\_1\_a\_02\\
Type: string\\
Default: nil
}

{\footnotesize Авторизовано доступ для осіб та ролей до функцій безпеки (роз- горнуті на програмному забезпеченні)}


{\footnotesize\bfseries
No: 3\\
Name: ac\_6\_1\_a\_03\\
Type: string\\
Default: nil
}

{\footnotesize Авторизовано доступ для осіб та ролей до функцій безпеки (розгорнуті на мікропрограмному забезпеченні)}


{\footnotesize\bfseries
No: 4\\
Name: ac\_6\_1\_b\\
Type: string\\
Default: nil
}

{\footnotesize Авторизовано доступ для осіб та ролей до інформації}


{\footnotesize\bfseries
No: 5\\
Name: ac\_6\_1\_odp\_01\\
Type: list\\
Default: [\textquotedbl{}admin\textquotedbl{}, \textquotedbl{}security\_officer\textquotedbl{}]
}

{\footnotesize Визначені особи або ролі з авторизованим доступом до функцій безпеки та інформації, що має відношення до безпеки}


{\footnotesize\bfseries
No: 6\\
Name: ac\_6\_1\_odp\_02\\
Type: string\\
Default: nil
}

{\footnotesize Визначені функції безпеки (розгорнуті в апаратному забезпеченні) для авторизованого доступу}


{\footnotesize\bfseries
No: 7\\
Name: ac\_6\_1\_odp\_03\\
Type: string\\
Default: nil
}

{\footnotesize Визначені функції безпеки (розгорнуті в програмному забезпеченні) для авторизованого доступу}


{\footnotesize\bfseries
No: 8\\
Name: ac\_6\_1\_odp\_04\\
Type: string\\
Default: nil
}

{\footnotesize Визначені функції безпеки (розгорнуті в мікропрограмному забезпеченні) для авторизованого доступу}


{\footnotesize\bfseries
No: 9\\
Name: ac\_6\_1\_odp\_05\\
Type: string\\
Default: nil
}

{\footnotesize Визначено інформацію, важливу для забезпечення безпеки, для авторизованого доступу}




\subsubsection{НЕПРИВІЛЕЙОВАНИЙ ДОСТУП ДО НЕЗАХИЩЕНИХ ФУНКЦІЙ (AC-6(2))}
Вимагати від користувачів облікових записів системи або ролей, які мають доступ до [Призначення: визначених організацією функцій безпеки або інформації, що стосується безпеки], використовувати непривілейовані облікові записи чи ролі під час доступу до незахищених функцій.

\vspace{10pt}
{\footnotesize\bfseries
No: 1\\
Name: ac\_6\_2\_01\\
Type: list\\
Default: [\textquotedbl{}admin\textquotedbl{}, \textquotedbl{}security\_officer\textquotedbl{}]
}

{\footnotesize Користувачі облікових записів (або ролей) системи з доступом до функцій безпеки або інформації, що стосується безпеки, повинні використовувати непривілейовані облікові записи або ролі під час доступу до незахищених функцій}


{\footnotesize\bfseries
No: 2\\
Name: ac\_6\_2\_odp\\
Type: string\\
Default: nil
}

{\footnotesize Визначені функції безпеки або інформація, що стосується безпеки}




\subsubsection{ПРИВІЛЕЙОВАНІ ОБЛІКОВІ ЗАПИСИ (AC-6(5))}
Обмежити привілейовані облікові записи в системі згідно з [Призначення: визначеним організацією персоналом або ролями].

\vspace{10pt}
{\footnotesize\bfseries
No: 1\\
Name: ac\_6\_5\_01\\
Type: list\\
Default: [\textquotedbl{}admin\textquotedbl{}, \textquotedbl{}security\_officer\textquotedbl{}]
}

{\footnotesize Привілейовані облікові записи в системі обмежено персонал або ролі}


{\footnotesize\bfseries
No: 2\\
Name: ac\_6\_5\_odp\\
Type: list\\
Default: [\textquotedbl{}admin\textquotedbl{}, \textquotedbl{}security\_officer\textquotedbl{}]
}

{\footnotesize Визначено персонал або ролі, яким мають бути обмежені привілейовані облікові записи в системі}




\subsubsection{ПЕРЕГЛЯД ПОВНОВАЖЕНЬ КОРИСТУВАЧА (AC-6(7))}
a) Переглядати [Призначення: з визначеною організацією частотою] повноваження призначених для [Призначення: визначених організацією посад або класів користувачів] для перевірки необхідності таких повноважень; b) За необхідності перепризначити або зняти повноваження, правильного відображення цілей (місії) організації та потреб організації. для

\vspace{10pt}
{\footnotesize\bfseries
No: 1\\
Name: ac\_6\_7\_a\\
Type: integer\\
Default: 30
}

{\footnotesize Повноваження, призначені ролям і класам, переглядаються з частотою для перевірки необхідності таких повноважень}


{\footnotesize\bfseries
No: 2\\
Name: ac\_6\_7\_b\\
Type: string\\
Default: nil
}

{\footnotesize Привілеї перепризначаються або знімаються, якщо це необхідно, для правильного відображення місії організації та потреб}


{\footnotesize\bfseries
No: 3\\
Name: ac\_6\_7\_odp\_01\\
Type: integer\\
Default: 30
}

{\footnotesize Визначено частоту перегляду повноважень, призначених ролям або класам користувачів}


{\footnotesize\bfseries
No: 4\\
Name: ac\_6\_7\_odp\_02\\
Type: list\\
Default: [\textquotedbl{}admin\textquotedbl{}, \textquotedbl{}security\_officer\textquotedbl{}]
}

{\footnotesize Визначено ролі або класи користувачів, яким призначено повноваження }




\subsubsection{АУДИТ ВИКОРИСТАННЯ ПРИВІЛЕЙОВАНИХ ФУНКЦІЙ (AC-6(9))}
Реєструвати виконання привілейованих функцій.

\vspace{10pt}
{\footnotesize\bfseries
No: 1\\
Name: ac\_6\_9\_01\\
Type: string\\
Default: nil
}

{\footnotesize Проводиться аудит виконання привілейованих функцій}




\subsubsection{ЗАБОРОНА НЕПРИВІЛЕЙОВАНИМ КОРИСТУВАЧАМ ВИКОНУВАТИ ПРИВІЛЕЙОВАНІ ФУНКЦІЇ (AC-6(10))}
Вжити заходи для запобігання можливості виконувати привілейовані функції непривілейованими користувачами.

\vspace{10pt}
{\footnotesize\bfseries
No: 1\\
Name: ac\_6\_10\_01\\
Type: string\\
Default: nil
}

{\footnotesize МІНІМІЗАЦІЯ ПОВНОВАЖЕНЬ - ЗАБОРОНА НЕПРИВІЛЕЙОВАНИМ КОРИСТУВАЧАМ ВИКОНУВАТИ ПРИВІЛЕЙОВАНІ ФУНКЦІЇ МЕТА ОЦІНКИ: Визначити, чи:}




\subsection{НЕВДАЛІ СПРОБИ ВХОДУ В СИСТЕМУ (AC-7)}
a. Встановити обмеження на [Призначення: визначену організацією кількість] послідовних неуспішних спроб входу користувача в систему впродовж [Призначення: визначеного організацією часового періоду].\\ 
b. Автоматично виконати [Вибір (один або декілька): блокування облікового запису/вузла на [Призначення: визначений організацією часовий період]; блокування облікового запису/вузла, доки він не буде розблокований адміністратором; затримання наступної команди входу в систему за [Надання: визначеним організацією алгоритмом затримки]; виконати [Призначення: визначені організацією дії]], коли перевищено максимальну кількість невдалих спроб входу в систему.

\vspace{10pt}
{\footnotesize\bfseries
No: 1\\
Name: ac\_7\_a\\
Type: integer\\
Default: 30
}

{\footnotesize Застосовано обмеження на кількість послідовних неуспішних спроб входу користувача протягом періоду часу}


{\footnotesize\bfseries
No: 2\\
Name: ac\_7\_odp\_01\\
Type: integer\\
Default: 30
}

{\footnotesize Визначається кількість послідовних неуспішних спроб входу користувача, дозволених протягом певного періоду часу}


{\footnotesize\bfseries
No: 3\\
Name: ac\_7\_odp\_02\\
Type: integer\\
Default: 30
}

{\footnotesize Визначено період часу, яким обмежується кількість послідовних неуспішних спроб входу користувача}


{\footnotesize\bfseries
No: 4\\
Name: ac\_7\_odp\_03\\
Type: string\\
Default: \textquotedbl{}AES-256-GCM\textquotedbl{}
}

{\footnotesize Вибрано одне або декілька з наступних ЗНАЧЕНЬ ПАРАМЕТРІВ: \{заблокувати обліковий запис або вузол на період часу; заблокувати обліковий запис або вузол до зняття адміністратором; затримати наступний запит на вхід за алгоритмом затримки; повідомити системного адміністратора; виконати іншу дію\}}


{\footnotesize\bfseries
No: 5\\
Name: ac\_7\_odp\_04\\
Type: integer\\
Default: 30
}

{\footnotesize Період часу, на який буде заблоковано обліковий запис або вузол (якщо вибрано)}


{\footnotesize\bfseries
No: 6\\
Name: ac\_7\_odp\_05\\
Type: string\\
Default: \textquotedbl{}AES-256-GCM\textquotedbl{}
}

{\footnotesize Визначено алгоритм затримки наступного запиту на вхід (якщо вибрано)}


{\footnotesize\bfseries
No: 7\\
Name: ac\_7\_odp\_06\\
Type: integer\\
Default: 3
}

{\footnotesize Інша дія, яка буде виконана після перевищення максимальної кількості невдалих спроб (якщо вибрано)}




\subsection{ПОПЕРЕДЖЕННЯ ПРО ВИКОРИСТАННЯ СИСТЕМИ (AC-8)}
a.\\ 
b. Демонструвати користувачам [Призначення: визначене організацією сповіщення або банер про використання системи] перед тим, як надавати доступ до системи, що забезпечує безпеку та приватність відповідно до чинних законів, нормативних документів, наказів, директив, політик, правил, стандартів і керівних принципів, які зазначають, що:\\ 
1. користувачі здійснюють доступ до урядової системи;\\ 
2. використання системи може контролюватися, реєструватися та підлягати аудиту;\\ 
3. несанкціоноване використання системи забороняється та приводить до кримінальної та цивільної відповідальності;\\ 
4. використання системи означає згоду на моніторинг і запис дій користувача. Зберігати сповіщення або банер на екрані, доки користувачі не визнають умови використання та не приймуть явних дій для входу в систему або подальшого доступу до системи.\\ 
c. Для загальнодоступних систем:\\ 
1. демонструвати інформацію про умови використання системи [Призначення: визначені організацією умови], перш ніж надавати подальший доступ до загальнодоступної системи;\\ 
2. демонструвати посилання, якщо такі є, на моніторинг, запис або аудит, які узгоджуються з акомодацією приватності для таких систем, які зазвичай забороняють такі дії;\\ 
3. мати опис авторизованого використання системи.

\vspace{10pt}
{\footnotesize\bfseries
No: 1\\
Name: ac\_8\_c\_02\\
Type: list\\
Default: [\textquotedbl{}default\_deny\_rule\textquotedbl{}, \textquotedbl{}abac\_rule\_1\textquotedbl{}]
}

{\footnotesize Визначається сповіщення або банер про використання системи, який система буде показувати користувачам перед наданням доступу до системи; визначено умови використання системи, які система відображатиме перед наданням подальшого доступу; повідомлення про використання системи відображається користувачам перед наданням доступу до системи, яке містить повідомлення про конфіденційність і безпеку, що відповідають чинним законам, нормативним документам, наказам, директивам, політикам, правилам, стандартам і керівним принципам; у повідомленні про використання системи зазначено, що користувачі отримують доступ до урядової системи; у повідомленні про використання системи зазначено, що використання системи може контролюватися, реєструватися та підлягати аудиту; у повідомленні про використання системи зазначено, що несанкціоноване використання системи заборонено і тягне за собою кримінальну та цивільну відповідальність; у повідомленні про використання системи зазначено, що використання системи означає згоду на моніторинг і запис дій; сповіщення або банер залишається на екрані до тих пір, поки користувачі не визнають умови використання та не приймуть явні дії для входу в систему або подальшого доступу до неї; для загальнодоступних систем інформація про використання системи умови відображається перед наданням подальшого доступу до загальнодоступної системи; для загальнодоступних систем відображаються будь-які посилання на моніторинг, запис або аудит, які узгоджуються з положеннями про конфіденційність таких систем, що зазвичай забороняють ці види діяльності; для загальнодоступних систем додається опис авторизованого використання системи}




\subsection{БЛОКУВАННЯ ПРИСТРОЮ (AC-11)}
a. Заборонити подальший доступ до системи шляхом ініціювання блокування пристрою після [Призначення: визначеного організацією періоду] бездіяльності або після отримання запиту від користувача.\\ 
b. Зберігати блокування пристрою, поки користувач не відновить доступ, використовуючи встановлені процедури ідентифікації та автентифікації.

\vspace{10pt}
{\footnotesize\bfseries
No: 1\\
Name: ac\_11\_b\\
Type: string\\
Default: nil
}

{\footnotesize Блокування пристрою зберігається доти, доки користувач не відновить доступ за допомогою встановлених процедур ідентифікації та автентифікації}


{\footnotesize\bfseries
No: 2\\
Name: ac\_11\_odp\_01\\
Type: integer\\
Default: 30
}

{\footnotesize Вибрано одне або декілька з наступних ЗНАЧЕНЬ ПАРАМЕТРІВ: \{ініціювання блокування пристрою після періоду неактивності; вимога до користувача ініціювати блокування пристрою перед тим, як залишити систему без нагляду\}}


{\footnotesize\bfseries
No: 3\\
Name: ac\_11\_odp\_02\\
Type: integer\\
Default: 30
}

{\footnotesize Часовий проміжок бездіяльності, після якого ініціюється блокування пристрою (якщо вибрано)}




\subsubsection{ПРИХОВАНІ ДИСПЛЕЇ (AC-11(1))}


\vspace{10pt}
{\footnotesize\bfseries
No: 1\\
Name: ac\_11\_1\_01\\
Type: string\\
Default: nil
}

{\footnotesize Система приховує (через блокування сеансу) інформацію, попередньо видиму на дисплеї, загальнодоступним зображенням}




\subsection{Припинення сеансу (AC-12)}
Сеанс користувача має завершуватися автоматично після [Призначення: визначених організацією умов або тригерних подій, що вимагають припинення сеансу].

\vspace{10pt}
{\footnotesize\bfseries
No: 1\\
Name: ac\_12\_01\\
Type: string\\
Default: nil
}

{\footnotesize Сеанс користувача автоматично завершується після виконання умов або подій}


{\footnotesize\bfseries
No: 2\\
Name: ac\_12\_odp\\
Type: list\\
Default: [\textquotedbl{}login\textquotedbl{}, \textquotedbl{}logout\textquotedbl{}, \textquotedbl{}failed\_attempt\textquotedbl{}]
}

{\footnotesize Визначено умови або події, що вимагають припинення сеансу}




\subsection{ВІДДАЛЕНИЙ ДОСТУП (AC-17)}
a. Встановити та задокументувати обмеження на використання, вимоги до конфігурації/підключення та рекомендації щодо здійснення кожного типу віддаленого доступу.\\ 
b. Авторизувати віддалений доступ до системи, перш ніж будуть дозволені такі підключення.

\vspace{10pt}
{\footnotesize\bfseries
No: 1\\
Name: ac\_17\_a\_02\\
Type: string\\
Default: nil
}

{\footnotesize Для кожного типу дозволеного віддаленого доступу встановлені та задокументовані вимоги до конфігурації/підключення}


{\footnotesize\bfseries
No: 2\\
Name: ac\_17\_a\_03\\
Type: string\\
Default: nil
}

{\footnotesize Для кожного типу дозволеного віддаленого доступу встановлені та задокументовані рекомендації}


{\footnotesize\bfseries
No: 3\\
Name: ac\_17\_b\\
Type: string\\
Default: nil
}

{\footnotesize Кожен тип віддаленого доступу до системи авторизується перед тим, як дозволити такі підключення}




\subsubsection{КЕРОВАНІ ТОЧКИ КОНТРОЛЮ ДОСТУПУ (AC-17(3))}


\vspace{10pt}
{\footnotesize\bfseries
No: 1\\
Name: ac\_17\_3\_01\\
Type: string\\
Default: nil
}

{\footnotesize Віддалений доступ маршрутизується через авторизовані та керовані точки контролю доступу до мережі}




\subsubsection{ПРИВІЛЕЙОВАНІ КОМАНДИ ТА ДОСТУП (AC-17(4))}


\vspace{10pt}
{\footnotesize\bfseries
No: 1\\
Name: ac\_17\_4\_a\_01\\
Type: string\\
Default: nil
}

{\footnotesize Виконання привілейованих команд за допомогою віддаленого доступу дозволено лише для наступних потреб: потреби }


{\footnotesize\bfseries
No: 2\\
Name: ac\_17\_4\_a\_02\\
Type: string\\
Default: nil
}

{\footnotesize Доступ до інформації, важливої для безпеки, за допомогою віддаленого доступу дозволяється лише для наступних потреб: потреби }


{\footnotesize\bfseries
No: 3\\
Name: ac\_17\_4\_b\\
Type: string\\
Default: nil
}

{\footnotesize Обґрунтування віддаленого доступу задокументовано в плані захисту інформації}


{\footnotesize\bfseries
No: 4\\
Name: ac\_17\_4\_odp\_01\\
Type: string\\
Default: nil
}

{\footnotesize Визначено потреби, що потребують виконання привілейованих команд за допомогою віддаленого доступу; AC-17(04)\_ODP[02] визначено потреби, що вимагають доступу до інформації, що стосується безпеки, за допомогою віддаленого доступу}




\subsection{БЕЗДРОТОВИЙ ДОСТУП (AC-18)}
a. Установити обмеження на використання, вимоги до конфігурації/підключення та рекомендації щодо здійснення бездротового доступу.\\ 
b. Авторизувати бездротовий доступ до системи, перш ніж будуть дозволені такі підключення.

\vspace{10pt}
{\footnotesize\bfseries
No: 1\\
Name: ac\_18\_a\_01\\
Type: string\\
Default: nil
}

{\footnotesize Встановлено обмеження на використання щодо здійснення бездротового доступу}


{\footnotesize\bfseries
No: 2\\
Name: ac\_18\_a\_02\\
Type: string\\
Default: nil
}

{\footnotesize Встановлено вимоги до конфігурації або підключення щодо здійснення бездротового доступу}


{\footnotesize\bfseries
No: 3\\
Name: ac\_18\_a\_03\\
Type: string\\
Default: nil
}

{\footnotesize Встановлено доступу}


{\footnotesize\bfseries
No: 4\\
Name: ac\_18\_b\\
Type: string\\
Default: nil
}

{\footnotesize Авторизується бездротовий доступ до системи перед тим, як дозволяти такі з'єднання. рекомендації щодо здійснення бездротового}




\subsubsection{АВТЕНТИФІКАЦІЯ ТА ШИФРУВАННЯ (AC-18(1))}


\vspace{10pt}
{\footnotesize\bfseries
No: 1\\
Name: ac\_18\_1\_02\\
Type: string\\
Default: \textquotedbl{}AES-256-GCM\textquotedbl{}
}

{\footnotesize Бездротовий доступ до системи захищений за допомогою шифрування}




\subsubsection{ВІДКЛЮЧЕННЯ БЕЗДРОТОВОЇ МЕРЕЖІ (AC-18(3))}


\vspace{10pt}
{\footnotesize\bfseries
No: 1\\
Name: ac\_18\_3\_01\\
Type: string\\
Default: nil
}

{\footnotesize Відключено, у разі відсутності необхідності у використанні, вбудовані в компоненти системи можливості бездротових мереж до їх виклику та розгортання}




\subsection{КОНТРОЛЬ ДОСТУПУ ДЛЯ МОБІЛЬНИХ ПРИСТРОЇВ (AC-19)}
a. Встановити обмеження на використання, вимоги до конфігурації, вимоги до підключення і рекомендації щодо впровадження мобільних пристроїв, контрольованих організацією.\\ 
b. Авторизувати підключення мобільних пристроїв до систем, які експлуатуються організацією.

\vspace{10pt}
{\footnotesize\bfseries
No: 1\\
Name: ac\_19\_a\_01\\
Type: string\\
Default: nil
}

{\footnotesize Встановлено вимоги до конфігурації мобільних пристроїв, що контролюються організацією, в тому числі, коли такі пристрої перебувають за межами контрольованої території}


{\footnotesize\bfseries
No: 2\\
Name: ac\_19\_a\_02\\
Type: string\\
Default: nil
}

{\footnotesize Встановлюються вимоги до підключення для мобільних пристроїв, що контролюються організацією, в тому числі, коли такі пристрої знаходяться за межами контрольованої території}


{\footnotesize\bfseries
No: 3\\
Name: ac\_19\_a\_03\\
Type: string\\
Default: nil
}

{\footnotesize Розроблено рекомендації щодо впровадження для мобільних пристроїв, які контролюються організацією, в тому числі, коли такі пристрої перебувають за межами контрольованої території}


{\footnotesize\bfseries
No: 4\\
Name: ac\_19\_b\\
Type: string\\
Default: nil
}

{\footnotesize Авторизовано підключення мобільних пристроїв до систем організації}




\subsubsection{ПОВНЕ ШИФРУВАННЯ ПРИСТРОЇВ ТА СХОВИЩ ІНФОРМАЦІЇ (AC-19(5))}


\vspace{10pt}
\textit{Немає параметрів для цього контролю.}
\vspace{10pt}


\subsection{ВИКОРИСТАННЯ ЗОВНІШНІХ СИСТЕМ (AC-20)}
a. [Вибір (один або кілька): Встановіть [Призначення: умови, визначені організацією]; Визначте [Призначення: визначені організацією засоби контролю, які, як стверджується, будуть реалізовані на зовнішніх системах]], узгоджені з довірчими відносинами, встановленими з іншими організаціями, які володіють, експлуатують та/або обслуговують зовнішні системи, дозволяючи уповноваженим особам:\\ 
1. доступ до системи із зовнішніх систем;\\ 
2. обробляти, зберігати або передавати керовану організацією інформацію за допомогою зовнішніх систем;\\ 
b. Заборонити використання [Призначення: організаційно-визначені типи зовнішніх систем].

\vspace{10pt}
{\footnotesize\bfseries
No: 1\\
Name: ac\_20\_odp\_01\\
Type: list\\
Default: []
}

{\footnotesize Вибрано одне або більше з наступних ЗНАЧЕНЬ ПАРАМЕТРІВ: \{встановити умови та положення; визначити заходи захисту\}}


{\footnotesize\bfseries
No: 2\\
Name: ac\_20\_odp\_02\\
Type: list\\
Default: []
}

{\footnotesize Визначено умови та положення, що відповідають довірчим відносинам, встановленим з іншими організаціями, які володіють, експлуатують та/або обслуговують зовнішні системи (якщо вибрано)}


{\footnotesize\bfseries
No: 3\\
Name: ac\_20\_odp\_03\\
Type: string\\
Default: nil
}

{\footnotesize Визначено заходи захисту, які мають бути застосовані до зовнішніх систем відповідно до довірчих відносин, встановлених з іншими організаціями, що володіють, експлуатують та/або обслуговують зовнішні системи (якщо обрано)}


{\footnotesize\bfseries
No: 4\\
Name: ac\_20\_odp\_04\\
Type: string\\
Default: nil
}

{\footnotesize Визначено типи зовнішніх систем, заборонених до використання}




\subsubsection{ОБМЕЖЕННЯ НА АВТОРИЗОВАНЕ ВИКОРИСТАННЯ (AC-20(1))}


\vspace{10pt}
{\footnotesize\bfseries
No: 1\\
Name: ac\_20\_1\_a\\
Type: list\\
Default: [\textquotedbl{}admin\textquotedbl{}, \textquotedbl{}security\_officer\textquotedbl{}]
}

{\footnotesize Авторизовані особи мають право використовувати зовнішню систему для доступу до системи або для обробки, зберігання чи передачі інформації, що контролюється організацією, лише після перевірки виконання заходів безпеки та конфіденційності, зазначених у політиці безпеки та конфіденційності організації, а також планах безпеки та конфіденційності (якщо такі застосовуються)}


{\footnotesize\bfseries
No: 2\\
Name: ac\_20\_1\_b\\
Type: list\\
Default: [\textquotedbl{}admin\textquotedbl{}, \textquotedbl{}security\_officer\textquotedbl{}]
}

{\footnotesize Авторизовані особи мають право використовувати зовнішню систему для доступу до системи або для обробки, зберігання чи передачі інформації, що контролюється організацією, лише після збереження погоджених угод про підключення або обробку системи з структурою організації, на якій розміщена зовнішня система}




\subsubsection{ПЕРЕНОСНІ ПРИСТРОЇ ЗБЕРІГАННЯ ДАНИХ (AC-20(2))}


\vspace{10pt}
{\footnotesize\bfseries
No: 1\\
Name: ac\_20\_2\_01\\
Type: list\\
Default: [\textquotedbl{}admin\textquotedbl{}, \textquotedbl{}security\_officer\textquotedbl{}]
}

{\footnotesize Використання портативних пристроїв носіїв інформації уповноваженими особами обмежено у зовнішніх системах за допомогою обмеження }


{\footnotesize\bfseries
No: 2\\
Name: ac\_20\_2\_odp\\
Type: list\\
Default: [\textquotedbl{}admin\textquotedbl{}, \textquotedbl{}security\_officer\textquotedbl{}]
}

{\footnotesize Визначено обмеження на використання авторизованими особами портативних носіїв інформації у зовнішніх системах}




\subsection{ПУБЛІЧНО ДОСТУПНИЙ КОНТЕНТ (AC-22)}
a. Призначити осіб, що уповноважені на розміщення інформації в загальнодоступній системі.\\ 
b. Навчати уповноважених осіб тому, щоб загальнодоступна інформація не містила інформацію з обмеженим доступом.\\ 
c. Переглядати запропонований зміст інформації до публікації в загальнодоступній системі, щоб гарантувати, що там не міститься інформація з обмеженим доступом.\\ 
d. Переглядати вміст загальнодоступної системи на предмет наявності там інформації з обмеженим доступом з [Призначення: визначеною організацією частотою]; така інформація має бути видалена в разі її виявлення.

\vspace{10pt}
{\footnotesize\bfseries
No: 1\\
Name: ac\_22\_a\\
Type: list\\
Default: [\textquotedbl{}admin\textquotedbl{}, \textquotedbl{}security\_officer\textquotedbl{}]
}

{\footnotesize Визначені особи уповноважені на розміщення інформації в загальнодоступній системі}


{\footnotesize\bfseries
No: 2\\
Name: ac\_22\_b\\
Type: list\\
Default: [\textquotedbl{}admin\textquotedbl{}, \textquotedbl{}security\_officer\textquotedbl{}]
}

{\footnotesize Уповноважені особи проходять навчання, щоб гарантувати, що загальнодоступна інформація не містить інформацію з обмеженим доступом}


{\footnotesize\bfseries
No: 3\\
Name: ac\_22\_c\\
Type: string\\
Default: nil
}

{\footnotesize Запропонований зміст інформації перевіряється до публікації в загальнодоступній системі, щоб гарантувати, що там не міститься інформація з обмеженим доступом}


{\footnotesize\bfseries
No: 4\\
Name: ac\_22\_d\_01\\
Type: integer\\
Default: 30
}

{\footnotesize Зміст у загальнодоступній системі перевіряється на наявність інформації з обмеженим доступом з частотою}


{\footnotesize\bfseries
No: 5\\
Name: ac\_22\_d\_02\\
Type: string\\
Default: nil
}

{\footnotesize Інформація з обмеженим доступом видаляється з загальнодоступної системи, якщо її виявлено}


{\footnotesize\bfseries
No: 6\\
Name: ac\_22\_odp\\
Type: integer\\
Default: 30
}

{\footnotesize Визначено частоту, з якою слід переглядати вміст загальнодоступної системи на предмет наявності там інформації з обмеженим доступом}




\section{AT}
Клас заходів захисту AT --- ОБІЗНАНІСТЬ ТА НАВЧАННЯ

Цей клас спрямований на забезпечення належного рівня знань персоналу щодо загроз інформаційній безпеці та їхніх обов'язків у цій сфері.

\textbf{Перелік заходів захисту:} Політика та процедури підвищення обізнаності та навчання (AT-1); Навчання з підвищення обізнаності (AT-2); НАВЧАННЯ З ПІДВИЩЕННЯ ОБІЗНАНОСТІ | НАВЧАННЯ ЩОДО РОЗПІЗНАВАННЯ ВНУТРІШНЬОЇ ЗАГРОЗИ; НАВЧАННЯ З ПІДВИЩЕННЯ ОБІЗНАНОСТІ | НАВЧАННЯ З ПИТАНЬ СОЦІАЛЬНОЇ ІНЖЕНЕРІЇ; Рольове навчання (AT-3).

\subsection{ПОЛІТИКА ТА ПРОЦЕДУРИ ПІДВИЩЕННЯ ОБІЗНАНОСТІ ТА НАВЧАННЯ (AT-1)}
a. Розробити, задокументувати та поширити [Призначення: серед визначеного організацією персоналу або ролей]:\\ 
1. [Вибір (один або декілька): Рівень організації; Рівень місії/бізнес-процесу; рівень системи] політики обізнаності та навчання у сфері забезпечення безпеки та приватності, яка:\\ 
(a) містить мету, сферу застосування, ролі, обов'язки, відповідальність керівництва, координацію між організаційними підрозділами та систему контролю відповідності (complaince);\\ 
(b) відповідає чинним законам, нормативним документам, директивам, нормам, політикам, стандартам та керівним документам.\\ 
2. Процедури, що сприяють реалізації політики підвищення обізнаності та професійної підготовки в галузі безпеки, приватності, а також пов'язаних з ними заходів захисту інформації та персональних даних.\\ 
b. Призначити [Призначення: визначену організацією посадову особу] для управління політикою та процедурами підвищення обізнаності та навчання у сфері забезпечення безпеки та приватності.\\ 
c. Переглядати та оновлювати:\\ 
1. Поточну політика [Призначення: частота, визначена організацією] і наступне [Призначення: події, визначені організацією];\\ 
2. Процедури [Призначення: частота, визначена організацією] та наступні [Призначення: події, визначені організацією].

\vspace{10pt}
\textit{Немає параметрів для цього контролю.}
\vspace{10pt}


\subsection{НАВЧАННЯ З ПІДВИЩЕННЯ ОБІЗНАНОСТІ (AT-2)}
Впровадити базові тренінги з підвищення обізнаності у сфері безпеки та приватності для користувачів системи (включно з менеджерами, керівниками компаній і підрядниками):\\ 
a. Забезпечити навчання грамотності з питань безпеки та конфіденційності для користувачів системи (включаючи менеджерів, керівників вищої ланки та підрядників):\\ 
1. як частину початкового навчання для нових користувачів і [Призначення: частота, визначена організацією] після цього;\\ 
2. якщо цього потребують системні зміни або наступні [Призначення: події, визначені організацією].\\ 
b. Використовувати наведені нижче методи, щоб підвищити рівень безпеки та конфіденційності користувачів системи [Завдання: визначені організацією методи поінформованості];\\ 
c. Оновлювати навчання грамотності та зміст обізнаності [Завдання: частота, визначена організацією] і наступні [Завдання: події, визначені організацією];\\ 
d. Включити уроки, отримані з внутрішніх або зовнішніх інцидентів безпеки або порушень, у навчання грамотності та методи підвищення обізнаності.

\vspace{10pt}
{\footnotesize\bfseries
No: 1\\
Name: at\_2\_b\\
Type: string\\
Default: nil
}

{\footnotesize Методи застосовуються для підвищення обізнаності користувачів системи щодо безпеки та конфіденційності}


{\footnotesize\bfseries
No: 2\\
Name: at\_2\_c\_01\\
Type: integer\\
Default: 30
}

{\footnotesize Оновлюється зміст навчання грамотності та підвищення обізнаності з частотою}


{\footnotesize\bfseries
No: 3\\
Name: at\_2\_c\_02\\
Type: string\\
Default: nil
}

{\footnotesize Оновлюється зміст навчання грамотності та підвищення обізнаності після подій}


{\footnotesize\bfseries
No: 4\\
Name: at\_2\_d\\
Type: string\\
Default: nil
}

{\footnotesize Уроки, отримані в результаті внутрішніх або зовнішніх інцидентів або порушень безпеки, включені в методи навчання та підвищення обізнаності}


{\footnotesize\bfseries
No: 5\\
Name: at\_2\_odp\_01\\
Type: string\\
Default: \textquotedbl{}щорічно\textquotedbl{}
}

{\footnotesize Визначено періодичність проведення навчання грамотності з питань безпеки для користувачів системи (в тому числі менеджерів, вищого керівництва та підрядників) після початкового тренінгу}


{\footnotesize\bfseries
No: 6\\
Name: at\_2\_odp\_02\\
Type: string\\
Default: \textquotedbl{}щорічно\textquotedbl{}
}

{\footnotesize Визначено періодичність проведення навчання грамотності з питань конфіденційності для користувачів системи (в тому числі менеджерів, вищого керівництва та підрядників) після початкового тренінгу}


{\footnotesize\bfseries
No: 7\\
Name: at\_2\_odp\_03\\
Type: list\\
Default: [\textquotedbl{}login\textquotedbl{}, \textquotedbl{}logout\textquotedbl{}, \textquotedbl{}failed\_attempt\textquotedbl{}]
}

{\footnotesize Визначено події, які потребують навчання користувачів системи грамотності з питань безпеки}


{\footnotesize\bfseries
No: 8\\
Name: at\_2\_odp\_05\\
Type: string\\
Default: nil
}

{\footnotesize Визначено методи, які слід застосовувати для підвищення обізнаності користувачів системи щодо безпеки та конфіденційністі}


{\footnotesize\bfseries
No: 9\\
Name: at\_2\_odp\_06\\
Type: list\\
Default: [\textquotedbl{}login\textquotedbl{}, \textquotedbl{}logout\textquotedbl{}, \textquotedbl{}failed\_attempt\textquotedbl{}]
}

{\footnotesize Визначено частоту оновлення навчання грамотності та змісту обізнаності; AT-02\_ODP[07] визначено події після яких необхідне оновлення навчання грамотності та змісту обізнаності; AT-02(a)[01][01] навчання з грамотності з питань безпеки надається користувачам системи (включаючи менеджерів, керівників вищої ланки та підрядників) як частина початкового навчання для нових користувачів; AT-02(a)[01][02] навчання з грамотності з питань конфіденційності надається користувачам системи (включаючи менеджерів, керівників вищої ланки та підрядників) як частина початкового навчання для нових користувачів; AT-02(a)[01][03] для користувачів системи (включно з менеджерами, вищим керівництвом та підрядниками) проводиться навчання з безпеки з періодичністю після цього; AT-02(a)[01][04] для користувачів системи (включно з менеджерами, вищим керівництвом та підрядниками) проводиться навчання з конфіденційності з періодичністю після цього; AT-02(a)[02][01] тренінги з грамотності з питань безпеки проводяться для користувачів системи (включаючи менеджерів, керівників вищої ланки та підрядників), коли цього вимагають зміни в системі або після подій; AT-02(a)[02][02] тренінги з грамотності з питань конфіденціності проводяться для користувачів системи (включаючи менеджерів, керівників вищої ланки та підрядників), коли цього вимагають зміни в системі або після подій}


{\footnotesize\bfseries
No: 10\\
Name: at\_2\_odp\_4\\
Type: list\\
Default: [\textquotedbl{}login\textquotedbl{}, \textquotedbl{}logout\textquotedbl{}, \textquotedbl{}failed\_attempt\textquotedbl{}]
}

{\footnotesize Визначено події, які потребують навчання користувачів системи грамотності з питань конфіденційності}




\subsubsection{НАВЧАННЯ З ПІДВИЩЕННЯ ОБІЗНАНОСТІ | НАВЧАННЯ ЩОДО РОЗПІЗНАВАННЯ ВНУТРІШНЬОЇ ЗАГРОЗИ}
Ввести до програми навчання вправи з розпізнавання та виявлення потенційних індикаторів внутрішніх загроз.

\vspace{10pt}
{\footnotesize\bfseries
No: 1\\
Name: at\_2\_2\_01\\
Type: string\\
Default: nil
}

{\footnotesize Введено до програми навчання вправи з розпізнавання потенційних індикаторів внутрішніх загроз}


{\footnotesize\bfseries
No: 2\\
Name: at\_2\_2\_02\\
Type: string\\
Default: nil
}

{\footnotesize Введено до програми навчання вправи з виявлення потенційних індикаторів внутрішніх загроз}




\subsubsection{НАВЧАННЯ З ПІДВИЩЕННЯ ОБІЗНАНОСТІ | НАВЧАННЯ З ПИТАНЬ СОЦІАЛЬНОЇ ІНЖЕНЕРІЇ}
Ввести до програми навчання вправи з підвищення обізнаності щодо розпізнавання та повідомлення про потенційні та фактичні атаки, з використанням методів соціальної інженерії та інтелектуального аналізу соціальних даних.

\vspace{10pt}
{\footnotesize\bfseries
No: 1\\
Name: at\_2\_3\_01\\
Type: string\\
Default: nil
}

{\footnotesize До програми навчання введено вправи з розпізнавання потенційних та фактичних випадків соціального інжинірингу}


{\footnotesize\bfseries
No: 2\\
Name: at\_2\_3\_02\\
Type: string\\
Default: nil
}

{\footnotesize До програми навчання введено вправи з повідомлення про потенційні та фактичні випадки соціального інжинірингу}


{\footnotesize\bfseries
No: 3\\
Name: at\_2\_3\_03\\
Type: string\\
Default: nil
}

{\footnotesize До програми навчання введено вправи з розпізнавання потенційних та фактичних випадків інтелектуального аналізу соціальних даних}


{\footnotesize\bfseries
No: 4\\
Name: at\_2\_3\_04\\
Type: string\\
Default: nil
}

{\footnotesize До програми навчання введено вправи з повідомлення про потенційні та фактичні випадки інтелектуального аналізу соціальних даних}




\subsection{РОЛЬОВЕ НАВЧАННЯ (AT-3)}
a. Забезпечити проведення навчання з питань безпеки та приватності на основі ролей для працівників з ролями та обов'язками: [Призначення: визначені організацією ролі та обов'язки]:\\ 
1. перед авторизацією доступу до системи, інформації або виконанням призначених обов'язків і [Призначення: частота, визначена організацією] після цього;\\ 
2. коли цього потребують системні зміни.\\ 
b. Оновити навчальний контент на основі ролей [Призначення: частота, визначена організацією] і наступні [Призначення: події, визначені організацією];\\ 
c. Включіть у рольове навчання, інформацію, отриману з внутрішніх або зовнішніх інцидентів та порушень безпеки.

\vspace{10pt}
{\footnotesize\bfseries
No: 1\\
Name: at\_3\_b\_01\\
Type: integer\\
Default: 30
}

{\footnotesize Оновлюється вміст навчання на основі ролей з частотою}


{\footnotesize\bfseries
No: 2\\
Name: at\_3\_b\_02\\
Type: string\\
Default: nil
}

{\footnotesize Оновлюється вміст навчання на основі ролей після подій}


{\footnotesize\bfseries
No: 3\\
Name: at\_3\_c\\
Type: string\\
Default: nil
}

{\footnotesize Інформація отримана з внутрішніх чи зовнішніх інцидентів або порушень безпеки, включається в навчання на основі ролей}


{\footnotesize\bfseries
No: 4\\
Name: at\_3\_odp\_01\\
Type: list\\
Default: [\textquotedbl{}admin\textquotedbl{}, \textquotedbl{}security\_officer\textquotedbl{}]
}

{\footnotesize Визначено ролі та обов'язки для тренінгів з безпеки на основі ролей}


{\footnotesize\bfseries
No: 5\\
Name: at\_3\_odp\_02\\
Type: list\\
Default: [\textquotedbl{}admin\textquotedbl{}, \textquotedbl{}security\_officer\textquotedbl{}]
}

{\footnotesize Визначено ролі та обов'язки для тренінгів з конфіденційності на основі ролей}


{\footnotesize\bfseries
No: 6\\
Name: at\_3\_odp\_03\\
Type: list\\
Default: [\textquotedbl{}admin\textquotedbl{}, \textquotedbl{}security\_officer\textquotedbl{}]
}

{\footnotesize Визначено частоту проведення тренінгів на основі ролей з безпеки та конфіденційності для призначеного персоналу після початкової підготовки}


{\footnotesize\bfseries
No: 7\\
Name: at\_3\_odp\_04\\
Type: integer\\
Default: 30
}

{\footnotesize Визначено частоту оновлення змісту навчання на основі ролей}


{\footnotesize\bfseries
No: 8\\
Name: at\_3\_odp\_05\\
Type: list\\
Default: [\textquotedbl{}admin\textquotedbl{}, \textquotedbl{}security\_officer\textquotedbl{}]
}

{\footnotesize Визначено події, які потребують оновлення змісту навчання на основі ролей; AT-03(a)[01][01] навчання з безпеки на основі ролей проводиться для ролей та обов'язків перед авторизацією доступу до системи, інформації або виконанням призначених обов'язків; AT-03(a)[01][02] навчання з конфіденційності на основі ролей проводиться для ролей та обов'язків перед авторизацією доступу до системи, інформації або виконанням призначених обов'язків; AT-03(a)[01][03] навчання з безпеки на основі ролей проводиться для ролей та обов'язків з частотою після цього; AT-03(a)[01][04] навчання з конфіденційності на основі ролей проводиться для ролей та обов'язків з частотою після цього; AT-03(a)[02][01] навчання з питань безпеки на основі ролей проводиться для персоналу, який виконує певні ролі та обов'язки у сфері безпеки, коли цього вимагають зміни в системі; AT-03(a)[02][02] навчання з питань конфіденційності на основі ролей проводиться для персоналу, який виконує певні ролі та обов'язки у сфері безпеки, коли цього вимагають зміни в системі}




\section{AU}
Клас заходів захисту AU --- АУДИТ ТА ПІДЗВІТНІСТЬ

Цей клас забезпечує можливість відстеження дій у системі, генерування, захист та аналіз записів аудиту для виявлення порушень політики безпеки.

\textbf{Перелік заходів захисту:} Політика та процедури аудиту та підзвітності (AU-1); Події аудиту (AU-2); Зміст записів аудиту (AU-3); Додаткова інформація про аудит (AU-3(1)); Реагування на відмови обробки даних аудиту (AU-5); Огляд, аналіз і звітність аудиту (AU-6); Зіставляння сховищ аудиту (AU-6(3)); Скорочення записів аудиту та формування звіту (AU-7); Позначка часу (AU-8); Захист інформації аудиту (AU-9); Доступ, який надається через членство в підмножини привілейованих користувачів (AU-9(4)); Збереження записів аудиту (AU-11); Генерація даних аудиту (AU-12).

\subsection{Політика та процедури аудиту та підзвітності (AU-1)}
a. Розробити, задокументувати та поширити [Призначення: серед персоналу або ролей, що їх визначила організація]:\\ 
1. [Вибір (один або декілька): Рівень організації; Рівень місії/бізнес-процесу; рівень системи] політика аудиту та підзвітності, яка:\\ 
(a) містить мету, сферу застосування, ролі, обов'язки, відповідальність керівництва, координацію між організаційними підрозділами та систему контролю відповідності (complaince);\\ 
(b) відповідає чинним законам, нормативним документам, директивам, нормам, політикам, стандартам та керівним документам.\\ 
2. Процедури, що сприяють здійсненню політики аудиту та підзвітності, а також пов'язані з ними заходи аудиту та підзвітності.\\ 
b. Призначити [Призначення: визначену організацією старшу посадову особу] для управління політикою та процедурами аудиту та підзвітності.\\ 
c. Переглядати та оновлювати поточний аудит та підзвітність:\\ 
1. політику [Призначення: частота, визначена організацією] та наступне [Призначення: події, визначені організацією];\\ 
2. процедури аудиту [Призначення: визначеною організацією частотою] та [Завдання: події, визначені організацією].

\vspace{10pt}
{\footnotesize\bfseries
No: 1\\
Name: au\_1\_a\_01\\
Type: list\\
Default: [\textquotedbl{}default\_deny\_rule\textquotedbl{}, \textquotedbl{}abac\_rule\_1\textquotedbl{}]
}

{\footnotesize Розроблено та задокументовано політику аудиту та підзвітності}


{\footnotesize\bfseries
No: 2\\
Name: au\_1\_a\_02\\
Type: list\\
Default: [\textquotedbl{}admin\textquotedbl{}, \textquotedbl{}security\_officer\textquotedbl{}]
}

{\footnotesize Політика аудиту та підзвітності доведена до персонал або ролі}


{\footnotesize\bfseries
No: 3\\
Name: au\_1\_a\_03\\
Type: list\\
Default: [\textquotedbl{}default\_deny\_rule\textquotedbl{}, \textquotedbl{}abac\_rule\_1\textquotedbl{}]
}

{\footnotesize Розроблені та задокументовані процедури аудиту та підзвітності, що сприяють впровадженню політики аудиту та підзвітності, а також відповідні заходи контролю аудиту та підзвітності}


{\footnotesize\bfseries
No: 4\\
Name: au\_1\_b\\
Type: list\\
Default: [\textquotedbl{}admin\textquotedbl{}, \textquotedbl{}security\_officer\textquotedbl{}]
}

{\footnotesize Посадова особа призначається для управління політикою та процедурами аудиту та підзвітності AU-01(c)[01][01] переглядається та оновлюється поточна політика аудиту та підзвітності з частота; AU-01(c)[01][02] переглядається та оновлюється поточна політика аудиту та підзвітності після подій; AU-01(c)[02][01] переглядається та оновлюється поточні процедури аудиту та підзвітності з частота; AU-01(c)[02][02] переглядається та оновлюється поточні процедури аудиту та підзвітності після подій}


{\footnotesize\bfseries
No: 5\\
Name: au\_1\_odp\_01\\
Type: list\\
Default: [\textquotedbl{}admin\textquotedbl{}, \textquotedbl{}security\_officer\textquotedbl{}]
}

{\footnotesize Визначено персонал або ролі, до яких має бути доведена політика аудиту та підзвітності}


{\footnotesize\bfseries
No: 6\\
Name: au\_1\_odp\_02\\
Type: list\\
Default: [\textquotedbl{}admin\textquotedbl{}, \textquotedbl{}security\_officer\textquotedbl{}]
}

{\footnotesize Визначено персонал або ролі, на які поширюються процедури аудиту та підзвітності}


{\footnotesize\bfseries
No: 7\\
Name: au\_1\_odp\_03\\
Type: string\\
Default: nil
}

{\footnotesize Вибрано одне або декілька з наступних ЗНАЧЕНЬ ПАРАМЕТРІВ: \{рівень організації; рівень місії/бізнес-процесу; рівень системи\}}


{\footnotesize\bfseries
No: 8\\
Name: au\_1\_odp\_04\\
Type: list\\
Default: [\textquotedbl{}admin\textquotedbl{}, \textquotedbl{}security\_officer\textquotedbl{}]
}

{\footnotesize Визначено посадову особу, яка управлятиме політикою та процедурами аудиту та підзвітності}


{\footnotesize\bfseries
No: 9\\
Name: au\_1\_odp\_05\\
Type: list\\
Default: [\textquotedbl{}default\_deny\_rule\textquotedbl{}, \textquotedbl{}abac\_rule\_1\textquotedbl{}]
}

{\footnotesize Визначено частоту, з якою переглядається та оновлюється поточна політика аудиту та підзвітності}


{\footnotesize\bfseries
No: 10\\
Name: au\_1\_odp\_06\\
Type: list\\
Default: [\textquotedbl{}default\_deny\_rule\textquotedbl{}, \textquotedbl{}abac\_rule\_1\textquotedbl{}]
}

{\footnotesize Визначено події, які потребують перегляду та оновлення поточної політики аудиту та підзвітності}


{\footnotesize\bfseries
No: 11\\
Name: au\_1\_odp\_07\\
Type: integer\\
Default: 30
}

{\footnotesize Визначено частоту, з якою переглядаються та оновлюються поточні процедури аудиту та підзвітності}


{\footnotesize\bfseries
No: 12\\
Name: au\_1\_odp\_08\\
Type: list\\
Default: [\textquotedbl{}login\textquotedbl{}, \textquotedbl{}logout\textquotedbl{}, \textquotedbl{}failed\_attempt\textquotedbl{}]
}

{\footnotesize Визначено події, які потребують перегляду та оновлення поточної процедури аудиту та підзвітності}




\subsection{ПОДІЇ АУДИТУ (AU-2)}
a. Визначити типи подій, які система може реєструвати для підтримки функції аудиту: [Призначення: типи подій, визначені організацією, які система здатна реєструвати];\\ 
b. Координувати функції аудиту безпеки з іншими організаційними підрозділами, які вимагають інформації, пов'язаної з аудитом, для посилення взаємної підтримки та допомоги у виборі типів подій, що перевіряються;\\ 
c. Визначити, які типи подій підлягають аудиту: [Призначення: визначені організацією події, що підлягають аудиту (підмножина подій, що підлягають аудиту, визначених в AU-2 a.), а також частота (або ситуація, що вимагає) проведення аудиту для кожної ідентифікованої події]\\ 
d. Обґрунтувати, чому типи подій, що перевіряються, вважаються достатніми для підтримки розслідувань інцидентів (постфактум), пов'язаних з безпекою та приватністю;\\ 
e. Перегляньте й оновіть типи подій, вибрані для журналювання [Призначення: частота, визначена організацією].

\vspace{10pt}
{\footnotesize\bfseries
No: 1\\
Name: au\_2\_a\\
Type: string\\
Default: nil
}

{\footnotesize Типи подій, які система здатна реєструвати, визначено для підтримки функції аудиту}


{\footnotesize\bfseries
No: 2\\
Name: au\_2\_b\\
Type: string\\
Default: nil
}

{\footnotesize Функція аудиту безпеки координується з іншими підрозділами організації, які вимагають інформації, пов'язаної з аудитом, длдля посилення взаємної підтримки та допомоги у виборі типів подій, що перевіряються}


{\footnotesize\bfseries
No: 3\\
Name: au\_2\_c\_01\\
Type: string\\
Default: nil
}

{\footnotesize Типи подій (підмножина 02\_ODP[01]) визначаються для реєстрації у системі; AU-}


{\footnotesize\bfseries
No: 4\\
Name: au\_2\_c\_02\\
Type: string\\
Default: \textquotedbl{}щорічно\textquotedbl{}
}

{\footnotesize Зазначені типи подій реєструються 02\_ODP[03] частота або ситуація>; <AU-}


{\footnotesize\bfseries
No: 5\\
Name: au\_2\_d\\
Type: string\\
Default: nil
}

{\footnotesize Надається обґрунтування, чому типи подій, що перевіряються, вважаються достатніми для підтримки розслідувань інцидентів (постфактум), пов'язаних з безпекою та конфіденційністю}


{\footnotesize\bfseries
No: 6\\
Name: au\_2\_e\\
Type: string\\
Default: \textquotedbl{}щорічно\textquotedbl{}
}

{\footnotesize Переглядаються та оновлюються типи подій, вибрані для реєстрації, частота. у системі}


{\footnotesize\bfseries
No: 7\\
Name: au\_2\_odp\_01\\
Type: string\\
Default: nil
}

{\footnotesize Визначено типи подій, які система може реєструвати для підтримки функції аудиту}


{\footnotesize\bfseries
No: 8\\
Name: au\_2\_odp\_02\\
Type: string\\
Default: nil
}

{\footnotesize Визначено типи подій (підмножина AU-02\_ODP[01]) що підлягають аудиту у системі}


{\footnotesize\bfseries
No: 9\\
Name: au\_2\_odp\_03\\
Type: list\\
Default: [\textquotedbl{}login\textquotedbl{}, \textquotedbl{}logout\textquotedbl{}, \textquotedbl{}failed\_attempt\textquotedbl{}]
}

{\footnotesize Визначено частоту або ситуацію, що вимагає проведення аудиту для кожної ідентифікованої події}


{\footnotesize\bfseries
No: 10\\
Name: au\_2\_odp\_04\\
Type: string\\
Default: \textquotedbl{}щорічно\textquotedbl{}
}

{\footnotesize Частота перегляду та оновлення типів подій, обраних для журналювання}




\subsection{ЗМІСТ ЗАПИСІВ АУДИТУ (AU-3)}
Переконатися, що записи аудиту містять інформацію, яка встановлює наступне:\\ 
a. який тип події стався;\\ 
b. коли відбулася подія;\\ 
c. де відбулася подія;\\ 
d. джерело події;\\ 
e. наслідки події;\\ 
f. результат події та ідентифікатор будь-яких осіб або суб'єктів, пов'язаних з подією.

\vspace{10pt}
{\footnotesize\bfseries
No: 1\\
Name: au\_3\_a\\
Type: list\\
Default: [\textquotedbl{}login\textquotedbl{}, \textquotedbl{}logout\textquotedbl{}, \textquotedbl{}failed\_attempt\textquotedbl{}]
}

{\footnotesize Записи аудиту містять інформацію, яка встановлює який тип події стався}


{\footnotesize\bfseries
No: 2\\
Name: au\_3\_b\\
Type: string\\
Default: nil
}

{\footnotesize Записи аудиту містять інформацію, яка встановлює коли подія сталася}


{\footnotesize\bfseries
No: 3\\
Name: au\_3\_c\\
Type: string\\
Default: nil
}

{\footnotesize Записи аудиту містять інформацію, яка встановлює де відбулася подія}


{\footnotesize\bfseries
No: 4\\
Name: au\_3\_d\\
Type: list\\
Default: [\textquotedbl{}login\textquotedbl{}, \textquotedbl{}logout\textquotedbl{}, \textquotedbl{}failed\_attempt\textquotedbl{}]
}

{\footnotesize Записи аудиту містять інформацію, яка встановлює джерело події}


{\footnotesize\bfseries
No: 5\\
Name: au\_3\_e\\
Type: list\\
Default: [\textquotedbl{}login\textquotedbl{}, \textquotedbl{}logout\textquotedbl{}, \textquotedbl{}failed\_attempt\textquotedbl{}]
}

{\footnotesize Записи аудиту містять інформацію, яка встановлює наслідки події}


{\footnotesize\bfseries
No: 6\\
Name: au\_3\_f\\
Type: list\\
Default: [\textquotedbl{}login\textquotedbl{}, \textquotedbl{}logout\textquotedbl{}, \textquotedbl{}failed\_attempt\textquotedbl{}]
}

{\footnotesize Записи аудиту містять інформацію, яка встановлює результат події та ідентифікатор будь-яких осіб або суб'єктів, пов'язаних з подією}




\subsubsection{ДОДАТКОВА ІНФОРМАЦІЯ ПРО АУДИТ (AU-3(1))}


\vspace{10pt}
{\footnotesize\bfseries
No: 1\\
Name: au\_3\_1\_01\\
Type: string\\
Default: nil
}

{\footnotesize ЗМІСТ ЗАПИСІВ АУДИТУ - ДОДАТКОВА ІНФОРМАЦІЯ ПРО АУДИТ МЕТА ОЦІНКИ: Визначити, чи:}


{\footnotesize\bfseries
No: 2\\
Name: au\_3\_1\_odp\\
Type: string\\
Default: nil
}

{\footnotesize Визначено додаткову інформацію, яка має бути включена до записів аудиту}




\subsection{РЕАГУВАННЯ НА ВІДМОВИ ОБРОБКИ ДАНИХ АУДИТУ (AU-5)}
a. Сповіщати [Призначення: визначені організацією персонал або посади] у разі збою обробки даних аудиту в [Призначення: визначений організацією період часу].\\ 
b. Виконати наступні додаткові дії: [Призначення: визначені організацією дії, які необхідно зробити].

\vspace{10pt}
{\footnotesize\bfseries
No: 1\\
Name: au\_5\_a\\
Type: list\\
Default: [\textquotedbl{}admin\textquotedbl{}, \textquotedbl{}security\_officer\textquotedbl{}]
}

{\footnotesize Персонал або ролі отримують сповіщення у разі збою процесу обробки даних аудиту періоду часу}


{\footnotesize\bfseries
No: 2\\
Name: au\_5\_b\\
Type: list\\
Default: [\textquotedbl{}login\textquotedbl{}, \textquotedbl{}logout\textquotedbl{}, \textquotedbl{}failed\_attempt\textquotedbl{}]
}

{\footnotesize Додаткові дії виконуються у разі збою процесу обробки даних аудиту}


{\footnotesize\bfseries
No: 3\\
Name: au\_5\_odp\_01\\
Type: list\\
Default: [\textquotedbl{}admin\textquotedbl{}, \textquotedbl{}security\_officer\textquotedbl{}]
}

{\footnotesize Визначено персонал або ролі, які отримують сповіщення про збої в процесі обробки даних аудиту }


{\footnotesize\bfseries
No: 4\\
Name: au\_5\_odp\_02\\
Type: list\\
Default: [\textquotedbl{}admin\textquotedbl{}, \textquotedbl{}security\_officer\textquotedbl{}]
}

{\footnotesize Визначено період часу, протягом якого персонал або ролі отримують сповіщення про збої в процесі обробки даних аудиту}


{\footnotesize\bfseries
No: 5\\
Name: au\_5\_odp\_03\\
Type: list\\
Default: [\textquotedbl{}login\textquotedbl{}, \textquotedbl{}logout\textquotedbl{}, \textquotedbl{}failed\_attempt\textquotedbl{}]
}

{\footnotesize Визначено додаткові дії, яких слід вжити у випадку збою в процесі обробки даних аудиту }




\subsection{ОГЛЯД, АНАЛІЗ І ЗВІТНІСТЬ АУДИТУ (AU-6)}
a. Переглядати та аналізувати записи системного аудиту [Призначення: з визначеною організацією частотою] для виявлення [Призначення: визначеної організацією неналежної або незвичайної діяльності].\\ 
b. Відправляти звіт про аудит [Призначення: визначеним організацією персоналу або посадам].\\ 
c. Налаштувати рівні огляду аудиту, аналізу та звітності в рамках системи, коли змінюється рівень ризику на основі інформації від правоохоронних органів, розвідувальної інформації або від інших достовірних джерел інформації.

\vspace{10pt}
{\footnotesize\bfseries
No: 1\\
Name: au\_6\_01\\
Type: string\\
Default: nil
}

{\footnotesize ОГЛЯД, АНАЛІЗ І ЗВІТНІСТЬ АУДИТУ МЕТА ОЦІНКИ: Визначити, чи:}


{\footnotesize\bfseries
No: 2\\
Name: au\_6\_a\\
Type: string\\
Default: \textquotedbl{}щотижня\textquotedbl{}
}

{\footnotesize Записи аудиту системи переглядаються та аналізуються частота для виявлення ознак неналежної або незвичної діяльності та потенційного впливу неналежної або незвичної діяльності}


{\footnotesize\bfseries
No: 3\\
Name: au\_6\_b\\
Type: list\\
Default: [\textquotedbl{}admin\textquotedbl{}, \textquotedbl{}security\_officer\textquotedbl{}]
}

{\footnotesize Звіт аудиту відправляється персоналу або ролям}


{\footnotesize\bfseries
No: 4\\
Name: au\_6\_c\\
Type: string\\
Default: nil
}

{\footnotesize Рівень перевірки, аналізу та звітування записів аудиту в системі коригується у разі зміни ризиків на основі інформації правоохоронних органів, розвідувальної інформації або інших достовірних джерел інформації}


{\footnotesize\bfseries
No: 5\\
Name: au\_6\_odp\_01\\
Type: integer\\
Default: 30
}

{\footnotesize Визначено частоту, з якою переглядаються та аналізуються записи аудиту системи}


{\footnotesize\bfseries
No: 6\\
Name: au\_6\_odp\_02\\
Type: string\\
Default: nil
}

{\footnotesize Визначена неналежна або незвична діяльність}


{\footnotesize\bfseries
No: 7\\
Name: au\_6\_odp\_03\\
Type: list\\
Default: [\textquotedbl{}admin\textquotedbl{}, \textquotedbl{}security\_officer\textquotedbl{}]
}

{\footnotesize Визначено персонал або ролі які отримують результати оглядів та аналізів системних записів}




\subsubsection{ЗІСТАВЛЯННЯ СХОВИЩ АУДИТУ (AU-6(3))}
Аналізувати та зіставляти записи аудиту в різних сховищах задля забезпечення ситуативної обізнаності в масштабах організації.

\vspace{10pt}
{\footnotesize\bfseries
No: 1\\
Name: au\_6\_3\_01\\
Type: string\\
Default: nil
}

{\footnotesize Аналізуєються та зіставляються записи аудиту в різних сховищах, задля забезпечення ситуативної обізнаності в масштабах організації}




\subsection{СКОРОЧЕННЯ ЗАПИСІВ АУДИТУ ТА ФОРМУВАННЯ ЗВІТУ (AU-7)}
Забезпечити та реалізувати можливості скорочення записів перевірок аудитом і звітів, до рівня, який:\\ 
a. підтримує перевірку, аналіз і звітність аудиту на вимогу та розслідування (постфактум) інцидентів безпеки;\\ 
b. не змінює оригінальний вміст або час упорядкування записів аудиту.

\vspace{10pt}
{\footnotesize\bfseries
No: 1\\
Name: au\_7\_a\_02\\
Type: string\\
Default: nil
}

{\footnotesize Реалізовано можливість скорочення записів перевірок аудитом та звітів, до рінвня що підтримує перевірку, аналіз і звітність аудиту на вимогу та розслідування (постфактум) інцидентів безпеки}


{\footnotesize\bfseries
No: 2\\
Name: au\_7\_b\_01\\
Type: integer\\
Default: 30
}

{\footnotesize Забезпечено можливість скорочення записів перевірок аудитом та звітів, які не змінюють оригінальний зміст або час упорядкування записів аудиту}


{\footnotesize\bfseries
No: 3\\
Name: au\_7\_b\_02\\
Type: integer\\
Default: 30
}

{\footnotesize Реалізовано можливість скорочення записів перевірок аудитом та звітів, які не змінюють оригінальний зміст або час упорядкування записів аудиту}




\subsection{ПОЗНАЧКА ЧАСУ (AU-8)}
a. Використовувати внутрішньосистемний годинник для створення позначок часу для записів аудиту.\\ 
b. Застосовувати позначки часу, які відповідають [Призначення: деталізація вимірювання часу, визначена організацією] і використовують всесвітній координований час, мають фіксоване зміщення місцевого часу відносно всесвітнього координованого часу або включають зміщення місцевого часу як частину позначки часу.

\vspace{10pt}
{\footnotesize\bfseries
No: 1\\
Name: au\_8\_a\\
Type: integer\\
Default: 30
}

{\footnotesize Внутрішній системний годинник використовується для створення позначок часу для записів аудиту}


{\footnotesize\bfseries
No: 2\\
Name: au\_8\_b\\
Type: integer\\
Default: 30
}

{\footnotesize Позначки часу застосовуються для записів аудиту, які відповідають деталізація вимірювання часу і які використовують всесвітній координований час, мають фіксоване місцеве зміщення місцевого часу від всесвітнього координованого часу або включають зміщення місцевого часу як частину позначки часу}


{\footnotesize\bfseries
No: 3\\
Name: au\_8\_odp\\
Type: integer\\
Default: 30
}

{\footnotesize Визначено деталізацію вимірювання часу для часових позначок записів аудиту}




\subsection{ЗАХИСТ ІНФОРМАЦІЇ АУДИТУ (AU-9)}
a. Захист інформації аудиту та інструментів несанкціонованого доступу, зміни та видалення; журналювання аудиту від\\ 
b. Сповіщення [Призначення: персонал або ролі, визначені організацією] у разі виявлення несанкціонованого доступу, зміни або видалення інформації аудиту.

\vspace{10pt}
{\footnotesize\bfseries
No: 1\\
Name: au\_9\_01\\
Type: string\\
Default: nil
}

{\footnotesize ЗАХИСТ ІНФОРМАЦІЇ АУДИТУ МЕТА ОЦІНКИ: Визначити, чи:}


{\footnotesize\bfseries
No: 2\\
Name: au\_9\_a\\
Type: string\\
Default: nil
}

{\footnotesize Інформація про аудит та інструменти журналювання аудиту захищені від несанкціонованого доступу, зміни та видалення}


{\footnotesize\bfseries
No: 3\\
Name: au\_9\_b\\
Type: list\\
Default: [\textquotedbl{}admin\textquotedbl{}, \textquotedbl{}security\_officer\textquotedbl{}]
}

{\footnotesize Персонал або ролі отримують сповіщення при виявленні несанкціонованого доступу, зміни або видалення інформації аудиту}


{\footnotesize\bfseries
No: 4\\
Name: au\_9\_odp\\
Type: list\\
Default: [\textquotedbl{}admin\textquotedbl{}, \textquotedbl{}security\_officer\textquotedbl{}]
}

{\footnotesize Визначено персонал або ролі, які мають бути сповіщені при виявленні несанкціонованого доступу, зміни або видалення інформації аудиту}




\subsubsection{ДОСТУП, ЯКИЙ НАДАЄТЬСЯ ЧЕРЕЗ ЧЛЕНСТВО В ПІДМНОЖИНИ ПРИВІЛЕЙОВАНИХ КОРИСТУВАЧІВ (AU-9(4))}
Авторизувати доступ до управління функціональністю аудиту тільки для [Призначення: визначеної організацією підмножини привілейованих користувачів].

\vspace{10pt}
{\footnotesize\bfseries
No: 1\\
Name: au\_9\_4\_01\\
Type: string\\
Default: nil
}

{\footnotesize Авторизувати доступ до управління функціональністю аудиту тільки для підмножини привілейованих користувачів}


{\footnotesize\bfseries
No: 2\\
Name: au\_9\_4\_odp\\
Type: string\\
Default: nil
}

{\footnotesize Визначено підмножину привілейованих користувачів}




\subsection{ЗБЕРЕЖЕННЯ ЗАПИСІВ АУДИТУ (AU-11)}
Зберігати записи аудиту впродовж [Призначення: визначеного організацією періоду часу, відповідно політиці зберігання записів], щоб забезпечити підтримку розслідувань (постфактум) інцидентів безпеки та приватності, а також для задоволення вимог нормативних і документів організації щодо збереження даних аудиту.

\vspace{10pt}
{\footnotesize\bfseries
No: 1\\
Name: au\_11\_01\\
Type: integer\\
Default: 30
}

{\footnotesize Записи аудиту зберігаються впродовж період часу, щоб забезпечити підтримку розслідування (постфактум) інцидентів безпеки та конфіденційності, а також відповідати нормативним та вимогам організації щодо збереження даних аудиту}


{\footnotesize\bfseries
No: 2\\
Name: au\_11\_odp\\
Type: list\\
Default: [\textquotedbl{}default\_deny\_rule\textquotedbl{}, \textquotedbl{}abac\_rule\_1\textquotedbl{}]
}

{\footnotesize Визначено період часу для зберігання записів аудиту, який узгоджується з політикою зберігання записів}




\subsection{ГЕНЕРАЦІЯ ДАНИХ АУДИТУ (AU-12)}
a. Забезпечити генерацію даних аудиту для типів подій, що перевіряються в AU-2а в [Призначення: визначених організацією компонентах системи].\\ 
b. Дозволити [Призначення: визначеному організацією персоналу або посадам] вибирати, які типи подій, що перевіряються, повинні перевірятися окремими компонентами системи;\\ 
c. Генерувати записи аудиту для типів подій, визначених в AU-2с. з вмістом згідно з AU-3.

\vspace{10pt}
{\footnotesize\bfseries
No: 1\\
Name: au\_12\_a\\
Type: string\\
Default: nil
}

{\footnotesize Можливість генерації записів аудиту для типів подій, які система здатна перевіряти (визначених у AU-02\_ODP[01]), забезпечується компонентами системи}


{\footnotesize\bfseries
No: 2\\
Name: au\_12\_b\\
Type: list\\
Default: [\textquotedbl{}admin\textquotedbl{}, \textquotedbl{}security\_officer\textquotedbl{}]
}

{\footnotesize Персонал або ролі може/можуть вибирати типи подій, які будуть реєструватися певними компонентами системи; AU-12(c) згенеровано записи аудиту для типів подій, визначених у AU02\_ODP[02], які включають вміст записів аудиту, визначений у AU03}


{\footnotesize\bfseries
No: 3\\
Name: au\_12\_odp\_01\\
Type: string\\
Default: nil
}

{\footnotesize Визначено компоненти системи, які забезпечують можливість генерації записів аудиту для типів подій (визначених у AU02\_ODP[02])}


{\footnotesize\bfseries
No: 4\\
Name: au\_12\_odp\_02\\
Type: list\\
Default: [\textquotedbl{}admin\textquotedbl{}, \textquotedbl{}security\_officer\textquotedbl{}]
}

{\footnotesize Визначено персонал або ролі, яким дозволено обирати типи подій, що мають реєструватися певними компонентами системи}




\section{CA}
Клас заходів захисту CA --- ОЦІНЮВАННЯ,

Цей клас регламентує процеси перевірки ефективності заходів захисту, авторизації систем та їх безперервного моніторингу.

\textbf{Перелік заходів захисту:} Політика і процедури оцінювання, акредитація та моніторинг безпеки (CA-1); Оцінювання (CA-2); Взаємодія систем (CA-3); План усунення недоліків та контрольні показники (CA-5); Безперервний моніторинг (CA-7).

\subsection{ПОЛІТИКА І ПРОЦЕДУРИ ОЦІНЮВАННЯ, АКРЕДИТАЦІЯ ТА МОНІТОРИНГ БЕЗПЕКИ (CA-1)}
a. Розробити, задокументувати та поширити серед [Призначення: визначеного організацією персоналу або посад]:\\ 
1. [Вибір (один або декілька): Рівень організації; Рівень місії/бізнес-процесу; рівень системи] політика оцінювання, авторизації та моніторингу, яка:\\ 
(a) містить мету, сферу застосування, ролі, обов'язки, відповідальність керівництва, координацію між організаційними підрозділами та систему контролю відповідності (complaince);\\ 
(b) відповідає чинним законам, нормативним документам, наказам, положенням, політиці, стандартам і керівним принципам.\\ 
2. Процедури, що сприяють реалізації політики оцінювання, авторизації та моніторингу безпеки та приватності, а також пов'язаних з ними заходів оцінювання, авторизації та моніторингу безпеки та приватності.\\ 
b. Призначити [Призначення: посадова особа, визначена організацією] для управління розробкою, документуванням і розповсюдженням політики та процедур оцінювання, авторизації та моніторингу;\\ 
c. Переглядати та оновлювати поточне оцінювання, авторизацію та моніторинг:\\ 
1. Політику [Призначення: частота, визначена організацією] та наступне [Призначення: події, визначені організацією];\\ 
2. Процедури [Призначення: частота, визначена організацією] та наступні [Призначення: події, визначені організацією].

\vspace{10pt}
{\footnotesize\bfseries
No: 1\\
Name: ca\_1\_a\_01\\
Type: list\\
Default: [\textquotedbl{}default\_deny\_rule\textquotedbl{}, \textquotedbl{}abac\_rule\_1\textquotedbl{}]
}

{\footnotesize Розроблено та задокументовано політику оцінювання, авторизації та моніторингу}


{\footnotesize\bfseries
No: 2\\
Name: ca\_1\_a\_02\\
Type: list\\
Default: [\textquotedbl{}admin\textquotedbl{}, \textquotedbl{}security\_officer\textquotedbl{}]
}

{\footnotesize Політика оцінювання, авторизації та моніторингу поширюється серед персоналу або ролей}


{\footnotesize\bfseries
No: 3\\
Name: ca\_1\_a\_03\\
Type: list\\
Default: [\textquotedbl{}default\_deny\_rule\textquotedbl{}, \textquotedbl{}abac\_rule\_1\textquotedbl{}]
}

{\footnotesize Розроблені та задокументовані процедури оцінювання, авторизації та моніторингу, що сприяють впровадженню політики оцінювання, авторизації та моніторингу, а також пов'язані з ними засоби контролю оцінювання, авторизації та моніторингу}


{\footnotesize\bfseries
No: 4\\
Name: ca\_1\_b\\
Type: list\\
Default: [\textquotedbl{}admin\textquotedbl{}, \textquotedbl{}security\_officer\textquotedbl{}]
}

{\footnotesize Посадова особа призначається для управління розробкою, документуванням та розповсюдженням політики та процедур оцінювання, авторизації та моніторингу; CA-01(c)[01][01] переглядається та оновлюється поточна політика оцінки, авторизації та моніторингу з частотою; CA-01(c)[01][02] переглядається та оновлюється поточна політика оцінки, авторизації та моніторингу після подій; CA-01(c)[02][01] переглядаються та оновлюються поточні процедури оцінки, авторизації та моніторингу з частотою; CA-01(c)[02][02] переглядаються та оновлюються поточні процедури оцінки, авторизації та моніторингу після подій; ПОТЕНЦІЙНІ МЕТОДИ ТА ОБ'ЄКТИ ОЦІНКИ:}


{\footnotesize\bfseries
No: 5\\
Name: ca\_1\_odp\_01\\
Type: list\\
Default: [\textquotedbl{}admin\textquotedbl{}, \textquotedbl{}security\_officer\textquotedbl{}]
}

{\footnotesize Визначено персонал або ролі, серед яких має бути поширена політика оцінювання, авторизації та моніторингу}


{\footnotesize\bfseries
No: 6\\
Name: ca\_1\_odp\_02\\
Type: list\\
Default: [\textquotedbl{}admin\textquotedbl{}, \textquotedbl{}security\_officer\textquotedbl{}]
}

{\footnotesize Визначено персонал або ролі, серед яких мають бути поширені процедури оцінювання, авторизації та моніторингу}


{\footnotesize\bfseries
No: 7\\
Name: ca\_1\_odp\_03\\
Type: string\\
Default: nil
}

{\footnotesize Вибрано одне або декілька з наступних ЗНАЧЕНЬ ПАРАМЕТРІВ: \{рівень організації; рівень місії/бізнес-процесу; рівень системи\}}


{\footnotesize\bfseries
No: 8\\
Name: ca\_1\_odp\_04\\
Type: list\\
Default: [\textquotedbl{}admin\textquotedbl{}, \textquotedbl{}security\_officer\textquotedbl{}]
}

{\footnotesize Визначено посадову особу, яка управлятиме розробкою, документуванням і розповсюдженням політики та процедур оцінювання, авторизації та моніторингу}


{\footnotesize\bfseries
No: 9\\
Name: ca\_1\_odp\_05\\
Type: list\\
Default: [\textquotedbl{}default\_deny\_rule\textquotedbl{}, \textquotedbl{}abac\_rule\_1\textquotedbl{}]
}

{\footnotesize Визначено частоту, з якою переглядається та оновлюється поточна політика оцінювання, авторизації та моніторингу}


{\footnotesize\bfseries
No: 10\\
Name: ca\_1\_odp\_06\\
Type: list\\
Default: [\textquotedbl{}default\_deny\_rule\textquotedbl{}, \textquotedbl{}abac\_rule\_1\textquotedbl{}]
}

{\footnotesize Визначено події, які потребують перегляду та оновлення поточної політики оцінки, авторизації та моніторингу}


{\footnotesize\bfseries
No: 11\\
Name: ca\_1\_odp\_07\\
Type: integer\\
Default: 30
}

{\footnotesize Визначено частоту, з якою переглядається та оновлюється поточні процедури оцінювання, авторизації та моніторингу}


{\footnotesize\bfseries
No: 12\\
Name: ca\_1\_odp\_08\\
Type: list\\
Default: [\textquotedbl{}login\textquotedbl{}, \textquotedbl{}logout\textquotedbl{}, \textquotedbl{}failed\_attempt\textquotedbl{}]
}

{\footnotesize Визначено події, які потребують перегляду та оновлення поточні процедури оцінки, авторизації та моніторингу}




\subsection{ОЦІНЮВАННЯ (CA-2)}
a. Виберіть відповідного оцінювача або команду з оцінки для типу оцінювання, яке буде проводитися;\\ 
b. Розробіть план контрольної оцінки, який описує обсяг оцінки, в тому числі:\\ 
1. заходи захисту та посилені заходи, що підлягають оцінюванню;\\ 
2. процедури оцінювання, ефективності заходів; які використовуватимуться для визначення\\ 
3. середовище оцінювання, групу оцінювання, ролі й обов'язки з оцінювання.\\ 
c. Забезпечити розгляд і затвердження плану оцінювання уповноваженою офіційною особою або призначеним для проведення оцінювання представником;\\ 
d. Оцінити заходи захисту в системі та в її середовищі функціонування з [Призначення: визначеною організацією частотою] для визначення, наскільки коректно реалізовані заходи безпеки і чи працюють вони за призначенням і дають бажаний результат щодо дотримання встановлених вимог безпеки та приватності;\\ 
e. Підготовити звіт оцінювання безпеки, який документує результати оцінювання;\\ 
f. Надати результати оцінювання з безпеки [Призначення: особам або ролям, визначеним організацією].

\vspace{10pt}
{\footnotesize\bfseries
No: 1\\
Name: ca\_2\_a\\
Type: string\\
Default: nil
}

{\footnotesize Обрано відповідного оцінювача або команду оцінювачів для проведення оцінювання}


{\footnotesize\bfseries
No: 2\\
Name: ca\_2\_b\_01\\
Type: list\\
Default: [\textquotedbl{}admin\textquotedbl{}, \textquotedbl{}security\_officer\textquotedbl{}]
}

{\footnotesize Розроблено план контрольної оцінки, який описує обсяг оцінки, включаючи заходи захисту та посилені заходи, що підлягають оцінюванню. CA-02(b)[02] розроблено план контрольної оцінки, який описує обсяг оцінки, включаючи процедури оцінювання, які використовуватимуться для визначення ефективності заходів. CA-02(b)[03][01] розроблено план контрольної оцінки, який описує обсяг оцінки, включаючи середовище оцінювання. CA-02(b)[03][02] розроблено план контрольної оцінки, який описує обсяг оцінки, включаючи групу оцінювання. CA-02(b)[03][03] розроблено план контрольної оцінки, який описує обсяг оцінки, включаючи ролі й обов'язки з оцінювання}


{\footnotesize\bfseries
No: 3\\
Name: ca\_2\_b\_02\\
Type: list\\
Default: [\textquotedbl{}admin\textquotedbl{}, \textquotedbl{}security\_officer\textquotedbl{}]
}

{\footnotesize Розроблено план контрольної оцінки, який описує обсяг оцінки, включаючи процедури оцінювання, які використовуватимуться для визначення ефективності заходів. CA-02(b)[03][01] розроблено план контрольної оцінки, який описує обсяг оцінки, включаючи середовище оцінювання. CA-02(b)[03][02] розроблено план контрольної оцінки, який описує обсяг оцінки, включаючи групу оцінювання. CA-02(b)[03][03] розроблено план контрольної оцінки, який описує обсяг оцінки, включаючи ролі й обов'язки з оцінювання}


{\footnotesize\bfseries
No: 4\\
Name: ca\_2\_c\\
Type: list\\
Default: [\textquotedbl{}admin\textquotedbl{}, \textquotedbl{}security\_officer\textquotedbl{}]
}

{\footnotesize План оцінки заходів захисту розглядається та затверджується уповноваженою посадовою особою або призначеним представником перед проведенням оцінки}


{\footnotesize\bfseries
No: 5\\
Name: ca\_2\_d\_01\\
Type: string\\
Default: \textquotedbl{}щорічно\textquotedbl{}
}

{\footnotesize Заходи захисту оцінюються в системі та середовищі її функціонування частота оцінки, щоб визначити, наскільки коректно реалізовані заходи захисту, чи працюють вони за призначенням і чи дають бажаний результат щодо дотримання встановлених вимог до безпеки}


{\footnotesize\bfseries
No: 6\\
Name: ca\_2\_d\_02\\
Type: string\\
Default: \textquotedbl{}щорічно\textquotedbl{}
}

{\footnotesize Заходи захисту оцінюються в системі та середовищі її функціонування частота оцінки, щоб визначити, наскільки коректно реалізовані заходи захисту, чи працюють вони за призначенням і чи дають бажаний результат щодо дотримання встановлених вимог до конфіденційності; CA-02(e) готується звіт оцінювання}


{\footnotesize\bfseries
No: 7\\
Name: ca\_2\_f\\
Type: list\\
Default: [\textquotedbl{}admin\textquotedbl{}, \textquotedbl{}security\_officer\textquotedbl{}]
}

{\footnotesize Результати оцінювання з безпеки надаються особам або ролям. оцінювання , який документує результати}


{\footnotesize\bfseries
No: 8\\
Name: ca\_2\_odp\_01\\
Type: integer\\
Default: 30
}

{\footnotesize Визначено частоту, з якою слід оцінювати засоби контролю в системі та середовищі її функціонування}


{\footnotesize\bfseries
No: 9\\
Name: ca\_2\_odp\_02\\
Type: list\\
Default: [\textquotedbl{}admin\textquotedbl{}, \textquotedbl{}security\_officer\textquotedbl{}]
}

{\footnotesize Визначені особи або ролі, яким мають бути надані результати оцінювання з безпеки}




\subsection{ВЗАЄМОДІЯ СИСТЕМ (CA-3)}
a. схвалити та керувати обміном інформацією між системою та іншими системами за допомогою [Вибір (один або кілька): угоди безпеки взаємозв'язку; договори безпеки обміну інформацією; меморандуми про взаєморозуміння; угоди про рівень обслуговування; угоди користувача; угоди про нерозголошення; [Доручення: тип договору, визначений організацією]];.\\ 
b. документувати, як частину угоди про обмін, характеристики інтерфейсу, вимоги до безпеки та приватності, засоби контролю та відповідальність для кожної системи, а також характер переданої інформації;\\ 
c. здійснювати перегляд та оновлення угод з [Призначення: визначеною організацією частотою].

\vspace{10pt}
{\footnotesize\bfseries
No: 1\\
Name: ca\_3\_b\_01\\
Type: integer\\
Default: 30
}

{\footnotesize Характеристики інтерфейсу задокументовані як частина кожної угоди про обмін}


{\footnotesize\bfseries
No: 2\\
Name: ca\_3\_b\_02\\
Type: integer\\
Default: 30
}

{\footnotesize Вимоги до безпеки задокументовані як частина кожної угоди про обмін}


{\footnotesize\bfseries
No: 3\\
Name: ca\_3\_b\_03\\
Type: integer\\
Default: 30
}

{\footnotesize Вимоги щодо конфіденційності задокументовані як частина кожної угоди про обмін}


{\footnotesize\bfseries
No: 4\\
Name: ca\_3\_b\_04\\
Type: integer\\
Default: 30
}

{\footnotesize Заходи захисту задокументовані як частина кожної угоди про обмін}


{\footnotesize\bfseries
No: 5\\
Name: ca\_3\_b\_05\\
Type: integer\\
Default: 30
}

{\footnotesize Відповідальність за кожну систему задокументована як частина кожної угоди про обмін}


{\footnotesize\bfseries
No: 6\\
Name: ca\_3\_b\_06\\
Type: integer\\
Default: 30
}

{\footnotesize Характер переданої інформації документується як частина кожної угоди про обмін}


{\footnotesize\bfseries
No: 7\\
Name: ca\_3\_c\\
Type: string\\
Default: \textquotedbl{}щорічно\textquotedbl{}
}

{\footnotesize Угоди переглядаються та оновлюються частота}


{\footnotesize\bfseries
No: 8\\
Name: ca\_3\_odp\_01\\
Type: string\\
Default: nil
}

{\footnotesize Вибрано одне або декілька з наступних ЗНАЧЕНЬ ПАРАМЕТРІВ: \{): угоди безпеки взаємозв'язку; договори безпеки обміну інформацією; меморандуми про взаєморозуміння; угоди про рівень обслуговування; угоди користувача; угоди про нерозголошення; тип угоди\}}


{\footnotesize\bfseries
No: 9\\
Name: ca\_3\_odp\_02\\
Type: string\\
Default: nil
}

{\footnotesize Визначено тип угоди, який використовується для схвалення та керування обміном інформацією (якщо вибрано)}


{\footnotesize\bfseries
No: 10\\
Name: ca\_3\_odp\_03\\
Type: integer\\
Default: 30
}

{\footnotesize Визначено частоту, з якою необхідно переглядати та оновлювати угоди}




\subsection{ПЛАН УСУНЕННЯ НЕДОЛІКІВ ТА КОНТРОЛЬНІ ПОКАЗНИКИ (CA-5)}
a. Розробити для системи план усунення недоліків та контрольні показники з метою документування запланованих коригувальних дій організації для усунення недоліків і зауважень, які виявлені в ході оцінювання заходів захисту, а також для зменшення або усунення відомих вразливостей у системі.\\ 
b. Оновлювати чинний план усунення недоліків та контрольні показники з [Призначення: визначеною організацією частотою] на основі результатів оцінювання заходів, незалежних аудитів та постійного моніторингу.

\vspace{10pt}
{\footnotesize\bfseries
No: 1\\
Name: ca\_5\_a\\
Type: list\\
Default: [\textquotedbl{}login\textquotedbl{}, \textquotedbl{}logout\textquotedbl{}, \textquotedbl{}failed\_attempt\textquotedbl{}]
}

{\footnotesize Розроблено план усунення недоліків та контрольні показники для системи, та задокументувано заплановані дії організації з коригування, спрямовані на усунення недоліків та зауважень, виявлених під час оцінки заходів захисту, а також на зменшення або усунення відомих вразливостей в системі}


{\footnotesize\bfseries
No: 2\\
Name: ca\_5\_b\\
Type: string\\
Default: \textquotedbl{}щорічно\textquotedbl{}
}

{\footnotesize Існуючий план оновлюються частота на основі результатів оцінювання заходів, незалежних аудитів та постійного моніторингу}


{\footnotesize\bfseries
No: 3\\
Name: ca\_5\_odp\\
Type: integer\\
Default: 30
}

{\footnotesize Визначено частоту оновлення чинного плану усунення недоліків та контрольних показників на основі результатів оцінювання заходів захисту, незалежних аудитів та постійного моніторингу}




\subsection{БЕЗПЕРЕРВНИЙ МОНІТОРИНГ (CA-7)}
Розробити стратегію безперервного моніторингу безпеки та приватності й упровадити програму безперервного моніторингу безпеки та приватності, яка охоплює:\\ 
a. встановлення показників безпеки та приватності, які необхідно відстежувати: [Призначення: визначені організацією метрики];\\ 
b. встановлення [Призначення: визначена організацією частота] для моніторингу та [Призначення: визначена організацією частота] для безперервного оцінювання ефективності заходів захисту;\\ 
c. поточні оцінювання заходів захисту відповідно до стратегії безперервного моніторингу організації;\\ 
d. постійний моніторинг стану безпеки та приватності відповідно до встановлених організацією метрик і відповідно до стратегії безперервного моніторингу організації;\\ 
e. зіставлення та аналіз інформації, отриманої в результаті оцінювання та моніторингу безпеки та приватності;\\ 
f. дії реагування за результатами аналізу інформації, пов'язаної з безпекою та приватністю;\\ 
g. повідомлення про статус безпеки та приватності системи [Призначення: визначені організацією персонал або ролі] з [Призначення: визначеною організацією частотою].

\vspace{10pt}
{\footnotesize\bfseries
No: 1\\
Name: ca\_7\_01\\
Type: string\\
Default: nil
}

{\footnotesize Розроблено стратегію безперервного моніторингу на системному рівні; CA-07[02] безперервний моніторинг на рівні системи здійснюється відповідно до стратегії безперервного моніторингу на рівні організації}


{\footnotesize\bfseries
No: 2\\
Name: ca\_7\_a\\
Type: string\\
Default: nil
}

{\footnotesize Безперервний моніторинг на рівні системи включає встановлення наступних метрик на рівні системи, які підлягають моніторингу: метрики системного рівня}


{\footnotesize\bfseries
No: 3\\
Name: ca\_7\_b\_01\\
Type: integer\\
Default: 30
}

{\footnotesize Безперервний моніторинг на рівні системи включає встановлені частоти для моніторингу ефективності заходів захисту}


{\footnotesize\bfseries
No: 4\\
Name: ca\_7\_b\_02\\
Type: integer\\
Default: 30
}

{\footnotesize Безперервний моніторинг на рівні системи включає встановлені частоти для оцінки ефективності заходів захисту}


{\footnotesize\bfseries
No: 5\\
Name: ca\_7\_c\\
Type: string\\
Default: nil
}

{\footnotesize Безперервний моніторинг на рівні системи включає поточні контрольні оцінки відповідно до стратегії безперервного моніторингу}


{\footnotesize\bfseries
No: 6\\
Name: ca\_7\_d\\
Type: string\\
Default: nil
}

{\footnotesize Безперервний моніторинг на рівні системи включає постійний моніторинг визначених системою та організацією показників відповідно до стратегії безперервного моніторингу}


{\footnotesize\bfseries
No: 7\\
Name: ca\_7\_e\\
Type: string\\
Default: nil
}

{\footnotesize Безперервний моніторинг на рівні системи включає зіставлення та аналіз інформації, отриманої в результаті оцінювання та моніторингу}


{\footnotesize\bfseries
No: 8\\
Name: ca\_7\_f\\
Type: list\\
Default: [\textquotedbl{}login\textquotedbl{}, \textquotedbl{}logout\textquotedbl{}, \textquotedbl{}failed\_attempt\textquotedbl{}]
}

{\footnotesize Безперервний моніторинг на рівні системи включає в себе дії з реагування на результати аналізу інформації, пов'язаної з безпекою та приватністю}


{\footnotesize\bfseries
No: 9\\
Name: ca\_7\_g\_01\\
Type: list\\
Default: [\textquotedbl{}admin\textquotedbl{}, \textquotedbl{}security\_officer\textquotedbl{}]
}

{\footnotesize Безперервний моніторинг на рівні системи включає повідомлення про статус безпеки системи для персоналу або ролей частота}


{\footnotesize\bfseries
No: 10\\
Name: ca\_7\_g\_02\\
Type: list\\
Default: [\textquotedbl{}admin\textquotedbl{}, \textquotedbl{}security\_officer\textquotedbl{}]
}

{\footnotesize Безперервний моніторинг на рівні системи включає повідомлення про стан конфіденційності системи персоналу або ролям частота}


{\footnotesize\bfseries
No: 11\\
Name: ca\_7\_odp\_01\\
Type: string\\
Default: nil
}

{\footnotesize Визначено метрики системного рівня, які підлягають моніторингу}


{\footnotesize\bfseries
No: 12\\
Name: ca\_7\_odp\_02\\
Type: integer\\
Default: 30
}

{\footnotesize Визначено частоту, з якою слід моніторити ефективність заходів захисту}


{\footnotesize\bfseries
No: 13\\
Name: ca\_7\_odp\_03\\
Type: integer\\
Default: 30
}

{\footnotesize Визначено частоту, з якою слід оцінювати ефективність заходів захисту}


{\footnotesize\bfseries
No: 14\\
Name: ca\_7\_odp\_04\\
Type: list\\
Default: [\textquotedbl{}admin\textquotedbl{}, \textquotedbl{}security\_officer\textquotedbl{}]
}

{\footnotesize Визначено персонал або ролі, яким повідомляється про стан безпеки системи}


{\footnotesize\bfseries
No: 15\\
Name: ca\_7\_odp\_05\\
Type: integer\\
Default: 30
}

{\footnotesize Визначено частоту, з якою повідомляється про стан безпеки системи}


{\footnotesize\bfseries
No: 16\\
Name: ca\_7\_odp\_06\\
Type: list\\
Default: [\textquotedbl{}admin\textquotedbl{}, \textquotedbl{}security\_officer\textquotedbl{}]
}

{\footnotesize Визначено персонал або ролі, яким повідомляється про стан конфіденційності системи}


{\footnotesize\bfseries
No: 17\\
Name: ca\_7\_odp\_07\\
Type: integer\\
Default: 30
}

{\footnotesize Визначено частоту, з якою повідомляється про стан конфіденційності системи}




\section{CM}
Клас заходів захисту CM --- УПРАВЛІННЯ КОНФІГУРАЦІЄЮ

Цей клас гарантує створення та підтримання базових налаштувань системи, а також строгий контроль за змінами в апаратному та програмному забезпеченні.

\textbf{Перелік заходів захисту:} Політика та процедури управління конфігурацією (CM-1); Базова конфігурація (CM-2); Конфігурація систем та компонентів для сфер з високим ризиком (CM-2(7)); Управління змінами конфігурації (CM-3); Аналіз впливу на безпеку та приватність (CM-4); Верифікація функцій безпеки та приватності (CM-4(2)); Обмеження доступу до зміни (CM-5); Налаштування конфігурації (CM-6); Мінімально необхідна функціональність (CM-7); Періодичний перегляд (CM-7(1)); Авторизоване програмне забезпечення -- білий список (CM-7(5)); Інвентаризація компонентів системи (CM-8); Оновлення під час встановлення та видалення (CM-8(1)); Розташування інформації (CM-12).

\subsection{Політика та процедури управління конфігурацією (CM-1)}
a. Розробити, задокументувати та поширити серед [Призначення: визначених організацією персоналу або ролей]:\\ 
1. [Вибір (один або декілька): Рівень організації; Рівень місії/бізнес-процесу; рівень системи] політики управління конфігурацією, яка:\\ 
2.\\ 
(a) містить мету, сферу застосування, ролі, обов'язки, відповідальність керівництва, координацію між організаційними підрозділами та систему контролю відповідності (complaince);\\ 
(b) відповідає чинним законам, нормативним документам, наказам, положенням, політикам, стандартам і керівним принципам; процедури, що сприяють реалізації політики управління конфігурацією та пов'язаних з нею заходів управління конфігурацією.\\ 
b. Призначити [Призначення: посадова особа, визначена організацією] для управління розробкою, документуванням і розповсюдженням політики та процедур керування конфігурацією.\\ 
c. Переглядати та оновлювати поточну політику управління конфігурацією:\\ 
1. Політика [Призначення: частота, визначена організацією] та наступні [Призначення: події, визначені організацією];\\ 
2. Процедури [Призначення: частота, визначена організацією] та наступні [Призначення: події, визначені організацією].

\vspace{10pt}
{\footnotesize\bfseries
No: 1\\
Name: cm\_1\_a\_01\\
Type: list\\
Default: [\textquotedbl{}default\_deny\_rule\textquotedbl{}, \textquotedbl{}abac\_rule\_1\textquotedbl{}]
}

{\footnotesize Розроблено та задокументовано політику управління конфігурацією}


{\footnotesize\bfseries
No: 2\\
Name: cm\_1\_a\_02\\
Type: list\\
Default: [\textquotedbl{}admin\textquotedbl{}, \textquotedbl{}security\_officer\textquotedbl{}]
}

{\footnotesize Політика управління конфігурацією поширюється на персонал або ролі}


{\footnotesize\bfseries
No: 3\\
Name: cm\_1\_a\_03\\
Type: list\\
Default: [\textquotedbl{}default\_deny\_rule\textquotedbl{}, \textquotedbl{}abac\_rule\_1\textquotedbl{}]
}

{\footnotesize Розроблені та задокументовані процедури управління , що сприяють реалізації політики управління конфігурацією та пов'язаних з нею заходів управління конфігурацією}


{\footnotesize\bfseries
No: 4\\
Name: cm\_1\_b\\
Type: list\\
Default: [\textquotedbl{}admin\textquotedbl{}, \textquotedbl{}security\_officer\textquotedbl{}]
}

{\footnotesize Посадова особа призначається для управління розробкою, документуванням і розповсюдженням політики та процедур керування конфігурацією; CM-01(c)[01][01] переглядається та оновлюється поточна політика управління конфігурацією частота; CM-01(c)[01][02] переглядається та оновлюється поточна політика управління конфігурацією після події; CM-01(c)[02][01] переглядаються та оновлюються поточні процедури управління конфігурацією частота; CM-01(c)[02][02] переглядаються та оновлюються поточні процедури управління конфігурацією після події; ЗНАЧЕННЯ конфігурацією}


{\footnotesize\bfseries
No: 5\\
Name: cm\_1\_odp\_01\\
Type: list\\
Default: [\textquotedbl{}admin\textquotedbl{}, \textquotedbl{}security\_officer\textquotedbl{}]
}

{\footnotesize Визначено персонал або ролі, на яких поширюється політика управління конфігурацією}


{\footnotesize\bfseries
No: 6\\
Name: cm\_1\_odp\_02\\
Type: list\\
Default: [\textquotedbl{}admin\textquotedbl{}, \textquotedbl{}security\_officer\textquotedbl{}]
}

{\footnotesize Визначено персонал або ролі, на яких поширюється процедури управління конфігурацією}


{\footnotesize\bfseries
No: 7\\
Name: cm\_1\_odp\_03\\
Type: string\\
Default: nil
}

{\footnotesize Вибрано одне або декілька з наступних ЗНАЧЕНЬ ПАРАМЕТРІВ: \{рівень організації; рівень місії/бізнес-процесу; рівень системи\}}


{\footnotesize\bfseries
No: 8\\
Name: cm\_1\_odp\_04\\
Type: list\\
Default: [\textquotedbl{}admin\textquotedbl{}, \textquotedbl{}security\_officer\textquotedbl{}]
}

{\footnotesize Визначено посадову особу, яка управляє розробкою, документуванням і розповсюдженням політики та процедур керування конфігурацією}


{\footnotesize\bfseries
No: 9\\
Name: cm\_1\_odp\_05\\
Type: list\\
Default: [\textquotedbl{}default\_deny\_rule\textquotedbl{}, \textquotedbl{}abac\_rule\_1\textquotedbl{}]
}

{\footnotesize Визначено частоту, з якою переглядається та оновлюється поточна політика управління конфігурацією}


{\footnotesize\bfseries
No: 10\\
Name: cm\_1\_odp\_06\\
Type: list\\
Default: [\textquotedbl{}default\_deny\_rule\textquotedbl{}, \textquotedbl{}abac\_rule\_1\textquotedbl{}]
}

{\footnotesize Визначено події, після яких переглядається та оновлюється поточна політика управління конфігурацією}


{\footnotesize\bfseries
No: 11\\
Name: cm\_1\_odp\_07\\
Type: integer\\
Default: 30
}

{\footnotesize Визначено частоту, з якою переглядаються та оновлюються поточні процедури управління конфігурацією}


{\footnotesize\bfseries
No: 12\\
Name: cm\_1\_odp\_08\\
Type: list\\
Default: [\textquotedbl{}login\textquotedbl{}, \textquotedbl{}logout\textquotedbl{}, \textquotedbl{}failed\_attempt\textquotedbl{}]
}

{\footnotesize Визначено події, після яких переглядаються та оновлюються поточні процедури управління конфігурацією}




\subsection{БАЗОВА КОНФІГУРАЦІЯ (CM-2)}
a. Розробити, задокументувати та підтримувати за допомогою заходів конфігурації поточні базові налаштування системи.\\ 
b. Переглядати та оновлювати базові налаштування системи:\\ 
1. з [Призначення: визначеною організацією частотою];\\ 
2. за потреби внаслідок [Призначення: визначених організацією обставин];\\ 
3. коли встановлені нові або оновлені компоненти системи.

\vspace{10pt}
{\footnotesize\bfseries
No: 1\\
Name: cm\_2\_a\_01\\
Type: string\\
Default: nil
}

{\footnotesize Розроблено та задокументовано поточні базові налаштування системи}


{\footnotesize\bfseries
No: 2\\
Name: cm\_2\_a\_02\\
Type: string\\
Default: nil
}

{\footnotesize Поточні базові налаштування системи підтримуються за допомогою заходів конфігурації}


{\footnotesize\bfseries
No: 3\\
Name: cm\_2\_b\_01\\
Type: string\\
Default: \textquotedbl{}щорічно\textquotedbl{}
}

{\footnotesize Переглядаються та оновлюються базові налаштування системи частота}


{\footnotesize\bfseries
No: 4\\
Name: cm\_2\_b\_02\\
Type: string\\
Default: nil
}

{\footnotesize Переглядаються та оновлюються базові налаштування системи після подій}


{\footnotesize\bfseries
No: 5\\
Name: cm\_2\_b\_03\\
Type: string\\
Default: nil
}

{\footnotesize Переглядаються та оновлюються базові налаштування системи коли встановлюються або модернізуються компоненти системи}


{\footnotesize\bfseries
No: 6\\
Name: cm\_2\_odp\_01\\
Type: integer\\
Default: 30
}

{\footnotesize Визначено частоту налаштувань; перегляду та оновлення базових}


{\footnotesize\bfseries
No: 7\\
Name: cm\_2\_odp\_02\\
Type: string\\
Default: nil
}

{\footnotesize Визначено обставини, що вимагають перегляду та оновлення базових налаштувань}




\subsubsection{КОНФІГУРАЦІЯ СИСТЕМ ТА КОМПОНЕНТІВ ДЛЯ СФЕР З ВИСОКИМ РИЗИКОМ (CM-2(7))}
(a) Видавати [Призначення: визначених організацією систем або компонентів систем] з [Призначенням: визначеними організацією конфігураціями] особам, що перебувають у місцях, які організація вважає місцями зі значним ризиком;\\ 
(b) Застосувати [Призначення: визначені організацією запобіжні заходи безпеки] до компонентів, коли особи повертаються з поїздки.

\vspace{10pt}
{\footnotesize\bfseries
No: 1\\
Name: cm\_2\_7\_a\\
Type: list\\
Default: [\textquotedbl{}admin\textquotedbl{}, \textquotedbl{}security\_officer\textquotedbl{}]
}

{\footnotesize Системи або компоненти системи з конфігураціями видаються особам, що перебувають у місцях, які організація вважає, становлять значний ризик}


{\footnotesize\bfseries
No: 2\\
Name: cm\_2\_7\_b\\
Type: list\\
Default: [\textquotedbl{}admin\textquotedbl{}, \textquotedbl{}security\_officer\textquotedbl{}]
}

{\footnotesize Заходи безпеки застосовуються до систем або компонентів системи, коли особи повертаються з поїздки}


{\footnotesize\bfseries
No: 3\\
Name: cm\_2\_7\_odp\_01\\
Type: list\\
Default: [\textquotedbl{}admin\textquotedbl{}, \textquotedbl{}security\_officer\textquotedbl{}]
}

{\footnotesize Визначено системи або компоненти систем, які мають видаватися особам, що перебувають у місцях зі значним ризиком}




\subsection{УПРАВЛІННЯ ЗМІНАМИ КОНФІГУРАЦІЇ (CM-3)}
a. Визначити типи змін у системі, які контролюються конфігурацією.\\ 
b. Переглядати запропоновані зміни в конфігурації, контрольовані системою, і схвалити або відхилити ці зміни з явним урахуванням аналізу наслідків безпеки.\\ 
c. Документувати рішення зі зміни конфігурації системи.\\ 
d. Впровадити схвалені зміни конфігурації в систему.\\ 
e. Зберігати записи змін конфігурації системі впродовж [Призначення: певного періоду часу, визначеного організацією].\\ 
f. Здійснювати моніторинг і аналіз дій, пов'язаних зі змінами конфігурації системи.\\ 
g. Координувати та впроваджувати нагляд за діяльністю з управління змінами конфігурації за допомогою [Призначення: елементу управління змінами конфігурації, визначеного організацією], який викликається [Вибір (один або кілька): [Призначення: з визначеною організацією частотою]; [Призначення: визначені організацією умови зміни конфігурації]].

\vspace{10pt}
{\footnotesize\bfseries
No: 1\\
Name: cm\_3\_b\_01\\
Type: string\\
Default: nil
}

{\footnotesize Розглядаються запропоновані зміни в конфігурації, що контролюються системою}


{\footnotesize\bfseries
No: 2\\
Name: cm\_3\_b\_02\\
Type: string\\
Default: nil
}

{\footnotesize Запропоновані зміни в конфігурації, що контролюються системою, схвалюються або відхиляються з урахуванням аналізу наслідків безпеки}


{\footnotesize\bfseries
No: 3\\
Name: cm\_3\_c\\
Type: string\\
Default: nil
}

{\footnotesize Рішення про зміну конфігурації системи документуються}


{\footnotesize\bfseries
No: 4\\
Name: cm\_3\_d\\
Type: string\\
Default: nil
}

{\footnotesize Впроваджуються схвалені зміни до конфігурації в систему}


{\footnotesize\bfseries
No: 5\\
Name: cm\_3\_e\\
Type: integer\\
Default: 30
}

{\footnotesize Записи про зміни конфігурації у системі зберігаються протягом <період часу CM-03\_ODP[01]>}


{\footnotesize\bfseries
No: 6\\
Name: cm\_3\_f\_01\\
Type: string\\
Default: nil
}

{\footnotesize Здійснюється моніторинг конфігурації системи}


{\footnotesize\bfseries
No: 7\\
Name: cm\_3\_f\_02\\
Type: string\\
Default: nil
}

{\footnotesize Здійснюється аналіз дій, пов'язаних зі змінами конфігурації системи}


{\footnotesize\bfseries
No: 8\\
Name: cm\_3\_g\_01\\
Type: string\\
Default: nil
}

{\footnotesize Діяльність з управління змінами конфігурації координується та контролюється елемент управління}


{\footnotesize\bfseries
No: 9\\
Name: cm\_3\_odp\_01\\
Type: integer\\
Default: 30
}

{\footnotesize Визначено період часу, протягом якого зберігатимуться записи про зміни конфігурації}


{\footnotesize\bfseries
No: 10\\
Name: cm\_3\_odp\_02\\
Type: string\\
Default: nil
}

{\footnotesize Визначено елементи управління змінами конфігурації, відповідальні за координацію та нагляд за діяльністю з управління змінами}


{\footnotesize\bfseries
No: 11\\
Name: cm\_3\_odp\_03\\
Type: list\\
Default: []
}

{\footnotesize Вибрано одне або декілька з наступних ЗНАЧЕНЬ ПАРАМЕТРІВ: \{частота; коли умови\}}


{\footnotesize\bfseries
No: 12\\
Name: cm\_3\_odp\_04\\
Type: integer\\
Default: 30
}

{\footnotesize Визначено частоту, з якою викликаються елементи управління змінами конфігурації (якщо вибрано)}


{\footnotesize\bfseries
No: 13\\
Name: cm\_3\_odp\_05\\
Type: list\\
Default: []
}

{\footnotesize Визначено умови, за яких викликаються елементи управління змінами конфігурації (якщо вибрано); CM-03(a) визначено та задокументовано типи змін до системи, які контролюються конфігурацією}




\subsection{АНАЛІЗ ВПЛИВУ НА БЕЗПЕКУ ТА ПРИВАТНІСТЬ (CM-4)}
Аналізувати зміни в системі, щоб визначити потенційну загрозу безпеці та приватності перед реалізацією змін.

\vspace{10pt}
{\footnotesize\bfseries
No: 1\\
Name: cm\_4\_01\\
Type: string\\
Default: nil
}

{\footnotesize Аналізуються зміни в системі, щоб визначити потенційну загрозу безпеці перед реалізацією змін}


{\footnotesize\bfseries
No: 2\\
Name: cm\_4\_02\\
Type: string\\
Default: nil
}

{\footnotesize Аналізуються зміни в системі, щоб визначити потенційну загрозу конфіденційності перед реалізацією змін}




\subsubsection{ВЕРИФІКАЦІЯ ФУНКЦІЙ БЕЗПЕКИ ТА ПРИВАТНОСТІ (CM-4(2))}
Після змін у системі переконайтеся, що відповідні заходи захисту реалізовано правильно і вони функціонують належним чином та дають бажаний результат щодо дотримання вимог безпеки та приватності для системи.

\vspace{10pt}
{\footnotesize\bfseries
No: 1\\
Name: cm\_4\_2\_01\\
Type: string\\
Default: nil
}

{\footnotesize Заходи захисту, на які було здійснено вплив, реалізовані правильно з точки зору відповідності вимогам безпеки системи після внесення змін до системи}


{\footnotesize\bfseries
No: 2\\
Name: cm\_4\_2\_02\\
Type: string\\
Default: nil
}

{\footnotesize Заходи захисту, на які було здійснено вплив, реалізовані правильно з точки зору відповідності вимогам конфіденційності системи після внесення змін до системи}


{\footnotesize\bfseries
No: 3\\
Name: cm\_4\_2\_03\\
Type: string\\
Default: nil
}

{\footnotesize Заходи захисту, на які було здійснено вплив, функціонують належним чином з точки зору відповідності вимогам безпеки системи після внесення змін до системи}


{\footnotesize\bfseries
No: 4\\
Name: cm\_4\_2\_04\\
Type: string\\
Default: nil
}

{\footnotesize Заходи захисту, на які було здійснено вплив, функціонують належним чином з точки зору відповідності вимогам конфіденційності системи після внесення змін до системи}


{\footnotesize\bfseries
No: 5\\
Name: cm\_4\_2\_05\\
Type: string\\
Default: nil
}

{\footnotesize Заходи захисту, на які було здійснено вплив, дають бажаний результат з точки зору відповідності вимогам безпеки системи після внесення змін до системи}


{\footnotesize\bfseries
No: 6\\
Name: cm\_4\_2\_06\\
Type: string\\
Default: nil
}

{\footnotesize Заходи захисту, на які було здійснено вплив, дають бажаний результат з точки зору відповідності вимогам конфіденційності системи після внесення змін до системи; ПОТЕНЦІЙНІ МЕТОДИ ТА ОБ'ЄКТИ ОЦІНКИ:}




\subsection{ОБМЕЖЕННЯ ДОСТУПУ ДО ЗМІНИ (CM-5)}
Визначити, задокументувати, затвердити та забезпечити застосування фізичних і логічних обмежень доступу, пов'язаних зі змінами в системі.

\vspace{10pt}
{\footnotesize\bfseries
No: 1\\
Name: cm\_5\_01\\
Type: string\\
Default: nil
}

{\footnotesize Визначені та задокументовані фізичні обмеження доступу, пов'язані зі змінами в системі}


{\footnotesize\bfseries
No: 2\\
Name: cm\_5\_02\\
Type: string\\
Default: nil
}

{\footnotesize Затверджені фізичні обмеження доступу, пов'язані зі змінами в системі}


{\footnotesize\bfseries
No: 3\\
Name: cm\_5\_03\\
Type: string\\
Default: nil
}

{\footnotesize Застосовуються фізичні обмеження доступу, пов'язані зі змінами в системі}


{\footnotesize\bfseries
No: 4\\
Name: cm\_5\_04\\
Type: string\\
Default: nil
}

{\footnotesize Визначені та задокументовані логічні обмеження доступу, пов'язані зі змінами в системі}


{\footnotesize\bfseries
No: 5\\
Name: cm\_5\_05\\
Type: string\\
Default: nil
}

{\footnotesize Затверджені логічні обмеження доступу, пов'язані зі змінами в системі}


{\footnotesize\bfseries
No: 6\\
Name: cm\_5\_06\\
Type: string\\
Default: nil
}

{\footnotesize Застосовуються логічні обмеження доступу, пов'язані зі змінами в системі}




\subsection{НАЛАШТУВАННЯ КОНФІГУРАЦІЇ (CM-6)}
a. Встановити та задокументувати параметри конфігурації компонентів, які застосовуються в системі, які відображають найбільш обмежений режим, що відповідає експлуатаційним вимогам, використовуючи [Призначення: визначені організацією загальні безпечні конфігурації].\\ 
b. Реалізувати конфігураційні установки.\\ 
c. Визначити, задокументувати та затвердити будь-які відхилення від встановлених конфігураційних параметрів конфігурації для [Призначення: визначених організацією компонентів системи] на основі [Призначення: визначених організацією експлуатаційних вимог].\\ 
d. Відстежувати та керувати змінами конфігураційних параметрів відповідно до організаційної політики та процедур.

\vspace{10pt}
{\footnotesize\bfseries
No: 1\\
Name: cm\_6\_a\\
Type: string\\
Default: nil
}

{\footnotesize Налаштування конфігурації, які відображають найбільш обмежувальний режим, що відповідає експлуатаційним вимогам, встановлені та задокументовані для компонентів, що застосовуються в системі з використанням безпечні конфігурації}


{\footnotesize\bfseries
No: 2\\
Name: cm\_6\_b\\
Type: string\\
Default: nil
}

{\footnotesize Реалізовано установки конфігурації, задокументовані в CM-06a}


{\footnotesize\bfseries
No: 3\\
Name: cm\_6\_c\_01\\
Type: string\\
Default: nil
}

{\footnotesize Будь-які відхилення від встановлених параметрів конфігурації для компонентів системи визначаються та документуються на основі експлуатаційних вимог}


{\footnotesize\bfseries
No: 4\\
Name: cm\_6\_c\_02\\
Type: string\\
Default: nil
}

{\footnotesize Будь-які відхилення від встановлених налаштувань конфігурації для компонентів системи затверджуються}


{\footnotesize\bfseries
No: 5\\
Name: cm\_6\_d\_01\\
Type: list\\
Default: [\textquotedbl{}default\_deny\_rule\textquotedbl{}, \textquotedbl{}abac\_rule\_1\textquotedbl{}]
}

{\footnotesize Зміни в налаштуваннях конфігурації відстежуються відповідно до політики та процедур організації}


{\footnotesize\bfseries
No: 6\\
Name: cm\_6\_d\_02\\
Type: list\\
Default: [\textquotedbl{}default\_deny\_rule\textquotedbl{}, \textquotedbl{}abac\_rule\_1\textquotedbl{}]
}

{\footnotesize Зміни налаштувань конфігурації керуються відповідно до політики та процедур організації}


{\footnotesize\bfseries
No: 7\\
Name: cm\_6\_odp\_01\\
Type: string\\
Default: nil
}

{\footnotesize Визначено загальні безпечні конфігурації для встановлення та документування параметрів конфігурації компонентів, які застосовуються в системі; CM-06\_ODP[02] визначено компоненти системи, для яких необхідно затвердити відхилення}


{\footnotesize\bfseries
No: 8\\
Name: cm\_6\_odp\_03\\
Type: string\\
Default: nil
}

{\footnotesize Визначені експлуатаційні вимоги, що вимагають затвердження відхилень}




\subsection{МІНІМАЛЬНО НЕОБХІДНА ФУНКЦІОНАЛЬНІСТЬ (CM-7)}
a. Налаштуйте систему для забезпечення лише [Призначення: основні функції, визначені організацією для місії];\\ 
b. Заборонити або обмежити використання таких функцій, портів, протоколів, програмного забезпечення та/або служб: [Призначення: визначені організацією заборонені або обмежені функції, системні порти, протоколи, програмне забезпечення та/або служби].

\vspace{10pt}
{\footnotesize\bfseries
No: 1\\
Name: cm\_7\_a\\
Type: string\\
Default: nil
}

{\footnotesize Система налаштована на забезпечення лише основні функції системи }


{\footnotesize\bfseries
No: 2\\
Name: cm\_7\_b\_01\\
Type: string\\
Default: nil
}

{\footnotesize Використання функцій заборонено або обмежено}


{\footnotesize\bfseries
No: 3\\
Name: cm\_7\_b\_02\\
Type: string\\
Default: nil
}

{\footnotesize Використання порти заборонено або обмежено}


{\footnotesize\bfseries
No: 4\\
Name: cm\_7\_b\_03\\
Type: string\\
Default: \textquotedbl{}TLS 1.3\textquotedbl{}
}

{\footnotesize Використання протоколи заборонено або обмежено}


{\footnotesize\bfseries
No: 5\\
Name: cm\_7\_b\_04\\
Type: string\\
Default: nil
}

{\footnotesize Використання програмне забезпечення заборонено або обмежено}


{\footnotesize\bfseries
No: 6\\
Name: cm\_7\_b\_05\\
Type: string\\
Default: nil
}

{\footnotesize Використання служби заборонено або обмежено}


{\footnotesize\bfseries
No: 7\\
Name: cm\_7\_odp\_01\\
Type: string\\
Default: nil
}

{\footnotesize Визначено основні функції системи, необхідні для виконання місії}


{\footnotesize\bfseries
No: 8\\
Name: cm\_7\_odp\_02\\
Type: string\\
Default: nil
}

{\footnotesize Визначено функції, які необхідно заборонити або обмежити}


{\footnotesize\bfseries
No: 9\\
Name: cm\_7\_odp\_03\\
Type: string\\
Default: nil
}

{\footnotesize Визначено системні порти, які необхідно заборонити або обмежити}


{\footnotesize\bfseries
No: 10\\
Name: cm\_7\_odp\_04\\
Type: string\\
Default: \textquotedbl{}TLS 1.3\textquotedbl{}
}

{\footnotesize Визначено протоколи, які необхідно заборонити або обмежити; CM-07\_ODP[05] визначено програмне забезпечення, яке необхідно заборонити або обмежити}


{\footnotesize\bfseries
No: 11\\
Name: cm\_7\_odp\_06\\
Type: string\\
Default: nil
}

{\footnotesize Визначено служби, які необхідно заборонити або обмежити}




\subsubsection{ПЕРІОДИЧНИЙ ПЕРЕГЛЯД (CM-7(1))}


\vspace{10pt}
{\footnotesize\bfseries
No: 1\\
Name: cm\_7\_1\_a\\
Type: string\\
Default: \textquotedbl{}TLS 1.3\textquotedbl{}
}

{\footnotesize Система переглядається частота для виявлення непотрібних та/або незахищених функцій, портів, протоколів і послуг}


{\footnotesize\bfseries
No: 2\\
Name: cm\_7\_1\_b\_01\\
Type: string\\
Default: nil
}

{\footnotesize Функції, які вважаються непотрібними та/або незахищеними, вимкнено}


{\footnotesize\bfseries
No: 3\\
Name: cm\_7\_1\_b\_02\\
Type: string\\
Default: nil
}

{\footnotesize Порти, які вважаються непотрібними та/або незахищеними, вимкнено}


{\footnotesize\bfseries
No: 4\\
Name: cm\_7\_1\_b\_03\\
Type: string\\
Default: \textquotedbl{}TLS 1.3\textquotedbl{}
}

{\footnotesize Протоколи, які вважаються непотрібними та/або незахищеними, вимкнено}


{\footnotesize\bfseries
No: 5\\
Name: cm\_7\_1\_b\_04\\
Type: string\\
Default: nil
}

{\footnotesize Послуги, які вважаються непотрібними та/або незахищеними, вимкнено}


{\footnotesize\bfseries
No: 6\\
Name: cm\_7\_1\_odp\_01\\
Type: string\\
Default: \textquotedbl{}TLS 1.3\textquotedbl{}
}

{\footnotesize Визначено частоту, з якою слід переглядати систему для виявлення непотрібних та/або незахищених функцій, портів, протоколів і послуг}


{\footnotesize\bfseries
No: 7\\
Name: cm\_7\_1\_odp\_02\\
Type: string\\
Default: nil
}

{\footnotesize Визначено функції, які слід вимкнути, якщо вони вважаються непотрібними або незахищеними; CM-07(01)\_ODP[03] визначено порти, які слід вимкнути, якщо вони вважаються непотрібними або незахищеними}


{\footnotesize\bfseries
No: 8\\
Name: cm\_7\_1\_odp\_04\\
Type: string\\
Default: \textquotedbl{}TLS 1.3\textquotedbl{}
}

{\footnotesize Визначено протоколи, які слід вимкнути, якщо вони вважаються непотрібними або незахищеними}


{\footnotesize\bfseries
No: 9\\
Name: cm\_7\_1\_odp\_05\\
Type: string\\
Default: nil
}

{\footnotesize Визначено послуги, які слід вимкнути, якщо вони вважаються непотрібними або незахищеними}




\subsubsection{АВТОРИЗОВАНЕ ПРОГРАМНЕ ЗАБЕЗПЕЧЕННЯ -- БІЛИЙ СПИСОК (CM-7(5))}
(a) Визначити [Призначення: визначені організацією програмне забезпечення, яке авторизовано виконується в системі]\\ 
(b) Впровадити політику «заборони всього, за винятком деяких», щоб дозволити виконання авторизованих програм у системі.\\ 
(c) Переглядати та оновлювати список авторизованих програм [Призначення: з визначеною організацією частотою].

\vspace{10pt}
{\footnotesize\bfseries
No: 1\\
Name: cm\_7\_5\_a\\
Type: string\\
Default: nil
}

{\footnotesize Визначено програмне забезпечення}


{\footnotesize\bfseries
No: 2\\
Name: cm\_7\_5\_b\\
Type: list\\
Default: [\textquotedbl{}default\_deny\_rule\textquotedbl{}, \textquotedbl{}abac\_rule\_1\textquotedbl{}]
}

{\footnotesize Політика \textquotedbl{}заборонити все, дозволити за винятком\textquotedbl{} застосовуэться, щоб дозволити виконання авторизованих програм у системі}


{\footnotesize\bfseries
No: 3\\
Name: cm\_7\_5\_c\\
Type: list\\
Default: []
}

{\footnotesize Переглядається та оновлюється список авторизованих програм частота}


{\footnotesize\bfseries
No: 4\\
Name: cm\_7\_5\_odp\_01\\
Type: string\\
Default: nil
}

{\footnotesize Визначено програмне забезпечення, яке авторизовано для виконання в системі}


{\footnotesize\bfseries
No: 5\\
Name: cm\_7\_5\_odp\_02\\
Type: integer\\
Default: 30
}

{\footnotesize Визначено частоту перегляду та оновлення списку авторизованих програм}




\subsection{ІНВЕНТАРИЗАЦІЯ КОМПОНЕНТІВ СИСТЕМИ (CM-8)}
a.\\ 
b. Розробити та задокументувати процес інвентаризації компонентів системи, який:\\ 
1. точно описує поточну систему;\\ 
2. охоплює всі компоненти в межах акредитації системи;\\ 
3. не включає повторний облік компонентів або компонентів, будь-якої іншої системи;\\ 
4. визначає рівень деталізації, який є необхідним для відстеження та звітування;\\ 
5. включає інформацію для досягнення підзвітності компонентів системи: [Призначення: визначена організацією інформація, необхідна для досягнення ефективної підзвітності компонентів системи]. Переглядати та оновлювати опис компонентів системи з [Призначення: визначеною організацією частотою].

\vspace{10pt}
{\footnotesize\bfseries
No: 1\\
Name: cm\_8\_a\_01\\
Type: string\\
Default: nil
}

{\footnotesize Розроблено та задокументовано процес інвентаризації компонентів системи, який точно описує поточну систему}


{\footnotesize\bfseries
No: 2\\
Name: cm\_8\_a\_02\\
Type: string\\
Default: nil
}

{\footnotesize Розроблено та задокументовано процес інвентаризації компонентів системи, який охоплює всі компоненти в межах акредитації системи}


{\footnotesize\bfseries
No: 3\\
Name: cm\_8\_a\_03\\
Type: string\\
Default: nil
}

{\footnotesize Розроблено та задокументовано процес інвентаризації компонентів системи, який не включає повторний облік компонентів або компонентів, будь-якої іншої системи}


{\footnotesize\bfseries
No: 4\\
Name: cm\_8\_a\_04\\
Type: string\\
Default: nil
}

{\footnotesize Розроблено та задокументовано процес інвентаризації компонентів системи, який визначає рівень деталізації, який є необхідним для відстеження та звітування}


{\footnotesize\bfseries
No: 5\\
Name: cm\_8\_a\_05\\
Type: string\\
Default: nil
}

{\footnotesize Розроблено та задокументовано процес інвентаризації компонентів системи, який включає інформацію}


{\footnotesize\bfseries
No: 6\\
Name: cm\_8\_b\\
Type: string\\
Default: \textquotedbl{}щорічно\textquotedbl{}
}

{\footnotesize Переглядається та оновлюється опис компонентів системи частота}


{\footnotesize\bfseries
No: 7\\
Name: cm\_8\_odp\_01\\
Type: string\\
Default: nil
}

{\footnotesize Визначено інформацію, яка вважається необхідною для досягнення ефективної підзвітності компонентів системи}


{\footnotesize\bfseries
No: 8\\
Name: cm\_8\_odp\_02\\
Type: integer\\
Default: 30
}

{\footnotesize Визначено частоту перегляду та оновлення опису компонентів системи}




\subsubsection{ОНОВЛЕННЯ ПІД ЧАС ВСТАНОВЛЕННЯ ТА ВИДАЛЕННЯ (CM-8(1))}


\vspace{10pt}
{\footnotesize\bfseries
No: 1\\
Name: cm\_8\_1\_01\\
Type: string\\
Default: nil
}

{\footnotesize Інвентаризація компонентів системи оновлюється в рамках інсталяцій компонентів системи}


{\footnotesize\bfseries
No: 2\\
Name: cm\_8\_1\_02\\
Type: string\\
Default: nil
}

{\footnotesize Інвентаризація компонентів системи оновлюється в рамках видалення компонентів системи}


{\footnotesize\bfseries
No: 3\\
Name: cm\_8\_1\_03\\
Type: string\\
Default: nil
}

{\footnotesize Інвентаризація компонентів системи оновлюється в рамках оновлення компонентів системи}




\subsection{РОЗТАШУВАННЯ ІНФОРМАЦІЇ (CM-12)}


\vspace{10pt}
{\footnotesize\bfseries
No: 1\\
Name: cm\_12\_a\_01\\
Type: string\\
Default: nil
}

{\footnotesize Місцезнаходження інформація визначено та задокументовано}


{\footnotesize\bfseries
No: 2\\
Name: cm\_12\_a\_02\\
Type: string\\
Default: nil
}

{\footnotesize Визначено та задокументовано конкретні компоненти системи, на яких обробляється інформація}


{\footnotesize\bfseries
No: 3\\
Name: cm\_12\_a\_03\\
Type: string\\
Default: nil
}

{\footnotesize Визначено та задокументовано конкретні компоненти системи, на яких зберігається інформація}


{\footnotesize\bfseries
No: 4\\
Name: cm\_12\_b\_01\\
Type: string\\
Default: nil
}

{\footnotesize Визначено та задокументовано користувачів, які мають доступ до системи та компонентів системи, де обробляється інформація}


{\footnotesize\bfseries
No: 5\\
Name: cm\_12\_b\_02\\
Type: string\\
Default: nil
}

{\footnotesize Визначено та задокументовано користувачів, які мають доступ до системи та компонентів системи, де зберігається інформація; CM-12(c)[01] задокументовано зміни розташування ( наприклад, системи або компонентів системи), де обробляється інформація}


{\footnotesize\bfseries
No: 6\\
Name: cm\_12\_c\_02\\
Type: string\\
Default: nil
}

{\footnotesize Задокументовано зміни розташування ( наприклад, системи або компонентів системи), де зберігається інформація}


{\footnotesize\bfseries
No: 7\\
Name: cm\_12\_odp\\
Type: string\\
Default: nil
}

{\footnotesize Визначено інформацію, місцезнаходження якої має бути визначено та задокументовано}




\section{CP}
Клас заходів захисту CP --- ПЛАНУВАННЯ БЕЗПЕРЕРВНОЇ РОБОТИ

Цей клас описує заходи для відновлення функціонування інформаційної системи та забезпечення її безперебійної роботи у разі надзвичайних ситуацій.

\textbf{Перелік заходів захисту:} Резервне копіювання (CP-9); Криптографічний захист (CP-9(8)).

\subsection{РЕЗЕРВНЕ КОПІЮВАННЯ (CP-9)}


\vspace{10pt}
{\footnotesize\bfseries
No: 1\\
Name: cp\_9\_01\\
Type: string\\
Default: nil
}

{\footnotesize РЕЗЕРВНЕ КОПІЮВАННЯ МЕТА ОЦІНКИ: Визначити, чи:}


{\footnotesize\bfseries
No: 2\\
Name: cp\_9\_a\\
Type: string\\
Default: \textquotedbl{}щорічно\textquotedbl{}
}

{\footnotesize Резервне копіювання інформації користувача, що міститься в компонентах системи, здійснюється частота}


{\footnotesize\bfseries
No: 3\\
Name: cp\_9\_b\\
Type: string\\
Default: \textquotedbl{}щорічно\textquotedbl{}
}

{\footnotesize Виконується резервне копіювання інформації системи, що міститься в системі частота; CP-09(с) створюються резервні копії документації системи, включаючи документацію, пов'язану з безпекою та конфіденційністю частота}


{\footnotesize\bfseries
No: 4\\
Name: cp\_9\_d\_01\\
Type: string\\
Default: nil
}

{\footnotesize Конфіденційність резервних копій інформації захищена}


{\footnotesize\bfseries
No: 5\\
Name: cp\_9\_d\_02\\
Type: string\\
Default: nil
}

{\footnotesize Цілісність резервних копій інформації захищена}


{\footnotesize\bfseries
No: 6\\
Name: cp\_9\_d\_03\\
Type: string\\
Default: nil
}

{\footnotesize Доступність резервних копій інформації захищена}


{\footnotesize\bfseries
No: 7\\
Name: cp\_9\_odp\_01\\
Type: string\\
Default: nil
}

{\footnotesize Визначено компоненти системи, для яких необхідно проводити резервне копіювання інформації користувачів}


{\footnotesize\bfseries
No: 8\\
Name: cp\_9\_odp\_02\\
Type: integer\\
Default: 30
}

{\footnotesize Визначено частоту, з якою слід проводити резервне копіювання інформації користувача відповідно до часу відновлення та цілей відновлення}


{\footnotesize\bfseries
No: 9\\
Name: cp\_9\_odp\_03\\
Type: integer\\
Default: 30
}

{\footnotesize Визначено частоту проведення резервного копіювання інформації системи, що відповідає завдань відновлення і встановлених цілей відновлення}


{\footnotesize\bfseries
No: 10\\
Name: cp\_9\_odp\_04\\
Type: integer\\
Default: 30
}

{\footnotesize Визначено частоту, з якою слід проводити резервне копіювання документації системи відповідно до часу відновлення та цілей точки відновлення}




\subsubsection{КРИПТОГРАФІЧНИЙ ЗАХИСТ (CP-9(8))}


\vspace{10pt}
{\footnotesize\bfseries
No: 1\\
Name: cp\_9\_8\_01\\
Type: string\\
Default: \textquotedbl{}AES-256-GCM\textquotedbl{}
}

{\footnotesize Реалізовано криптографічні механізми для запобігання несанкціонованому розкриттю та зміні резервної інформації}


{\footnotesize\bfseries
No: 2\\
Name: cp\_9\_8\_odp\\
Type: string\\
Default: nil
}

{\footnotesize Визначено резервні копії інформації для захисту від несанкціонованого розкриття та змін}




\section{IA}
Клас заходів захисту IA --- ІДЕНТИФІКАЦІЯ ТА АВТЕНТИФІКАЦІЯ

Цей клас відповідає за однозначне розпізнавання користувачів або пристроїв та підтвердження їхньої справжності перед наданням доступу до системи.

\textbf{Перелік заходів захисту:} Політика та процедури ідентифікації та автентифікації (IA-1); Ідентифікація та автентифікація користувачів (IA-2); Багатофакторна автентифікація привілейованих облікових записів (IA-2(1)); Багатофакторна автентифікація непривілейованих (IA-2(2)); Доступ до облікових записів --- стійкість до відтворення (IA-2(8)); Ідентифікація та автентифікація пристроїв (IA-3); Управління ідентифікацією (IA-4); Ідентифікація статусу користувача (IA-4(4)); АВТЕНТИФІКАТОР УПРАВЛІННЯ; Автентифікація на основі пароля (IA-5(1)); ПРИХОВУВАННЯ ЗВОРОТНОГО ЗВ'ЯЗКУ АВТЕНТИФІКАТОРА; Повторна автентифікація (IA-11).

\subsection{Політика та процедури ідентифікації та автентифікації (IA-1)}


\vspace{10pt}
{\footnotesize\bfseries
No: 1\\
Name: ia\_1\_a\_01\\
Type: list\\
Default: [\textquotedbl{}default\_deny\_rule\textquotedbl{}, \textquotedbl{}abac\_rule\_1\textquotedbl{}]
}

{\footnotesize Розроблено та задокументовано політику ідентифікації та автентифікації}


{\footnotesize\bfseries
No: 2\\
Name: ia\_1\_a\_02\\
Type: list\\
Default: [\textquotedbl{}admin\textquotedbl{}, \textquotedbl{}security\_officer\textquotedbl{}]
}

{\footnotesize Політика ідентифікації та автентифікації поширюється на персонал або ролі}


{\footnotesize\bfseries
No: 3\\
Name: ia\_1\_a\_03\\
Type: list\\
Default: [\textquotedbl{}default\_deny\_rule\textquotedbl{}, \textquotedbl{}abac\_rule\_1\textquotedbl{}]
}

{\footnotesize Розроблені та задокументовані процедури ідентифікації та автентифікації, що сприяють впровадженню політики ідентифікації та автентифікації, а також відповідні заходи ідентифікації та перевірки автентичності}


{\footnotesize\bfseries
No: 4\\
Name: ia\_1\_b\\
Type: list\\
Default: [\textquotedbl{}admin\textquotedbl{}, \textquotedbl{}security\_officer\textquotedbl{}]
}

{\footnotesize Посадова особа призначається для управління політикою та процедурами ідентифікації та автентифікації; IA-01(c)[01][01] переглядається та оновлюється поточна політика ідентифікації та автентифікації частота; IA-01(c)[01][02] переглядається та оновлюється поточна політика ідентифікації та автентифікації після подій; IA-01(c)[02][01] переглядаються та оновлюються поточні процедури ідентифікації та автентифікації частота; IA-01(c)[02][02] переглядаються та оновлюються поточні процедури ідентифікації та автентифікації після подій}


{\footnotesize\bfseries
No: 5\\
Name: ia\_1\_odp\_01\\
Type: list\\
Default: [\textquotedbl{}admin\textquotedbl{}, \textquotedbl{}security\_officer\textquotedbl{}]
}

{\footnotesize Визначено персонал або ролі, на які поширюється політика ідентифікації та автентифікації}


{\footnotesize\bfseries
No: 6\\
Name: ia\_1\_odp\_02\\
Type: list\\
Default: [\textquotedbl{}admin\textquotedbl{}, \textquotedbl{}security\_officer\textquotedbl{}]
}

{\footnotesize Визначено персонал або ролі, на які поширюються процедури ідентифікації та автентифікації}


{\footnotesize\bfseries
No: 7\\
Name: ia\_1\_odp\_03\\
Type: string\\
Default: nil
}

{\footnotesize Вибрано одне або декілька з наступних ЗНАЧЕНЬ ПАРАМЕТРІВ: \{рівень організації; рівень місії/бізнес-процесів; рівень системи\}}


{\footnotesize\bfseries
No: 8\\
Name: ia\_1\_odp\_04\\
Type: list\\
Default: [\textquotedbl{}admin\textquotedbl{}, \textquotedbl{}security\_officer\textquotedbl{}]
}

{\footnotesize Визначено посадову особу, яка управляє політикою та процедурами ідентифікації та автентифікації}


{\footnotesize\bfseries
No: 9\\
Name: ia\_1\_odp\_05\\
Type: list\\
Default: [\textquotedbl{}default\_deny\_rule\textquotedbl{}, \textquotedbl{}abac\_rule\_1\textquotedbl{}]
}

{\footnotesize Визначено частоту, з якою переглядається та оновлюється поточна політика ідентифікації та автентифікації}


{\footnotesize\bfseries
No: 10\\
Name: ia\_1\_odp\_06\\
Type: list\\
Default: [\textquotedbl{}default\_deny\_rule\textquotedbl{}, \textquotedbl{}abac\_rule\_1\textquotedbl{}]
}

{\footnotesize Визначено події, які потребують перегляду та оновлення поточної політики ідентифікації та автентифікації}


{\footnotesize\bfseries
No: 11\\
Name: ia\_1\_odp\_07\\
Type: integer\\
Default: 30
}

{\footnotesize Визначено частоту, з якою переглядаються та оновлюються поточні процедури ідентифікації та автентифікації}


{\footnotesize\bfseries
No: 12\\
Name: ia\_1\_odp\_08\\
Type: list\\
Default: [\textquotedbl{}login\textquotedbl{}, \textquotedbl{}logout\textquotedbl{}, \textquotedbl{}failed\_attempt\textquotedbl{}]
}

{\footnotesize Визначено події, які потребують перегляду та оновлення процедур ідентифікації та автентифікації}




\subsection{ІДЕНТИФІКАЦІЯ ТА АВТЕНТИФІКАЦІЯ КОРИСТУВАЧІВ (IA-2)}


\vspace{10pt}
{\footnotesize\bfseries
No: 1\\
Name: ia\_2\_01\\
Type: string\\
Default: nil
}

{\footnotesize ІДЕНТИФІКАЦІЯ ТА АВТЕНТИФІКАЦІЯ (КОРИСТУВАЧІВ ОРГАНІЗАЦІЇ) МЕТА ОЦІНКИ: Визначити, чи:}


{\footnotesize\bfseries
No: 2\\
Name: ia\_2\_02\\
Type: string\\
Default: nil
}

{\footnotesize Процеси що діють від імені користувачів унікально ідентифіковані та автентифіковані}




\subsubsection{БАГАТОФАКТОРНА АВТЕНТИФІКАЦІЯ ПРИВІЛЕЙОВАНИХ ОБЛІКОВИХ ЗАПИСІВ (IA-2(1))}


\vspace{10pt}
{\footnotesize\bfseries
No: 1\\
Name: ia\_2\_1\_01\\
Type: string\\
Default: nil
}

{\footnotesize Реалізувано багатофакторну автентифікацію для доступу до привілейованих облікових записів}




\subsubsection{БАГАТОФАКТОРНА АВТЕНТИФІКАЦІЯ НЕПРИВІЛЕЙОВАНИХ (IA-2(2))}


\vspace{10pt}
{\footnotesize\bfseries
No: 1\\
Name: ia\_2\_2\_01\\
Type: string\\
Default: nil
}

{\footnotesize Реалізовано багатофакторну автентифікацію для доступу до непривілейованих облікових записів}




\subsubsection{Доступ до облікових записів --- стійкість до відтворення (IA-2(8))}


\vspace{10pt}
{\footnotesize\bfseries
No: 1\\
Name: ia\_2\_8\_odp\\
Type: string\\
Default: nil
}

{\footnotesize Вибрано одне або декілька з наступних ЗНАЧЕНЬ ПАРАМЕТРІВ: \{привілейовані облікові записи; непривілейовані облікові записи\}}




\subsection{ІДЕНТИФІКАЦІЯ ТА АВТЕНТИФІКАЦІЯ ПРИСТРОЇВ (IA-3)}


\vspace{10pt}
{\footnotesize\bfseries
No: 1\\
Name: ia\_3\_odp\_01\\
Type: string\\
Default: nil
}

{\footnotesize Визначені пристрої та/або типи пристроїв, які повинні бути унікально ідентифіковані та автентифіковані перед установкою підключення}


{\footnotesize\bfseries
No: 2\\
Name: ia\_3\_odp\_02\\
Type: string\\
Default: nil
}

{\footnotesize Вибрано одне або декілька з наступних ЗНАЧЕНЬ ПАРАМЕТРІВ: \{локальний; віддалений; мережевий\}}




\subsection{УПРАВЛІННЯ ІДЕНТИФІКАЦІЄЮ (IA-4)}


\vspace{10pt}
{\footnotesize\bfseries
No: 1\\
Name: ia\_4\_a\\
Type: list\\
Default: [\textquotedbl{}admin\textquotedbl{}, \textquotedbl{}security\_officer\textquotedbl{}]
}

{\footnotesize Управління ідентифікаторами здійснюється шляхом отримання дозволу від персоналу або ролей на призначення ідентифікатора особі, групі, ролі або пристрою}


{\footnotesize\bfseries
No: 2\\
Name: ia\_4\_b\\
Type: list\\
Default: [\textquotedbl{}admin\textquotedbl{}, \textquotedbl{}security\_officer\textquotedbl{}]
}

{\footnotesize Управління ідентифікаторами здійснюється шляхом вибору ідентифікатора, який ідентифікує окрему особу, групу, ролі або пристрій}


{\footnotesize\bfseries
No: 3\\
Name: ia\_4\_c\\
Type: list\\
Default: [\textquotedbl{}admin\textquotedbl{}, \textquotedbl{}security\_officer\textquotedbl{}]
}

{\footnotesize Управління ідентифікаторами здійснюється шляхом призначення ідентифікатора особі, групі, ролі або пристрою; IA-04(d) ідентифікатори управляються шляхом запобігання повторному використанню ідентифікаторів впродовж період часу}


{\footnotesize\bfseries
No: 4\\
Name: ia\_4\_odp\_01\\
Type: list\\
Default: [\textquotedbl{}admin\textquotedbl{}, \textquotedbl{}security\_officer\textquotedbl{}]
}

{\footnotesize Визначено персонал або ролі, від яких необхідно отримати дозвіл на призначення ідентифікатора}


{\footnotesize\bfseries
No: 5\\
Name: ia\_4\_odp\_02\\
Type: integer\\
Default: 30
}

{\footnotesize Визначено період часу для запобігання повторному використанню ідентифікаторів}




\subsubsection{ІДЕНТИФІКАЦІЯ СТАТУСУ КОРИСТУВАЧА (IA-4(4))}


\vspace{10pt}
{\footnotesize\bfseries
No: 1\\
Name: ia\_4\_4\_01\\
Type: string\\
Default: nil
}

{\footnotesize Управління індивідуальними ідентифікаторами, однозначно ідентифікуючи кожного індивідуума ознака, що ідентифікує індивідуальний статус}


{\footnotesize\bfseries
No: 2\\
Name: ia\_4\_4\_odp\\
Type: string\\
Default: nil
}

{\footnotesize Визначено ознаку, що ідентифікує індивідуальний статус}




\subsection{АВТЕНТИФІКАТОР УПРАВЛІННЯ}


\vspace{10pt}
{\footnotesize\bfseries
No: 1\\
Name: ia\_5\_a\\
Type: list\\
Default: [\textquotedbl{}admin\textquotedbl{}, \textquotedbl{}security\_officer\textquotedbl{}]
}

{\footnotesize Управління системними автентифікаторами здійснюється шляхом перевірки, як частини початкового розподілу автентифікатора, особи, групи, ролі або пристрою, який отримує автентифікатор}


{\footnotesize\bfseries
No: 2\\
Name: ia\_5\_b\\
Type: string\\
Default: nil
}

{\footnotesize Управління системними автентифікаторами здійснюється шляхом створення вихідного вмісту автентифікатора для будь-яких автентифікаторів, виданих організацією}


{\footnotesize\bfseries
No: 3\\
Name: ia\_5\_c\\
Type: string\\
Default: nil
}

{\footnotesize Управління системними автентифікаторами здійснюється шляхом забезпечення того, щоб автентифікатори мали достатню стійкість механізму для їх використання за призначенням; IA-05(d) управління системними автентифікаторами здійснюється шляхом створення та реалізація адміністративних процедур для первинного розповсюдження автентифікаторів, для втрачених/скомпрометованих або пошкоджених автентифікаторів, а також для відкликання автентифікаторів}


{\footnotesize\bfseries
No: 4\\
Name: ia\_5\_e\\
Type: string\\
Default: nil
}

{\footnotesize Управління системними автентифікаторами здійснюється шляхом зміни типових автентифікаторів перед першим використанням}


{\footnotesize\bfseries
No: 5\\
Name: ia\_5\_f\\
Type: list\\
Default: [\textquotedbl{}login\textquotedbl{}, \textquotedbl{}logout\textquotedbl{}, \textquotedbl{}failed\_attempt\textquotedbl{}]
}

{\footnotesize Управління системними автентифікаторами здійснюється шляхом зміни/оновлення автентифікаторів у встановлений період часу або коли відбуваються події}


{\footnotesize\bfseries
No: 6\\
Name: ia\_5\_g\\
Type: string\\
Default: nil
}

{\footnotesize Управління системними автентифікаторами здійснюється шляхом захисту вмісту автентифікатора від несанкціонованого розкриття та модифікацій}


{\footnotesize\bfseries
No: 7\\
Name: ia\_5\_h\\
Type: string\\
Default: nil
}

{\footnotesize Управління системними автентифікаторами здійснюється шляхом вимоги до осіб, які використовують пристрої, використовувати спеціальні заходи безпеки для захисту автентифікаторів}


{\footnotesize\bfseries
No: 8\\
Name: ia\_5\_i\\
Type: string\\
Default: nil
}

{\footnotesize Управління системними автентифікаторами здійснюється шляхом вимоги змінювати автентифікатори для облікових записів груп/ролей при зміні членства в цих облікових записах}


{\footnotesize\bfseries
No: 9\\
Name: ia\_5\_odp\_01\\
Type: integer\\
Default: 30
}

{\footnotesize Визначено період часу для зміни або оновлення автентифікаторів за типом автентифікатора}




\subsubsection{АВТЕНТИФІКАЦІЯ НА ОСНОВІ ПАРОЛЯ (IA-5(1))}


\vspace{10pt}
{\footnotesize\bfseries
No: 1\\
Name: ia\_5\_1\_a\\
Type: list\\
Default: [\textquotedbl{}admin\textquotedbl{}, \textquotedbl{}security\_officer\textquotedbl{}]
}

{\footnotesize Для автентифікації на основі паролів підтримується та оновлюється список часто використовуваних, очікуваних або скомпрометованих паролів частота, а також коли є підозра, що паролі організації були скомпрометовані прямо чи опосередковано}


{\footnotesize\bfseries
No: 2\\
Name: ia\_5\_1\_b\\
Type: list\\
Default: [\textquotedbl{}admin\textquotedbl{}, \textquotedbl{}security\_officer\textquotedbl{}]
}

{\footnotesize Для автентифікації на основі паролів, коли паролі створюються або оновлюються користувачами, паролі перевіряються на відсутність у списку загальновживаних, очікуваних або скомпрометованих паролів в IA-05(01)(a)}


{\footnotesize\bfseries
No: 3\\
Name: ia\_5\_1\_c\\
Type: list\\
Default: [\textquotedbl{}admin\textquotedbl{}, \textquotedbl{}security\_officer\textquotedbl{}]
}

{\footnotesize Для автентифікації на основі паролів, паролі передаються лише криптографічно захищеними каналами}


{\footnotesize\bfseries
No: 4\\
Name: ia\_5\_1\_d\\
Type: list\\
Default: [\textquotedbl{}admin\textquotedbl{}, \textquotedbl{}security\_officer\textquotedbl{}]
}

{\footnotesize Для автентифікації на основі паролів паролі зберігаються за допомогою затвердженого алгоритму гешування, переважно використовуючи ключову геш-функцію}


{\footnotesize\bfseries
No: 5\\
Name: ia\_5\_1\_e\\
Type: string\\
Default: nil
}

{\footnotesize Для автентифікації на основі пароля після відновлення облікового запису потрібно негайно вибрати новий пароль}


{\footnotesize\bfseries
No: 6\\
Name: ia\_5\_1\_f\\
Type: list\\
Default: [\textquotedbl{}admin\textquotedbl{}, \textquotedbl{}security\_officer\textquotedbl{}]
}

{\footnotesize Для автентифікації на основі пароля дозволяється вибір користувачем довгих паролів і фраз, що включають пробіли та всі друковані символи}


{\footnotesize\bfseries
No: 7\\
Name: ia\_5\_1\_g\\
Type: list\\
Default: [\textquotedbl{}admin\textquotedbl{}, \textquotedbl{}security\_officer\textquotedbl{}]
}

{\footnotesize Для автентифікації на основі пароля використовуються автоматизовані інструменти, які допомагають користувачеві у виборі надійних автентифікаторів паролів}


{\footnotesize\bfseries
No: 8\\
Name: ia\_5\_1\_h\\
Type: list\\
Default: [\textquotedbl{}default\_deny\_rule\textquotedbl{}, \textquotedbl{}abac\_rule\_1\textquotedbl{}]
}

{\footnotesize Для автентифікації на основі пароля застосовуються склад та правила складності}


{\footnotesize\bfseries
No: 9\\
Name: ia\_5\_1\_odp\_01\\
Type: list\\
Default: [\textquotedbl{}admin\textquotedbl{}, \textquotedbl{}security\_officer\textquotedbl{}]
}

{\footnotesize Визначено частоту оновлення списку часто використовуваних, очікуваних або скомпрометованих паролів}


{\footnotesize\bfseries
No: 10\\
Name: ia\_5\_1\_odp\_02\\
Type: list\\
Default: [\textquotedbl{}default\_deny\_rule\textquotedbl{}, \textquotedbl{}abac\_rule\_1\textquotedbl{}]
}

{\footnotesize Визначено склад та правила складності автентифікатора}




\subsection{ПРИХОВУВАННЯ ЗВОРОТНОГО ЗВ'ЯЗКУ АВТЕНТИФІКАТОРА}


\vspace{10pt}
{\footnotesize\bfseries
No: 1\\
Name: ia\_6\_01\\
Type: list\\
Default: [\textquotedbl{}admin\textquotedbl{}, \textquotedbl{}security\_officer\textquotedbl{}]
}

{\footnotesize Забезпечино приховану зворотну передачу інформації автентифікації в про- цесі автентифікації для забезпечення захисту інформації від можливої експлуатації та використання неавторизованими особами}




\subsection{ПОВТОРНА АВТЕНТИФІКАЦІЯ (IA-11)}


\vspace{10pt}
{\footnotesize\bfseries
No: 1\\
Name: ia\_11\_01\\
Type: string\\
Default: nil
}

{\footnotesize Користувачі повинні повторно автентифікуватися, коли обставини або ситуації}


{\footnotesize\bfseries
No: 2\\
Name: ia\_11\_odp\\
Type: string\\
Default: nil
}

{\footnotesize Визначено обставини або ситуації, що вимагають повторної автен- тифікації}




\section{IR}
Клас заходів захисту IR --- РЕАГУВАННЯ НА ІНЦИДЕНТИ

Цей клас визначає процедури виявлення, аналізу, локалізації та усунення наслідків інцидентів інформаційної безпеки.

\textbf{Перелік заходів захисту:} Політика та процедури реагування на інциденти (IR-1); Навчання з реагування на інциденти (IR-2); Перевірка реагувань на інциденти (IR-3); Обробка інциденту (IR-4); Моніторинг інциденту (IR-5); Звітність про інциденти (IR-6); Підтримка реагування на інциденти (IR-7); План реагування на інциденти (IR-8).

\subsection{Політика та процедури реагування на інциденти (IR-1)}


\vspace{10pt}
{\footnotesize\bfseries
No: 1\\
Name: ir\_1\_a\_01\\
Type: list\\
Default: [\textquotedbl{}default\_deny\_rule\textquotedbl{}, \textquotedbl{}abac\_rule\_1\textquotedbl{}]
}

{\footnotesize Розроблено та задокументовано політику реагування на інциденти}


{\footnotesize\bfseries
No: 2\\
Name: ir\_1\_a\_02\\
Type: list\\
Default: [\textquotedbl{}admin\textquotedbl{}, \textquotedbl{}security\_officer\textquotedbl{}]
}

{\footnotesize Політика реагування на інциденти поширюється серед персоналу або ролей}


{\footnotesize\bfseries
No: 3\\
Name: ir\_1\_a\_03\\
Type: list\\
Default: [\textquotedbl{}default\_deny\_rule\textquotedbl{}, \textquotedbl{}abac\_rule\_1\textquotedbl{}]
}

{\footnotesize Розроблені та задокументовані процедури реагування на інциденти, що сприяють впровадженню політики реагування на інциденти та пов'язаних з нею заходів захисту з реагування на інциденти}


{\footnotesize\bfseries
No: 4\\
Name: ir\_1\_b\\
Type: list\\
Default: [\textquotedbl{}admin\textquotedbl{}, \textquotedbl{}security\_officer\textquotedbl{}]
}

{\footnotesize Посадова особа призначається для управління розробкою, документуванням та розповсюдженням політики та процедур реагування на інциденти; IR-01(c)[01][01] переглядається та оновлюється поточна політика реагування на інциденти частота; IR-01(c)[01][02] поточна політика реагування на інциденти переглядається та оновлюється після подій; IR-01(c)[02][01] переглядаються та оновлюються поточні процедури реагування на інциденти частота; IR-01(c)[02][02] поточні процедури реагування на інциденти переглядаються та оновлюються після подій}


{\footnotesize\bfseries
No: 5\\
Name: ir\_1\_odp\_01\\
Type: list\\
Default: [\textquotedbl{}admin\textquotedbl{}, \textquotedbl{}security\_officer\textquotedbl{}]
}

{\footnotesize Визначено персонал або ролі, до яких має бути доведена політика реагування на інциденти}


{\footnotesize\bfseries
No: 6\\
Name: ir\_1\_odp\_02\\
Type: list\\
Default: [\textquotedbl{}admin\textquotedbl{}, \textquotedbl{}security\_officer\textquotedbl{}]
}

{\footnotesize Визначено персонал або ролі, до яких мають бути доведені процедури реагування на інциденти}


{\footnotesize\bfseries
No: 7\\
Name: ir\_1\_odp\_03\\
Type: string\\
Default: nil
}

{\footnotesize Вибрано одне або декілька з наступних ЗНАЧЕНЬ ПАРАМЕТРІВ: \{рівень організації; рівень місії/бізнес-процесу; рівень системи\}}


{\footnotesize\bfseries
No: 8\\
Name: ir\_1\_odp\_04\\
Type: list\\
Default: [\textquotedbl{}admin\textquotedbl{}, \textquotedbl{}security\_officer\textquotedbl{}]
}

{\footnotesize Визначено посадову особу, яка керуватиме політикою та процедурами реагування на інциденти}


{\footnotesize\bfseries
No: 9\\
Name: ir\_1\_odp\_05\\
Type: list\\
Default: [\textquotedbl{}default\_deny\_rule\textquotedbl{}, \textquotedbl{}abac\_rule\_1\textquotedbl{}]
}

{\footnotesize Визначено частоту, з якою переглядається та оновлюється поточна політика реагування на інциденти}


{\footnotesize\bfseries
No: 10\\
Name: ir\_1\_odp\_06\\
Type: list\\
Default: [\textquotedbl{}default\_deny\_rule\textquotedbl{}, \textquotedbl{}abac\_rule\_1\textquotedbl{}]
}

{\footnotesize Визначаються події, які потребують перегляду та оновлення поточної політики реагування на інциденти}


{\footnotesize\bfseries
No: 11\\
Name: ir\_1\_odp\_07\\
Type: integer\\
Default: 30
}

{\footnotesize Визначено частоту, з якою переглядаються та оновлюються поточні процедури реагування на інциденти}


{\footnotesize\bfseries
No: 12\\
Name: ir\_1\_odp\_08\\
Type: list\\
Default: [\textquotedbl{}login\textquotedbl{}, \textquotedbl{}logout\textquotedbl{}, \textquotedbl{}failed\_attempt\textquotedbl{}]
}

{\footnotesize Визначено події, які потребують перегляду та оновлення процедур реагування на інциденти}




\subsection{НАВЧАННЯ З РЕАГУВАННЯ НА ІНЦИДЕНТИ (IR-2)}


\vspace{10pt}
{\footnotesize\bfseries
No: 1\\
Name: ir\_2\_a\_01\\
Type: list\\
Default: [\textquotedbl{}admin\textquotedbl{}, \textquotedbl{}security\_officer\textquotedbl{}]
}

{\footnotesize Навчання з реагування на інциденти надається користувачам системи відповідно до призначених ролей та обов'язків протягом періоду часу з моменту прийняття на себе ролі або обов'язків з реагування на інциденти або отримання доступу до системи}


{\footnotesize\bfseries
No: 2\\
Name: ir\_2\_a\_02\\
Type: string\\
Default: nil
}

{\footnotesize Навчання з реагування на інциденти надається користувачам системи відповідно до призначених ролей та обов'язків, коли цього вимагають зміни в системі}


{\footnotesize\bfseries
No: 3\\
Name: ir\_2\_a\_03\\
Type: string\\
Default: \textquotedbl{}щорічно\textquotedbl{}
}

{\footnotesize Користувачам системи надається навчання з реагування на інциденти відповідно до призначених ролей та обов'язків частота}


{\footnotesize\bfseries
No: 4\\
Name: ir\_2\_b\_01\\
Type: string\\
Default: \textquotedbl{}щорічно\textquotedbl{}
}

{\footnotesize Зміст навчання з реагування на інциденти переглядається та оновлюється частота}


{\footnotesize\bfseries
No: 5\\
Name: ir\_2\_b\_02\\
Type: string\\
Default: nil
}

{\footnotesize Зміст навчання з реагування на інциденти переглядається та оновлюється після подій}


{\footnotesize\bfseries
No: 6\\
Name: ir\_2\_odp\_02\\
Type: integer\\
Default: 30
}

{\footnotesize Визначено частоту, з якою користувачі повинні проходити навчання з реагування на інциденти}


{\footnotesize\bfseries
No: 7\\
Name: ir\_2\_odp\_03\\
Type: integer\\
Default: 30
}

{\footnotesize Визначено частоту перегляду та оновлення змісту навчання з реагування на інциденти}


{\footnotesize\bfseries
No: 8\\
Name: ir\_2\_odp\_04\\
Type: list\\
Default: [\textquotedbl{}login\textquotedbl{}, \textquotedbl{}logout\textquotedbl{}, \textquotedbl{}failed\_attempt\textquotedbl{}]
}

{\footnotesize Визначено події, які ініціюють перегляд змісту навчання з реагування на інциденти}




\subsection{ПЕРЕВІРКА РЕАГУВАНЬ НА ІНЦИДЕНТИ (IR-3)}


\vspace{10pt}
{\footnotesize\bfseries
No: 1\\
Name: ir\_3\_01\\
Type: string\\
Default: \textquotedbl{}щорічно\textquotedbl{}
}

{\footnotesize Ефективність реагування системи на інциденти перевіряється частота за допомогою тестів}


{\footnotesize\bfseries
No: 2\\
Name: ir\_3\_odp\_01\\
Type: integer\\
Default: 30
}

{\footnotesize Визначено частоту, з якою необхідно перевіряти ефективність реагування системи на інциденти}


{\footnotesize\bfseries
No: 3\\
Name: ir\_3\_odp\_02\\
Type: string\\
Default: nil
}

{\footnotesize Визначено тести, що використовуються для перевірки ефективності реагування на інциденти в системі}




\subsection{Обробка інциденту (IR-4)}


\vspace{10pt}
{\footnotesize\bfseries
No: 1\\
Name: ir\_4\_a\_02\\
Type: string\\
Default: nil
}

{\footnotesize Впроваджено можливість обробки інцидентів безпеки включно з виявленням}


{\footnotesize\bfseries
No: 2\\
Name: ir\_4\_a\_03\\
Type: string\\
Default: nil
}

{\footnotesize Впроваджено можливість обробки інцидентів безпеки включно з аналізом}


{\footnotesize\bfseries
No: 3\\
Name: ir\_4\_a\_04\\
Type: string\\
Default: nil
}

{\footnotesize Впроваджено можливість обробки інцидентів безпеки включно з локалізацією}


{\footnotesize\bfseries
No: 4\\
Name: ir\_4\_a\_05\\
Type: string\\
Default: nil
}

{\footnotesize Впроваджено можливість обробки інцидентів безпеки включно з ліквідацією}


{\footnotesize\bfseries
No: 5\\
Name: ir\_4\_a\_06\\
Type: string\\
Default: nil
}

{\footnotesize Впроваджено можливість обробки інцидентів безпеки включно з відновленням}


{\footnotesize\bfseries
No: 6\\
Name: ir\_4\_b\\
Type: string\\
Default: nil
}

{\footnotesize Діяльність з обробки інцидентів координується із заходами із забезпечення безперервності функціонування}


{\footnotesize\bfseries
No: 7\\
Name: ir\_4\_c\_01\\
Type: string\\
Default: nil
}

{\footnotesize Уроки, отримані з поточних дій з обробки інцидентів, включаються в процедури реагування на інциденти, навчання та тестування}


{\footnotesize\bfseries
No: 8\\
Name: ir\_4\_c\_02\\
Type: string\\
Default: nil
}

{\footnotesize Зміни, що випливають відповідним чином}


{\footnotesize\bfseries
No: 9\\
Name: ir\_4\_d\_01\\
Type: string\\
Default: nil
}

{\footnotesize Строгість заходів з обробки інцидентів передбачуваною в межах всієї організації}


{\footnotesize\bfseries
No: 10\\
Name: ir\_4\_d\_02\\
Type: string\\
Default: nil
}

{\footnotesize Інтенсивність заходів з обробки інцидентів є порівнянною та передбачуваною в межах всієї організації}


{\footnotesize\bfseries
No: 11\\
Name: ir\_4\_d\_03\\
Type: string\\
Default: nil
}

{\footnotesize Обсяг заходів з обробки інцидентів є порівнянною та передбачуваною в межах всієї організації}


{\footnotesize\bfseries
No: 12\\
Name: ir\_4\_d\_04\\
Type: string\\
Default: nil
}

{\footnotesize Результати діяльності заходів з обробки інцидентів є порівнянною та передбачуваною в межах всієї організації; з отриманих уроків, є впроваджуються порівнянною}




\subsection{Моніторинг інциденту (IR-5)}


\vspace{10pt}
{\footnotesize\bfseries
No: 1\\
Name: ir\_5\_01\\
Type: string\\
Default: nil
}

{\footnotesize Відстежуються інциденти безпеки та приватності}


{\footnotesize\bfseries
No: 2\\
Name: ir\_5\_02\\
Type: string\\
Default: nil
}

{\footnotesize Документуються інциденти безпеки та приватності. ПОТЕНЦІЙНІ МЕТОДИ ТА ОБ'ЄКТИ ОЦІНКИ:}




\subsection{ЗВІТНІСТЬ ПРО ІНЦИДЕНТИ (IR-6)}


\vspace{10pt}
{\footnotesize\bfseries
No: 1\\
Name: ir\_6\_a\\
Type: list\\
Default: [\textquotedbl{}admin\textquotedbl{}, \textquotedbl{}security\_officer\textquotedbl{}]
}

{\footnotesize Персонал зобов'язаний повідомляти про підозрілі інциденти протягом періоду часу}


{\footnotesize\bfseries
No: 2\\
Name: ir\_6\_b\\
Type: string\\
Default: nil
}

{\footnotesize Інформацію про інцидент повідомляється органам}


{\footnotesize\bfseries
No: 3\\
Name: ir\_6\_odp\_01\\
Type: list\\
Default: [\textquotedbl{}admin\textquotedbl{}, \textquotedbl{}security\_officer\textquotedbl{}]
}

{\footnotesize Визначено період часу, протягом якого персонал повинен повідомляти про підозрілі інциденти до уповноваженого органу}


{\footnotesize\bfseries
No: 4\\
Name: ir\_6\_odp\_02\\
Type: string\\
Default: nil
}

{\footnotesize Визначені органи, до яких слід повідомляти інформацію про інцидент}




\subsection{ПІДТРИМКА РЕАГУВАННЯ НА ІНЦИДЕНТИ (IR-7)}


\vspace{10pt}
{\footnotesize\bfseries
No: 1\\
Name: ir\_7\_01\\
Type: string\\
Default: nil
}

{\footnotesize Ресурс підтримки реагування на інциденти містить поради та допомогу користувачам системи для обробки та формування звітності про інциденти}




\subsection{План реагування на інциденти (IR-8)}


\vspace{10pt}
{\footnotesize\bfseries
No: 1\\
Name: ir\_8\_a\_01\\
Type: string\\
Default: nil
}

{\footnotesize Розроблено план реагування на інциденти, який надає організації дорожню карту для впровадження її можливостей реагування на інциденти}


{\footnotesize\bfseries
No: 2\\
Name: ir\_8\_a\_02\\
Type: string\\
Default: nil
}

{\footnotesize Розроблено план реагування на інциденти, який описує структуру та організацію спроможності реагування на інциденти}


{\footnotesize\bfseries
No: 3\\
Name: ir\_8\_a\_03\\
Type: string\\
Default: nil
}

{\footnotesize Розроблено план реагування на інциденти, який надає високорівневий підхід до того, як здатність реагування на інциденти вписується в загальну практику організації}


{\footnotesize\bfseries
No: 4\\
Name: ir\_8\_a\_04\\
Type: string\\
Default: nil
}

{\footnotesize Розроблено план реагування на інциденти, який відповідає унікальним вимогам організації, які пов'язані із завданнями, розміром, структурою і функціями}


{\footnotesize\bfseries
No: 5\\
Name: ir\_8\_a\_05\\
Type: string\\
Default: nil
}

{\footnotesize Розроблено план реагування на інциденти, який визначає підзвітні інциденти}


{\footnotesize\bfseries
No: 6\\
Name: ir\_8\_a\_06\\
Type: string\\
Default: nil
}

{\footnotesize Розроблено план реагування на інциденти, який надає показники для вимірювання можливостей реагування на інциденти всередині організації}


{\footnotesize\bfseries
No: 7\\
Name: ir\_8\_a\_07\\
Type: string\\
Default: nil
}

{\footnotesize Розроблено план реагування на інциденти, який визначає ресурси та управлінську підтримку, необхідну для ефективної підтримки та розвитку можливостей реагування на інциденти}


{\footnotesize\bfseries
No: 8\\
Name: ir\_8\_a\_08\\
Type: string\\
Default: nil
}

{\footnotesize Розроблено план реагування на інциденти, який вирішує питання обміну інформацією про інциденти}


{\footnotesize\bfseries
No: 9\\
Name: ir\_8\_a\_09\\
Type: list\\
Default: [\textquotedbl{}admin\textquotedbl{}, \textquotedbl{}security\_officer\textquotedbl{}]
}

{\footnotesize Розроблено план реагування на інцидент, який розглядається та затверджується персоналом або ролями частота}


{\footnotesize\bfseries
No: 10\\
Name: ir\_8\_a\_10\\
Type: list\\
Default: [\textquotedbl{}admin\textquotedbl{}, \textquotedbl{}security\_officer\textquotedbl{}]
}

{\footnotesize Розроблено план реагування на інциденти, в якому чітко визначено відповідальність за реагування на інциденти для організацій, персоналу або ролей}


{\footnotesize\bfseries
No: 11\\
Name: ir\_8\_b\_01\\
Type: list\\
Default: [\textquotedbl{}admin\textquotedbl{}, \textquotedbl{}security\_officer\textquotedbl{}]
}

{\footnotesize Копії плану реагування на інцидент розповсюджуються серед <IR- 08\_ODP[04] персоналу з реагування на інциденти>}


{\footnotesize\bfseries
No: 12\\
Name: ir\_8\_b\_02\\
Type: string\\
Default: nil
}

{\footnotesize Копії плану реагування на інцидент розповсюджуються серед елементів організації}


{\footnotesize\bfseries
No: 13\\
Name: ir\_8\_c\\
Type: integer\\
Default: 30
}

{\footnotesize План реагування на інциденти оновлюється з урахуванням змін у системі та організації або проблем, що виникають під час впровадження, виконання або тестування плану}


{\footnotesize\bfseries
No: 14\\
Name: ir\_8\_d\_01\\
Type: list\\
Default: [\textquotedbl{}admin\textquotedbl{}, \textquotedbl{}security\_officer\textquotedbl{}]
}

{\footnotesize Зміни в плані реагування на інцидент повідомляються персоналу з реагування на інциденти}


{\footnotesize\bfseries
No: 15\\
Name: ir\_8\_d\_02\\
Type: string\\
Default: nil
}

{\footnotesize Зміни в плані реагування на інциденти надсилаються до елементів організації}


{\footnotesize\bfseries
No: 16\\
Name: ir\_8\_e\_01\\
Type: string\\
Default: nil
}

{\footnotesize План реагування на інциденти захищений від несанкціонованого розкриття}


{\footnotesize\bfseries
No: 17\\
Name: ir\_8\_e\_02\\
Type: string\\
Default: nil
}

{\footnotesize План реагування на інциденти захищений від несанкціонованої модифікації}


{\footnotesize\bfseries
No: 18\\
Name: ir\_8\_odp\_01\\
Type: list\\
Default: [\textquotedbl{}admin\textquotedbl{}, \textquotedbl{}security\_officer\textquotedbl{}]
}

{\footnotesize Визначено персонал або ролі, які переглядають та затверджують план реагування на інциденти}


{\footnotesize\bfseries
No: 19\\
Name: ir\_8\_odp\_02\\
Type: string\\
Default: \textquotedbl{}щорічно\textquotedbl{}
}

{\footnotesize Визначено періодичність перегляду та затвердження плану реагування на інциденти}


{\footnotesize\bfseries
No: 20\\
Name: ir\_8\_odp\_03\\
Type: list\\
Default: [\textquotedbl{}admin\textquotedbl{}, \textquotedbl{}security\_officer\textquotedbl{}]
}

{\footnotesize Визначені організації, персонал або ролі, які несуть відповідальність за реагування на інциденти; IR-08\_ODP[04] визначено персонал з реагування на інцидент (ідентифікований за іменами та/або за ролями), якому мають бути роздані копії плану реагування на інцидент}


{\footnotesize\bfseries
No: 21\\
Name: ir\_8\_odp\_05\\
Type: string\\
Default: nil
}

{\footnotesize Визначено елементи організації, серед яких мають розповсюджені копії плану реагування на інцидент; бути}


{\footnotesize\bfseries
No: 22\\
Name: ir\_8\_odp\_06\\
Type: list\\
Default: [\textquotedbl{}admin\textquotedbl{}, \textquotedbl{}security\_officer\textquotedbl{}]
}

{\footnotesize Визначено персонал з реагування на інцидент (ідентифікований за іменами та/або ролями), якому повідомляються зміни до плану реагування на інцидент}


{\footnotesize\bfseries
No: 23\\
Name: ir\_8\_odp\_07\\
Type: string\\
Default: nil
}

{\footnotesize Визначено елементи організації, яким повідомляється про зміни в плані реагування на інцидент}




\section{MA}
Клас заходів захисту MA --- ТЕХНІЧНЕ ОБСЛУГОВУВАННЯ

Цей клас регулює процеси планового та позапланового технічного обслуговування компонентів системи для запобігання збоям.

\textbf{Перелік заходів захисту:} Політика та процедури технічного обслуговування (MA-1); Інструменти для обслуговування (MA-3); Перевірка інструментів (MA-3(1)); Перевірка носіїв інформації (MA-3(2)); Запобігання несанкціонованому переміщенню (MA-3(3)); Віддалене обслуговування (MA-4); Технічний персонал (MA-5).

\subsection{ПОЛІТИКА ТА ПРОЦЕДУРИ ТЕХНІЧНОГО ОБСЛУГОВУВАННЯ (MA-1)}
a. Планувати, документувати та переглядати записи з технічного обслуговування, ремонту або заміни компонентів системи відповідно до вимог виробника та постачальників та/або вимог організації.\\ 
b. Затвердити та здійснювати моніторинг усіх заходів з технічного обслуговування, незалежно від того, виконуються вони на місці або віддалено, а також чи обслуговуються системи або системні компоненти на місці, чи переміщуються в інше місце.\\ 
c. Вимагати, щоб [Призначення: визначені організацією персонал чи ролі] явно схвалили видалення системи або компоненту системи з організаційного обладнання для технічного обслуговування, ремонту чи заміни поза об'єктами експлуатації.\\ 
d. Очищати обладнання з погляду видалення всієї інформації з носіїв до вилучення обладнання організації для технічного обслуговування, ремонту чи заміни поза об'єктами експлуатації.\\ 
e. Перевірити всі потенційно порушені заходи захисту, щоб переконатися, що вони, як і раніше, працюють належним чином після дій з обслуговування, ремонту або заміни.\\ 
f. Вносити [Призначення: визначену організацією інформацію, пов'язану з технічним обслуговуванням] до записів з технічного обслуговування.

\vspace{10pt}
{\footnotesize\bfseries
No: 1\\
Name: ma\_1\_a\_01\\
Type: list\\
Default: [\textquotedbl{}default\_deny\_rule\textquotedbl{}, \textquotedbl{}abac\_rule\_1\textquotedbl{}]
}

{\footnotesize Розроблено та задокументовано політику технічного обслуговування}


{\footnotesize\bfseries
No: 2\\
Name: ma\_1\_a\_02\\
Type: list\\
Default: [\textquotedbl{}admin\textquotedbl{}, \textquotedbl{}security\_officer\textquotedbl{}]
}

{\footnotesize Політика технічного обслуговування поширюється на персонал або ролі}


{\footnotesize\bfseries
No: 3\\
Name: ma\_1\_a\_03\\
Type: list\\
Default: [\textquotedbl{}default\_deny\_rule\textquotedbl{}, \textquotedbl{}abac\_rule\_1\textquotedbl{}]
}

{\footnotesize Розроблені та задокументовані процедури технічного обслуговування, що сприяють впровадженню політики технічного обслуговування та пов'язаних з нею заходів технічного обслуговування}


{\footnotesize\bfseries
No: 4\\
Name: ma\_1\_b\\
Type: list\\
Default: [\textquotedbl{}admin\textquotedbl{}, \textquotedbl{}security\_officer\textquotedbl{}]
}

{\footnotesize Посадова особа призначається для управління розробкою, документуванням та розповсюдженням політики та процедур технічного обслуговування; MA-01(c)[01][01] переглядається та оновлюється поточна політика обслуговування частота; MA-01(c)[01][02] переглядається та оновлюється поточна політика обслуговування після подій; MA-01(c)[02][01] переглядаються та оновлюються поточні процедури технічного обслуговування частота; MA-01(c)[02][02] переглядаються та оновлюються поточні процедури технічного обслуговування після подій}


{\footnotesize\bfseries
No: 5\\
Name: ma\_1\_odp\_01\\
Type: list\\
Default: [\textquotedbl{}admin\textquotedbl{}, \textquotedbl{}security\_officer\textquotedbl{}]
}

{\footnotesize Визначено персонал або ролі, на які поширюється політика технічного обслуговування}


{\footnotesize\bfseries
No: 6\\
Name: ma\_1\_odp\_02\\
Type: list\\
Default: [\textquotedbl{}admin\textquotedbl{}, \textquotedbl{}security\_officer\textquotedbl{}]
}

{\footnotesize Визначено персонал або ролі, на які поширюються процедури технічного обслуговування}


{\footnotesize\bfseries
No: 7\\
Name: ma\_1\_odp\_03\\
Type: string\\
Default: nil
}

{\footnotesize Вибрано одне або декілька з наступних ЗНАЧЕНЬ ПАРАМЕТРІВ: \{рівень організації; рівень місії/бізнес-процесу; рівень системи\}}


{\footnotesize\bfseries
No: 8\\
Name: ma\_1\_odp\_04\\
Type: list\\
Default: [\textquotedbl{}admin\textquotedbl{}, \textquotedbl{}security\_officer\textquotedbl{}]
}

{\footnotesize Визначено посадову особу, яка керуватиме політикою та процедурами технічного обслуговування}


{\footnotesize\bfseries
No: 9\\
Name: ma\_1\_odp\_05\\
Type: list\\
Default: [\textquotedbl{}default\_deny\_rule\textquotedbl{}, \textquotedbl{}abac\_rule\_1\textquotedbl{}]
}

{\footnotesize Визначено частоту, з якою переглядається та оновлюється поточна політика технічного обслуговування}


{\footnotesize\bfseries
No: 10\\
Name: ma\_1\_odp\_06\\
Type: list\\
Default: [\textquotedbl{}default\_deny\_rule\textquotedbl{}, \textquotedbl{}abac\_rule\_1\textquotedbl{}]
}

{\footnotesize Визначено події, які потребують перегляду та оновлення поточної політики технічного обслуговування}


{\footnotesize\bfseries
No: 11\\
Name: ma\_1\_odp\_07\\
Type: integer\\
Default: 30
}

{\footnotesize Визначено частоту, з якою переглядаються та оновлюються поточні процедури технічного обслуговування}


{\footnotesize\bfseries
No: 12\\
Name: ma\_1\_odp\_08\\
Type: list\\
Default: [\textquotedbl{}login\textquotedbl{}, \textquotedbl{}logout\textquotedbl{}, \textquotedbl{}failed\_attempt\textquotedbl{}]
}

{\footnotesize Визначено події, які потребують перегляду та оновлення процедур технічного обслуговування}




\subsection{ІНСТРУМЕНТИ ДЛЯ ОБСЛУГОВУВАННЯ (MA-3)}
a. Затвердити, контролювати та відстежувати використання засобів технічного обслуговування.\\ 
b. Переглядати раніше затверджені інструменти технічного [Призначення: з частотою, визначеною організацією]. обслуговування

\vspace{10pt}
{\footnotesize\bfseries
No: 1\\
Name: ma\_3\_a\_01\\
Type: string\\
Default: nil
}

{\footnotesize Використання засобів технічного обслуговування затверджено}


{\footnotesize\bfseries
No: 2\\
Name: ma\_3\_a\_02\\
Type: string\\
Default: nil
}

{\footnotesize Використання засобів технічного обслуговування контролюється}


{\footnotesize\bfseries
No: 3\\
Name: ma\_3\_a\_03\\
Type: string\\
Default: nil
}

{\footnotesize Використання засобів технічного обслуговування відстажуються}


{\footnotesize\bfseries
No: 4\\
Name: ma\_3\_b\\
Type: string\\
Default: \textquotedbl{}щорічно\textquotedbl{}
}

{\footnotesize Переглядаються раніше затверджені інструменти технічного обслуговування частота}


{\footnotesize\bfseries
No: 5\\
Name: ma\_3\_odp\\
Type: integer\\
Default: 30
}

{\footnotesize Визначено частоту, з якою слід переглядати раніше затверджені інструменти технічного обслуговування}




\subsubsection{ПЕРЕВІРКА ІНСТРУМЕНТІВ (MA-3(1))}


\vspace{10pt}
{\footnotesize\bfseries
No: 1\\
Name: ma\_3\_1\_01\\
Type: list\\
Default: [\textquotedbl{}admin\textquotedbl{}, \textquotedbl{}security\_officer\textquotedbl{}]
}

{\footnotesize Оглядаються інструменти для технічного обслуговування, які доставлені на об'єкт обслуговуючим персоналом, на предмет неправильних або несанкціонованих модифікацій}




\subsubsection{ПЕРЕВІРКА НОСІЇВ ІНФОРМАЦІЇ (MA-3(2))}
Перед використанням носіїв у системі перевірити носії, що містять діагностичні та тестові програми на наявність шкідливого коду.

\vspace{10pt}
{\footnotesize\bfseries
No: 1\\
Name: ma\_3\_2\_01\\
Type: string\\
Default: nil
}

{\footnotesize Перед використанням носіїв у системі перевіряються носії, що містять діагностичні та тестові програми на наявність шкідливого коду}




\subsubsection{ЗАПОБІГАННЯ НЕСАНКЦІОНОВАНОМУ ПЕРЕМІЩЕННЮ (MA-3(3))}
Запобігти переміщенню обладнання для технічного обслуговування, що містить організаційну інформацію, шляхом:\\ 
(a) перевірки відсутності організаційної інформації, розміщеної на обладнанні;\\ 
(b) очищення або знищення обладнання;\\ 
(c) утримання обладнання на об'єкті;\\ 
(d) отримання дозволу від [Призначення: визначених організацією персоналу чи ролей], які явно дозволяють переміщення обладнання з об'єкта.

\vspace{10pt}
{\footnotesize\bfseries
No: 1\\
Name: ma\_3\_3\_a\\
Type: string\\
Default: nil
}

{\footnotesize Переміщення обладнання для технічного обслуговування, що містить інформацію організації, запобігається шляхом перевірки того, що на обладнанні не міститься ніякої інформації організації; або}


{\footnotesize\bfseries
No: 2\\
Name: ma\_3\_3\_b\\
Type: string\\
Default: nil
}

{\footnotesize Переміщення обладнання для технічного обслуговування, що містить інформацію організації, запобігається шляхом очищення або знищення обладнання; або}


{\footnotesize\bfseries
No: 3\\
Name: ma\_3\_3\_c\\
Type: string\\
Default: nil
}

{\footnotesize Переміщення обладнання для технічного обслуговування, що містить інформацію організації, запобігається шляхом утримання обладнання на об'єкті; або}


{\footnotesize\bfseries
No: 4\\
Name: ma\_3\_3\_d\\
Type: list\\
Default: [\textquotedbl{}admin\textquotedbl{}, \textquotedbl{}security\_officer\textquotedbl{}]
}

{\footnotesize Переміщення обладнання для технічного обслуговування, що містить інформацію організації, запобігається шляхом отримання дозволу від персонал або ролі}


{\footnotesize\bfseries
No: 5\\
Name: ma\_3\_3\_odp\\
Type: list\\
Default: [\textquotedbl{}admin\textquotedbl{}, \textquotedbl{}security\_officer\textquotedbl{}]
}

{\footnotesize Визначено персонал або ролі, які можуть надавати дозвіл на переміщення обладнання з об'єкту}




\subsection{ВІДДАЛЕНЕ ОБСЛУГОВУВАННЯ (MA-4)}


\vspace{10pt}
{\footnotesize\bfseries
No: 1\\
Name: ma\_4\_a\_01\\
Type: list\\
Default: [\textquotedbl{}login\textquotedbl{}, \textquotedbl{}logout\textquotedbl{}, \textquotedbl{}failed\_attempt\textquotedbl{}]
}

{\footnotesize Впроваджено віддалені дії з обслуговування та діагностики}


{\footnotesize\bfseries
No: 2\\
Name: ma\_4\_a\_02\\
Type: list\\
Default: [\textquotedbl{}login\textquotedbl{}, \textquotedbl{}logout\textquotedbl{}, \textquotedbl{}failed\_attempt\textquotedbl{}]
}

{\footnotesize Відстежуються віддалені дії з обслуговування та діагностики}


{\footnotesize\bfseries
No: 3\\
Name: ma\_4\_b\_01\\
Type: list\\
Default: [\textquotedbl{}default\_deny\_rule\textquotedbl{}, \textquotedbl{}abac\_rule\_1\textquotedbl{}]
}

{\footnotesize Використання віддалених засобів технічного обслуговування та діагностики дозволено лише відповідно до політики організації}


{\footnotesize\bfseries
No: 4\\
Name: ma\_4\_b\_02\\
Type: string\\
Default: nil
}

{\footnotesize Використання віддалених засобів технічного обслуговування та діагностики задокументовано в плані захисту інформації}


{\footnotesize\bfseries
No: 5\\
Name: ma\_4\_c\\
Type: string\\
Default: nil
}

{\footnotesize Надійна автентифікація використовується при встановленні віддалених технічних та діагностичних сеансів}


{\footnotesize\bfseries
No: 6\\
Name: ma\_4\_d\\
Type: string\\
Default: nil
}

{\footnotesize Ведеться облік віддалених дій з обслуговування та діагностики}


{\footnotesize\bfseries
No: 7\\
Name: ma\_4\_e\_01\\
Type: string\\
Default: nil
}

{\footnotesize Сесія припиняється, коли завершено віддалене обслуговування}


{\footnotesize\bfseries
No: 8\\
Name: ma\_4\_e\_02\\
Type: string\\
Default: nil
}

{\footnotesize Мережеве з'єднання припиняється, коли завершено віддалене обслуговування}




\subsection{ТЕХНІЧНИЙ ПЕРСОНАЛ (MA-5)}
a. Встановити процедуру авторизації технічного персоналу та вести перелік авторизованих організацій технічного обслуговування або персоналу.\\ 
b. Перевіряти, що персонал, який не супроводжується та виконує технічне обслуговування в системі, має необхідні дозволи на доступ.\\ 
c. Визначити персонал організації з необхідними повноваженнями щодо доступу та технічною компетенцією для нагляду за персоналом з технічного обслуговування, який не має необхідних дозволів на доступ.

\vspace{10pt}
{\footnotesize\bfseries
No: 1\\
Name: ma\_5\_a\_01\\
Type: list\\
Default: [\textquotedbl{}admin\textquotedbl{}, \textquotedbl{}security\_officer\textquotedbl{}]
}

{\footnotesize Запроваджено процес авторизації технічного персоналу}


{\footnotesize\bfseries
No: 2\\
Name: ma\_5\_a\_02\\
Type: list\\
Default: [\textquotedbl{}admin\textquotedbl{}, \textquotedbl{}security\_officer\textquotedbl{}]
}

{\footnotesize Ведеться перелік авторизованих організацій або персоналу з технічного обслуговування}


{\footnotesize\bfseries
No: 3\\
Name: ma\_5\_b\\
Type: list\\
Default: [\textquotedbl{}admin\textquotedbl{}, \textquotedbl{}security\_officer\textquotedbl{}]
}

{\footnotesize Персонал без супроводу, який виконує технічне обслуговування системи, має необхідні дозволи на доступ}


{\footnotesize\bfseries
No: 4\\
Name: ma\_5\_c\\
Type: list\\
Default: [\textquotedbl{}admin\textquotedbl{}, \textquotedbl{}security\_officer\textquotedbl{}]
}

{\footnotesize Персонал організації з необхідними повноваженнями доступу та технічною компетентністю призначений/призначені для нагляду за діяльністю з технічного обслуговування персоналу, який не має необхідних дозволів на доступ}




\section{MP}
Клас заходів захисту MP --- ЗАХИСТ НОСІЇВ ІНФОРМАЦІЇ

Цей клас спрямований на безпечне зберігання, транспортування, використання та знищення як цифрових, так і паперових носіїв інформації.

\textbf{Перелік заходів захисту:} Політика та процедури щодо захисту носіїв інформації (MP-1); Доступ до носіїв інформації (MP-2); Маркування носіїв інформації (MP-3); Зберігання носіїв інформації (MP-4); Транспортування носіїв інформації (MP-5); Знищення інформації на носіях інформації (MP-6); Використання носіїв інформації (MP-7).

\subsection{ПОЛІТИКА ТА ПРОЦЕДУРИ ЩОДО ЗАХИСТУ НОСІЇВ ІНФОРМАЦІЇ (MP-1)}
a. Розробити, задокументувати та поширити серед [Призначення: визначеного організацією персоналу або посад]:\\ 
1. 2. політику захисту носіїв інформації, яка:\\ 
(a) містить мету, сферу застосування, ролі, обов'язки, відповідальність керівництва, координацію між організаційними підрозділами та систему контролю відповідності (complaince);\\ 
(b) відповідає чинному законодавству, виконавчим наказам, директивам, нормам, політикам, стандартам та керівним принципам; процедури, які сприяють здійсненню політики та заходів захисту носіїв інформації.\\ 
b. Призначити [Призначення: визначену організацією посадову особу] для управління розробкою, документування, та розповсюдження політики та процедурами захисту носіїв інформації.\\ 
c. Переглядати та оновлювати чинну систему захисту носіїв інформації:\\ 
1. поточну політику захисту носіїв інформації [Призначення: з визначеною організацією частотою];\\ 
2. поточні процедури захисту носіїв інформації [Призначення: з визначеною організацією частотою].

\vspace{10pt}
{\footnotesize\bfseries
No: 1\\
Name: mp\_1\_a\_01\\
Type: list\\
Default: [\textquotedbl{}default\_deny\_rule\textquotedbl{}, \textquotedbl{}abac\_rule\_1\textquotedbl{}]
}

{\footnotesize Розроблено та задокументовано політику захисту носіїв інформації}


{\footnotesize\bfseries
No: 2\\
Name: mp\_1\_a\_02\\
Type: list\\
Default: [\textquotedbl{}admin\textquotedbl{}, \textquotedbl{}security\_officer\textquotedbl{}]
}

{\footnotesize Політика захисту носіїв інформації поширюється на персонал або ролі}


{\footnotesize\bfseries
No: 3\\
Name: mp\_1\_a\_03\\
Type: list\\
Default: [\textquotedbl{}default\_deny\_rule\textquotedbl{}, \textquotedbl{}abac\_rule\_1\textquotedbl{}]
}

{\footnotesize Розроблено та задокументовано процедури захисту носіїв інформації, що сприятимуть реалізації політики захисту носіїв інформації та заходів захисту носіїв інформації}


{\footnotesize\bfseries
No: 4\\
Name: mp\_1\_b\\
Type: list\\
Default: [\textquotedbl{}admin\textquotedbl{}, \textquotedbl{}security\_officer\textquotedbl{}]
}

{\footnotesize Посадова особа призначається для управління розробкою, документуванням та розповсюдженням політики та процедур захисту носіїв інформації. MP-01(c)[01][01] переглядається та оновлюється поточна політика захисту носіїв інформації частота; MP-01(c)[01][02] переглядається та оновлюється поточна політика захисту носіїв інформації після подій; MP-01(c)[02][01] переглядаються та оновлюються поточні процедури захисту носіїв інформації частота; MP-01(c)[02][02] переглядаються та оновлюються поточні процедури захисту носіїв інформації після подій}


{\footnotesize\bfseries
No: 5\\
Name: mp\_1\_odp\_01\\
Type: list\\
Default: [\textquotedbl{}admin\textquotedbl{}, \textquotedbl{}security\_officer\textquotedbl{}]
}

{\footnotesize Визначено персонал або ролі, серед яких має бути поширена політика захисту носіїв інформації}


{\footnotesize\bfseries
No: 6\\
Name: mp\_1\_odp\_02\\
Type: list\\
Default: [\textquotedbl{}admin\textquotedbl{}, \textquotedbl{}security\_officer\textquotedbl{}]
}

{\footnotesize Визначено персонал або ролі, серед яких мають бути поширені процедури захисту носіїв інформації}


{\footnotesize\bfseries
No: 7\\
Name: mp\_1\_odp\_03\\
Type: string\\
Default: nil
}

{\footnotesize Вибрано одне або декілька з наступних ЗНАЧЕНЬ ПАРАМЕТРІВ: \{рівень організації; рівень місії/бізнес-процесу; рівень системи\}}


{\footnotesize\bfseries
No: 8\\
Name: mp\_1\_odp\_04\\
Type: list\\
Default: [\textquotedbl{}admin\textquotedbl{}, \textquotedbl{}security\_officer\textquotedbl{}]
}

{\footnotesize Визначено посадову особу, яка керуватиме політикою та процедурами захисту носіїв інформації}


{\footnotesize\bfseries
No: 9\\
Name: mp\_1\_odp\_05\\
Type: list\\
Default: [\textquotedbl{}default\_deny\_rule\textquotedbl{}, \textquotedbl{}abac\_rule\_1\textquotedbl{}]
}

{\footnotesize Визначено частоту, з якою переглядається та оновлюється поточна політика захисту носіїв інформації}


{\footnotesize\bfseries
No: 10\\
Name: mp\_1\_odp\_06\\
Type: list\\
Default: [\textquotedbl{}default\_deny\_rule\textquotedbl{}, \textquotedbl{}abac\_rule\_1\textquotedbl{}]
}

{\footnotesize Визначено події, які потребують перегляду та оновлення чинної політики захисту носіїв інформації}


{\footnotesize\bfseries
No: 11\\
Name: mp\_1\_odp\_07\\
Type: integer\\
Default: 30
}

{\footnotesize Визначено частоту, з якою переглядаються та оновлюються чинні процедури захисту носіїв інформації}


{\footnotesize\bfseries
No: 12\\
Name: mp\_1\_odp\_08\\
Type: list\\
Default: [\textquotedbl{}login\textquotedbl{}, \textquotedbl{}logout\textquotedbl{}, \textquotedbl{}failed\_attempt\textquotedbl{}]
}

{\footnotesize Визначено події, які потребують перегляду та оновлення процедур захисту носіїв інформації}




\subsection{ДОСТУП ДО НОСІЇВ ІНФОРМАЦІЇ (MP-2)}
Обмежити доступ до [Призначення: визначених організацією типів цифрових та/або нецифрових носіїв інформації] [Призначення: визначеним організацією персоналом або ролями].

\vspace{10pt}
{\footnotesize\bfseries
No: 1\\
Name: mp\_2\_01\\
Type: list\\
Default: [\textquotedbl{}admin\textquotedbl{}, \textquotedbl{}security\_officer\textquotedbl{}]
}

{\footnotesize Доступ до типів цифрових носіїв інформації обмежено для персоналу або ролей}


{\footnotesize\bfseries
No: 2\\
Name: mp\_2\_02\\
Type: list\\
Default: [\textquotedbl{}admin\textquotedbl{}, \textquotedbl{}security\_officer\textquotedbl{}]
}

{\footnotesize Доступ до типів нецифрових носіїв інформації обмежено для персоналу або ролей}


{\footnotesize\bfseries
No: 3\\
Name: mp\_2\_odp\_01\\
Type: string\\
Default: nil
}

{\footnotesize Визначено типи цифрових носіїв інформації, доступ до яких обмежено}


{\footnotesize\bfseries
No: 4\\
Name: mp\_2\_odp\_02\\
Type: list\\
Default: [\textquotedbl{}admin\textquotedbl{}, \textquotedbl{}security\_officer\textquotedbl{}]
}

{\footnotesize Визначено персонал або ролі, уповноважені на доступ до цифрових носіїв інформації}


{\footnotesize\bfseries
No: 5\\
Name: mp\_2\_odp\_03\\
Type: string\\
Default: nil
}

{\footnotesize Визначено типи нецифрових носіїв інформації, доступ до яких обмежено}


{\footnotesize\bfseries
No: 6\\
Name: mp\_2\_odp\_04\\
Type: list\\
Default: [\textquotedbl{}admin\textquotedbl{}, \textquotedbl{}security\_officer\textquotedbl{}]
}

{\footnotesize Визначено персонал або ролі, уповноважені на доступ до нецифрових носіїв інформації}




\subsection{МАРКУВАННЯ НОСІЇВ ІНФОРМАЦІЇ (MP-3)}
a. Наносити на носії інформації маркування, що вказують на обмеження поширення, обробки, а також застереження та відповідні мітки безпеки (якщо такі є) інформації.\\ 
b. Звільнити [Призначення: визначені організацією типи носіїв системи] від маркування, якщо носії залишаються в межах [Призначення: визначених організацією контрольованих зон].

\vspace{10pt}
{\footnotesize\bfseries
No: 1\\
Name: mp\_3\_01\\
Type: string\\
Default: nil
}

{\footnotesize МАРКУВАННЯ НОСІЇВ ІНФОРМАЦІЇ МЕТА ОЦІНКИ: Визначити, чи:}


{\footnotesize\bfseries
No: 2\\
Name: mp\_3\_a\\
Type: string\\
Default: nil
}

{\footnotesize Носії інформації маркуються, щоб вказати на обмеження поширення, обробки, а також застереження та відповідні мітки безпеки (якщо такі є) інформації}


{\footnotesize\bfseries
No: 3\\
Name: mp\_3\_b\\
Type: string\\
Default: nil
}

{\footnotesize Типи носіїв інформації залишаються в межах контрольованих зон}


{\footnotesize\bfseries
No: 4\\
Name: mp\_3\_odp\_01\\
Type: integer\\
Default: 30
}

{\footnotesize Визначено типи носіїв інформації, які звільняються від маркування під час перебування на контрольованих зонах}


{\footnotesize\bfseries
No: 5\\
Name: mp\_3\_odp\_02\\
Type: string\\
Default: nil
}

{\footnotesize Визначено контрольовані зони, де носії інформації звільняються від маркування}




\subsection{ЗБЕРІГАННЯ НОСІЇВ ІНФОРМАЦІЇ (MP-4)}
a. Фізично контролювати та безпечно зберігати [Призначення: визначені організацією типи цифрових та/або нецифрових носіїв інформації] в межах [Призначення: визначених організацією контрольованих зон].\\ 
b. Захищати системні носії, які визначені в МР-4 до того часу, як носії знищуються або очищаються, з використанням затвердженого обладнання, методів та процедур.

\vspace{10pt}
{\footnotesize\bfseries
No: 1\\
Name: mp\_4\_a\_01\\
Type: string\\
Default: nil
}

{\footnotesize Типи цифрових носіїв контролюються фізично}


{\footnotesize\bfseries
No: 2\\
Name: mp\_4\_a\_02\\
Type: string\\
Default: nil
}

{\footnotesize Типи нецифрових носіїв контролюються фізично}


{\footnotesize\bfseries
No: 3\\
Name: mp\_4\_a\_03\\
Type: string\\
Default: nil
}

{\footnotesize Типи цифрових носіїв безпечно зберігаються в контрольованих зонах}


{\footnotesize\bfseries
No: 4\\
Name: mp\_4\_a\_04\\
Type: string\\
Default: nil
}

{\footnotesize Типи нецифрових носіїв безпечно зберігаються в контрольованих зонах}


{\footnotesize\bfseries
No: 5\\
Name: mp\_4\_b\\
Type: string\\
Default: nil
}

{\footnotesize Типи носіїв інформації (визначені в MP-04\_ODP[01], MP04\_ODP[02], MP-04\_ODP[03], MP-04\_ODP[04]) захищені доти, доки носії інформації не будуть знищені або очищені за допомогою визначеного обладнання, методик та процедур}


{\footnotesize\bfseries
No: 6\\
Name: mp\_4\_odp\_01\\
Type: string\\
Default: nil
}

{\footnotesize Визначено типи цифрових носіїв інформації, які підлягають фізичному контролю (якщо вибрано)}


{\footnotesize\bfseries
No: 7\\
Name: mp\_4\_odp\_02\\
Type: string\\
Default: nil
}

{\footnotesize Визначено типи нецифрових носіїв інформації, які підлягають фізичному контролю (якщо вибрано)}


{\footnotesize\bfseries
No: 8\\
Name: mp\_4\_odp\_03\\
Type: string\\
Default: nil
}

{\footnotesize Визначено типи цифрових носіїв інформації для безпечного зберігання (якщо вибрано)}


{\footnotesize\bfseries
No: 9\\
Name: mp\_4\_odp\_04\\
Type: string\\
Default: nil
}

{\footnotesize Визначено типи нецифрових носіїв інформації для безпечного зберігання (якщо вибрано)}


{\footnotesize\bfseries
No: 10\\
Name: mp\_4\_odp\_05\\
Type: string\\
Default: nil
}

{\footnotesize Визначено контрольовані зони, в яких можна безпечно зберігати цифрові носії інформації; MP-04\_ODP[06] визначено контрольовані зони, в яких можна безпечно зберігати нецифрові носії інформації}




\subsection{Транспортування носіїв інформації (MP-5)}
a. Захищати та контролювати [Призначення: визначені організацією типи носіїв системи] під час транспортування за межами контрольованих зон, використовуючи [Призначення: визначені організацією заходи безпеки].\\ 
b. Вести облік носіїв системи інформації під час транспортування за межами контрольованих зон.\\ 
c. Документувати дії, пов'язані з транспортуванням носіїв системи.\\ 
d. Обмежити діяльність уповноваженого персоналу, пов'язану з транспортуванням носіїв системи.

\vspace{10pt}
{\footnotesize\bfseries
No: 1\\
Name: mp\_5\_a\_01\\
Type: integer\\
Default: 30
}

{\footnotesize Типи носіїв інформації системи захищаються під час транспортування за межі контрольованих зон за допомогою заходів безпеки}


{\footnotesize\bfseries
No: 2\\
Name: mp\_5\_a\_02\\
Type: integer\\
Default: 30
}

{\footnotesize Типи носіїв інформації системи контролюються під час транспортування за межі контрольованих зон за допомогою заходів безпеки}


{\footnotesize\bfseries
No: 3\\
Name: mp\_5\_b\\
Type: integer\\
Default: 30
}

{\footnotesize Під час транспортування за межі контрольованих зон ведеться облік носіїв інформації системи}


{\footnotesize\bfseries
No: 4\\
Name: mp\_5\_c\\
Type: string\\
Default: nil
}

{\footnotesize Діяльність, пов'язана з транспортуванням носіїв інформації системи, задокументована}


{\footnotesize\bfseries
No: 5\\
Name: mp\_5\_d\_01\\
Type: list\\
Default: [\textquotedbl{}admin\textquotedbl{}, \textquotedbl{}security\_officer\textquotedbl{}]
}

{\footnotesize Визначено персонал, уповноважений здійснювати діяльність з транспортування носіїв інформації}


{\footnotesize\bfseries
No: 6\\
Name: mp\_5\_d\_02\\
Type: list\\
Default: [\textquotedbl{}admin\textquotedbl{}, \textquotedbl{}security\_officer\textquotedbl{}]
}

{\footnotesize Діяльність, пов'язана з транспортуванням носіїв інформації системи, обмежується визначеним уповноваженим персоналом}


{\footnotesize\bfseries
No: 7\\
Name: mp\_5\_odp\_01\\
Type: integer\\
Default: 30
}

{\footnotesize Визначено типи носіїв інформації системи для захисту та контролю під час транспортування за межі контрольованих зон}


{\footnotesize\bfseries
No: 8\\
Name: mp\_5\_odp\_02\\
Type: string\\
Default: nil
}

{\footnotesize Визначено заходи безпеки, що використовуються для захисту носіїв інформації системи поза контрольованими зонами}


{\footnotesize\bfseries
No: 9\\
Name: mp\_5\_odp\_03\\
Type: string\\
Default: nil
}

{\footnotesize Визначено заходи безпеки, що використовуються для контролю носіїв інформації системи за межами контрольованих зон}




\subsection{ЗНИЩЕННЯ ІНФОРМАЦІЇ НА НОСІЯХ ІНФОРМАЦІЇ (MP-6)}
a. Очищувати [Призначення: визначені організацією системні носії] перед утилізацією, випуском за межі організаційного контролю, або перед повторним використанням [Призначення: методами та процедурами очищення, визначеними організацією].\\ 
b. Використовувати механізми очищення зі стійкістю та цілісністю, що відповідає категорії безпеки або рівню секретності інформації.

\vspace{10pt}
{\footnotesize\bfseries
No: 1\\
Name: mp\_6\_a\_01\\
Type: string\\
Default: nil
}

{\footnotesize Носії інформації системи перед утилізацією піддаються очищенню за допомогою методів та процедур очищення}


{\footnotesize\bfseries
No: 2\\
Name: mp\_6\_a\_02\\
Type: string\\
Default: nil
}

{\footnotesize Носії інформації системи очищаються за допомогою методів та процедур очищення перед випуском за межі контрольованої зони}


{\footnotesize\bfseries
No: 3\\
Name: mp\_6\_a\_03\\
Type: string\\
Default: nil
}

{\footnotesize Носії інформації системи очищуються за допомогою методів і процедур очищення перед повторним використанням}


{\footnotesize\bfseries
No: 4\\
Name: mp\_6\_b\\
Type: string\\
Default: \textquotedbl{}автоматизований засіб моніторингу\textquotedbl{}
}

{\footnotesize Застосовуються механізми очищення, надійність і цілісність яких відповідає категорії безпеки або рівню секретності інформації}


{\footnotesize\bfseries
No: 5\\
Name: mp\_6\_odp\_01\\
Type: string\\
Default: nil
}

{\footnotesize Визначено носії інформації системи, які підлягають очищенню перед утилізацією}


{\footnotesize\bfseries
No: 6\\
Name: mp\_6\_odp\_02\\
Type: string\\
Default: nil
}

{\footnotesize Визначено носії інформації системи, які підлягають очищенню перед випуском за межі контрольованої зони}


{\footnotesize\bfseries
No: 7\\
Name: mp\_6\_odp\_03\\
Type: string\\
Default: nil
}

{\footnotesize Визначені носії інформації системи, що підлягають очищенню перед повторним використанням}


{\footnotesize\bfseries
No: 8\\
Name: mp\_6\_odp\_04\\
Type: string\\
Default: nil
}

{\footnotesize Визначено методи та процедури очищення, які слід використовувати для очищення перед утилізацією}


{\footnotesize\bfseries
No: 9\\
Name: mp\_6\_odp\_05\\
Type: string\\
Default: nil
}

{\footnotesize Визначено методи та процедури очищення, які слід використовувати для очищення перед випуском за межі контрольованої зони}


{\footnotesize\bfseries
No: 10\\
Name: mp\_6\_odp\_06\\
Type: string\\
Default: nil
}

{\footnotesize Визначено методи та процедури очищення, які слід використовувати для очищення перед повторним використанням}




\subsection{ВИКОРИСТАННЯ НОСІЇВ ІНФОРМАЦІЇ (MP-7)}
a. [Вибір: обмежити; заборонити] використання [Призначення: визначених організацією типів носіїв системи] на [Призначення: визначені організацією системи або компоненти системи], використовуючи [Призначення: визначені організацією заходи безпеки].\\ 
b. Заборонити використання портативних пристроїв зберігання даних в системах організації, якщо такі пристрої не мають визначеного власника.

\vspace{10pt}
{\footnotesize\bfseries
No: 1\\
Name: mp\_7\_01\\
Type: string\\
Default: nil
}

{\footnotesize ВИКОРИСТАННЯ НОСІЇВ ІНФОРМАЦІЇ МЕТА ОЦІНКИ: Визначити, чи:}


{\footnotesize\bfseries
No: 2\\
Name: mp\_7\_b\\
Type: string\\
Default: nil
}

{\footnotesize Використання зовнішніх носіїв інформації в системах організації заборонено, якщо такі пристрої не мають власника, якого можна ідентифікувати}


{\footnotesize\bfseries
No: 3\\
Name: mp\_7\_odp\_01\\
Type: string\\
Default: nil
}

{\footnotesize Визначено типи носіїв інформації, які мають бути обмежені або заборонені до використання в системі або компонентах системи}


{\footnotesize\bfseries
No: 4\\
Name: mp\_7\_odp\_02\\
Type: string\\
Default: nil
}

{\footnotesize Вибрано одне з наступних ЗНАЧЕНЬ ПАРАМЕТРА: \{обмежити; заборонити\}}


{\footnotesize\bfseries
No: 5\\
Name: mp\_7\_odp\_03\\
Type: string\\
Default: nil
}

{\footnotesize Системи або компоненти системи, для яких визначено використання певних типів носіїв інформації, що підлягають обмеженню або забороні}


{\footnotesize\bfseries
No: 6\\
Name: mp\_7\_odp\_04\\
Type: string\\
Default: nil
}

{\footnotesize Визначено заходи безпеки для обмеження або заборони використання певних типів носіїв інформації в системах або компонентах системи}




\section{PE}
Клас заходів захисту PE --- ФІЗИЧНИЙ ЗАХИСТ І ЗАХИСТ РОБОЧОГО

Цей клас охоплює заходи контролю фізичного доступу до об'єктів організації та захисту обладнання від загроз навколишнього середовища.

\textbf{Перелік заходів захисту:} Ре-1 фізичного захисту та захисту робочого середовища авторизація фізичного (PE-1); Ре-2 доступу (PE-2); ФІЗИЧНИЙ ДОСТУП ДО СИСТЕМИ; Ре-4 ліній електроживлення (PE-4); КОНТРОЛЬ ДОСТУПУ В ПРИМІЩЕННЯ ДЛЯ ВІДОБРАЖЕННЯ ІНФОРМАЦІЇ; МОНІТОРИНГ ФІЗИЧНОГО ДОСТУПУ; Альтернативне робоче місце (PE-17).

\subsection{РЕ-1 фізичного захисту та захисту робочого середовища Авторизація фізичного (PE-1)}


\vspace{10pt}
{\footnotesize\bfseries
No: 1\\
Name: pe\_1\_a\_01\\
Type: list\\
Default: [\textquotedbl{}default\_deny\_rule\textquotedbl{}, \textquotedbl{}abac\_rule\_1\textquotedbl{}]
}

{\footnotesize Розроблено та задокументовано політику фізичного захисту та захисту робочого середовища}


{\footnotesize\bfseries
No: 2\\
Name: pe\_1\_a\_02\\
Type: list\\
Default: [\textquotedbl{}admin\textquotedbl{}, \textquotedbl{}security\_officer\textquotedbl{}]
}

{\footnotesize Політика фізичного захисту та захисту робочого середовища розповсюджується серед персоналу або ролей}


{\footnotesize\bfseries
No: 3\\
Name: pe\_1\_b\\
Type: list\\
Default: [\textquotedbl{}admin\textquotedbl{}, \textquotedbl{}security\_officer\textquotedbl{}]
}

{\footnotesize Посадова особа призначається для управління розробкою, документуванням та розповсюдженням політики та процедур фізичного захисту та захисту робочого середовища; PE-01(c)[01][01] переглядається та оновлюється поточна політика фізичного захисту та захисту робочого середовища частота; PE-01(c)[01][02] поточна політика фізичного захисту та захисту робочого середовища переглядається та оновлюється після подій; PE-01(c)[02][01] переглядаються та оновлюються поточні процедури фізичного захисту та захисту робочого середовища частота; поточні процедури фізичного та екологічного захисту переглядаються та оновлюються після подій. PE-01(c)[02][02]}


{\footnotesize\bfseries
No: 4\\
Name: pe\_1\_odp\_01\\
Type: list\\
Default: [\textquotedbl{}admin\textquotedbl{}, \textquotedbl{}security\_officer\textquotedbl{}]
}

{\footnotesize Визначено персонал або ролі, до яких має бути доведена політика фізичного захисту та захисту робочого середовища}


{\footnotesize\bfseries
No: 5\\
Name: pe\_1\_odp\_02\\
Type: list\\
Default: [\textquotedbl{}admin\textquotedbl{}, \textquotedbl{}security\_officer\textquotedbl{}]
}

{\footnotesize Визначено персонал або ролі, на які поширюються процедури фізичного захисту та захисту робочого середовища}


{\footnotesize\bfseries
No: 6\\
Name: pe\_1\_odp\_03\\
Type: string\\
Default: nil
}

{\footnotesize Вибрано одне або декілька з наступних ЗНАЧЕНЬ ПАРАМЕТРІВ: \{рівень організації; рівень місії/бізнес-процесу; рівень системи\}}


{\footnotesize\bfseries
No: 7\\
Name: pe\_1\_odp\_04\\
Type: list\\
Default: [\textquotedbl{}admin\textquotedbl{}, \textquotedbl{}security\_officer\textquotedbl{}]
}

{\footnotesize Визначено посадову особу, яка керуватиме політикою та процедурами фізичного захисту та захисту робочого середовища}


{\footnotesize\bfseries
No: 8\\
Name: pe\_1\_odp\_05\\
Type: list\\
Default: [\textquotedbl{}default\_deny\_rule\textquotedbl{}, \textquotedbl{}abac\_rule\_1\textquotedbl{}]
}

{\footnotesize Визначено частоту, з якою переглядається та оновлюється поточна політика фізичного захисту та захисту робочого середовища}


{\footnotesize\bfseries
No: 9\\
Name: pe\_1\_odp\_06\\
Type: list\\
Default: [\textquotedbl{}default\_deny\_rule\textquotedbl{}, \textquotedbl{}abac\_rule\_1\textquotedbl{}]
}

{\footnotesize Визначено події, які потребують перегляду та оновлення поточної політики фізичного захисту та захисту робочого середовища}


{\footnotesize\bfseries
No: 10\\
Name: pe\_1\_odp\_07\\
Type: integer\\
Default: 30
}

{\footnotesize Визначено частоту, з якою переглядаються та оновлюються поточні процедури фізичного захисту та захисту робочого середовища}


{\footnotesize\bfseries
No: 11\\
Name: pe\_1\_odp\_08\\
Type: list\\
Default: [\textquotedbl{}login\textquotedbl{}, \textquotedbl{}logout\textquotedbl{}, \textquotedbl{}failed\_attempt\textquotedbl{}]
}

{\footnotesize Визначено події, які потребують перегляду та оновлення процедур фізичного захисту та захисту робочого середовища}




\subsection{РЕ-2 доступу (PE-2)}


\vspace{10pt}
{\footnotesize\bfseries
No: 1\\
Name: pe\_2\_a\_01\\
Type: list\\
Default: []
}

{\footnotesize Розроблено перелік осіб, які мають право авторизованого доступу до об'єкта, де перебуває система}


{\footnotesize\bfseries
No: 2\\
Name: pe\_2\_a\_02\\
Type: list\\
Default: []
}

{\footnotesize Затверджено перелік осіб, які мають право авторизованого доступу до об'єкта, де перебуває система}


{\footnotesize\bfseries
No: 3\\
Name: pe\_2\_a\_03\\
Type: list\\
Default: []
}

{\footnotesize Ведеться перелік осіб, які мають право авторизованого доступу до об'єкта, де перебуває система}


{\footnotesize\bfseries
No: 4\\
Name: pe\_2\_b\\
Type: string\\
Default: nil
}

{\footnotesize Для доступу до об'єкта надаються повноваження}


{\footnotesize\bfseries
No: 5\\
Name: pe\_2\_c\\
Type: list\\
Default: [\textquotedbl{}admin\textquotedbl{}, \textquotedbl{}security\_officer\textquotedbl{}]
}

{\footnotesize Переглядається список доступу, , у якому закріплений перелік персоналу або ролей, яким дозволений санкціонований доступ до об'єкта частота}


{\footnotesize\bfseries
No: 6\\
Name: pe\_2\_d\\
Type: list\\
Default: [\textquotedbl{}admin\textquotedbl{}, \textquotedbl{}security\_officer\textquotedbl{}]
}

{\footnotesize Особи видаляються зі списку доступу до об'єкта, коли доступ більше не потрібен}


{\footnotesize\bfseries
No: 7\\
Name: pe\_2\_odp\\
Type: list\\
Default: [\textquotedbl{}admin\textquotedbl{}, \textquotedbl{}security\_officer\textquotedbl{}]
}

{\footnotesize Визначено періодичність перегляду списку доступу, у якому закріплений перелік персоналу або ролей, яким дозволений санкціонований доступ до об'єкта}




\subsection{ФІЗИЧНИЙ ДОСТУП ДО СИСТЕМИ}


\vspace{10pt}
{\footnotesize\bfseries
No: 1\\
Name: pe\_3\_a\_01\\
Type: string\\
Default: nil
}

{\footnotesize Авторизація фізичного доступу забезпечується в пунктах входу і виходу шляхом перевірки індивідуальних дозволів доступу}


{\footnotesize\bfseries
No: 2\\
Name: pe\_3\_b\\
Type: string\\
Default: nil
}

{\footnotesize Журнали контролю фізичного доступу ведуться для 03\_ODP[04] точок входу або виходу>}


{\footnotesize\bfseries
No: 3\\
Name: pe\_3\_c\\
Type: string\\
Default: nil
}

{\footnotesize Доступ в зони всередині об'єкту, визначені як загальнодоступні, підтримується шляхом впровадження заходів захисту}


{\footnotesize\bfseries
No: 4\\
Name: pe\_3\_d\_01\\
Type: string\\
Default: nil
}

{\footnotesize Відвідувачів супроводжують}


{\footnotesize\bfseries
No: 5\\
Name: pe\_3\_d\_02\\
Type: list\\
Default: []
}

{\footnotesize Активність відвідувачів контролюється умови}


{\footnotesize\bfseries
No: 6\\
Name: pe\_3\_e\_01\\
Type: string\\
Default: nil
}

{\footnotesize Ключі захищені}


{\footnotesize\bfseries
No: 7\\
Name: pe\_3\_e\_02\\
Type: string\\
Default: nil
}

{\footnotesize Коди доступу захищені}


{\footnotesize\bfseries
No: 8\\
Name: pe\_3\_e\_03\\
Type: string\\
Default: nil
}

{\footnotesize Інші пристрої фізичного доступу захищені}


{\footnotesize\bfseries
No: 9\\
Name: pe\_3\_f\\
Type: string\\
Default: \textquotedbl{}щорічно\textquotedbl{}
}

{\footnotesize Пристрої фізичного доступу інвентаризуються частота}


{\footnotesize\bfseries
No: 10\\
Name: pe\_3\_g\_01\\
Type: list\\
Default: [\textquotedbl{}admin\textquotedbl{}, \textquotedbl{}security\_officer\textquotedbl{}]
}

{\footnotesize Коди доступу змінюється частота, коли код скомпрометовано, або коли особи, які володіють кодом, переводяться або звільняються; <PE- PE-03(g)[02] ключі змінюються частота, коли ключі втрачено, або коли особи, що володіють ключами, переводяться або звільняються}


{\footnotesize\bfseries
No: 11\\
Name: pe\_3\_odp\_01\\
Type: string\\
Default: nil
}

{\footnotesize Визначено точки входу та виходу в об'єкт, в якому знаходиться система}


{\footnotesize\bfseries
No: 12\\
Name: pe\_3\_odp\_02\\
Type: string\\
Default: nil
}

{\footnotesize Вибрано одне або декілька з наступних ЗНАЧЕНЬ ПАРАМЕТРІВ: \{системи або пристрої; охоронці\}}


{\footnotesize\bfseries
No: 13\\
Name: pe\_3\_odp\_03\\
Type: string\\
Default: nil
}

{\footnotesize Визначено фізичні системи або пристрої контролю доступу, що використовуються для контролю входу та виходу на об'єкт (якщо вибрано); PE-03\_ODP[04] визначено точки входу або виходу, для яких ведуться журнали контролю фізичного доступу}


{\footnotesize\bfseries
No: 14\\
Name: pe\_3\_odp\_05\\
Type: string\\
Default: nil
}

{\footnotesize Визначено заходи захисту для контролю доступу в зони всередині об'єкту, позначені як загальнодоступні}


{\footnotesize\bfseries
No: 15\\
Name: pe\_3\_odp\_06\\
Type: list\\
Default: []
}

{\footnotesize Визначено умови, що вимагають супроводу відвідувачів та моніторингу активності відвідувачів}


{\footnotesize\bfseries
No: 16\\
Name: pe\_3\_odp\_07\\
Type: string\\
Default: nil
}

{\footnotesize Визначені пристрої фізичного доступу, що підлягають інвентаризації}


{\footnotesize\bfseries
No: 17\\
Name: pe\_3\_odp\_08\\
Type: integer\\
Default: 30
}

{\footnotesize Визначено частоту проведення інвентаризації пристроїв фізичного доступу}


{\footnotesize\bfseries
No: 18\\
Name: pe\_3\_odp\_09\\
Type: integer\\
Default: 30
}

{\footnotesize Визначено частоту, з якою потрібно змінювати коди доступу}


{\footnotesize\bfseries
No: 19\\
Name: pe\_3\_odp\_10\\
Type: integer\\
Default: 30
}

{\footnotesize Визначено частоту, з якою потрібно змінювати ключі}




\subsection{РЕ-4 ліній електроживлення (PE-4)}


\vspace{10pt}
{\footnotesize\bfseries
No: 1\\
Name: pe\_4\_01\\
Type: string\\
Default: nil
}

{\footnotesize Фізичний доступ до систем розподілу та постачання живлення в межах об'єктів організації контролюється за допомогою заходів захисту}


{\footnotesize\bfseries
No: 2\\
Name: pe\_4\_odp\_01\\
Type: string\\
Default: nil
}

{\footnotesize Визначені системи розподілу та постачання живлення, які потребують фізичного контролю доступу}


{\footnotesize\bfseries
No: 3\\
Name: pe\_4\_odp\_02\\
Type: string\\
Default: nil
}

{\footnotesize Визначено заходи захисту, які необхідно впровадити для контролю фізичного доступу до систем розподілу та постачання живлення в межах об'єкту організації}




\subsection{КОНТРОЛЬ ДОСТУПУ В ПРИМІЩЕННЯ ДЛЯ ВІДОБРАЖЕННЯ ІНФОРМАЦІЇ}


\vspace{10pt}
{\footnotesize\bfseries
No: 1\\
Name: pe\_5\_odp\\
Type: string\\
Default: nil
}

{\footnotesize Визначено пристрої для виведення інформації над якими необхідний контроль над фізичним доступом до вихідних даних}




\subsection{МОНІТОРИНГ ФІЗИЧНОГО ДОСТУПУ}


\vspace{10pt}
{\footnotesize\bfseries
No: 1\\
Name: pe\_6\_a\\
Type: string\\
Default: nil
}

{\footnotesize Фізичний доступ до об'єкту, де знаходиться система, моніториться з метою виявлення та реагування на інциденти фізичної безпеки}


{\footnotesize\bfseries
No: 2\\
Name: pe\_6\_b\_01\\
Type: string\\
Default: \textquotedbl{}щорічно\textquotedbl{}
}

{\footnotesize Переглядаються журнали фізичного доступу частота}


{\footnotesize\bfseries
No: 3\\
Name: pe\_6\_b\_02\\
Type: string\\
Default: nil
}

{\footnotesize Журнали фізичного доступу переглядаються при виникненні подій}


{\footnotesize\bfseries
No: 4\\
Name: pe\_6\_c\_01\\
Type: string\\
Default: nil
}

{\footnotesize Результати переглядів узгоджуються з можливостями організації щодо реагування на інциденти}


{\footnotesize\bfseries
No: 5\\
Name: pe\_6\_c\_02\\
Type: string\\
Default: nil
}

{\footnotesize Результати розслідувань узгоджуються з можливостями організації щодо реагування на інциденти}


{\footnotesize\bfseries
No: 6\\
Name: pe\_6\_odp\_01\\
Type: integer\\
Default: 30
}

{\footnotesize Визначено частоту перегляду журналів фізичного доступу}


{\footnotesize\bfseries
No: 7\\
Name: pe\_6\_odp\_02\\
Type: list\\
Default: [\textquotedbl{}login\textquotedbl{}, \textquotedbl{}logout\textquotedbl{}, \textquotedbl{}failed\_attempt\textquotedbl{}]
}

{\footnotesize Визначено події або потенційні ознаки подій, що вимагають перегляду журналів фізичного доступу}




\subsection{АЛЬТЕРНАТИВНЕ РОБОЧЕ МІСЦЕ (PE-17)}


\vspace{10pt}
{\footnotesize\bfseries
No: 1\\
Name: pe\_17\_odp\_01\\
Type: string\\
Default: nil
}

{\footnotesize Визначені альтернативні робочі місця, дозволені для використання працівниками}


{\footnotesize\bfseries
No: 2\\
Name: pe\_17\_odp\_02\\
Type: list\\
Default: [\textquotedbl{}admin\textquotedbl{}, \textquotedbl{}security\_officer\textquotedbl{}]
}

{\footnotesize Визначаються заходи захисту, які будуть застосовуватися на альтернативних робочих місцях; РЕ-17(a) альтернативні робочі місця визначені та задокументовані; РЕ-17(b) заходи захисту впроваджені на альтернативних робочих місцях; РЕ-17(c) оцінюється ефективність заходів захисту на альтернативних робочих місцях; РЕ-17(d) працівникам надаються засоби комунікації з персоналом служби інформаційної безпеки на випадок інцидентів}




\section{PL}
Клас заходів захисту PL --- ПЛАНУВАННЯ БЕЗПЕКИ

Цей клас регламентує розробку, документування та регулярне оновлення планів захисту інформації, архітектури безпеки та приватності.

\textbf{Перелік заходів захисту:} Політики та процедури планування безпеки (PL-1); Плани захисту інформації та персональних даних (PL-2); Правила поведінки (PL-4).

\subsection{ПОЛІТИКИ ТА ПРОЦЕДУРИ ПЛАНУВАННЯ БЕЗПЕКИ (PL-1)}
a. Розробити, задокументувати та поширити серед [Призначення: визначеного організацією персоналу або ролей]:\\ 
1. 2. [вибір (один або декілька): рівень організації; рівень місії/бізнес процесу; рівень системи] політику планування, яка:\\ 
(a) містить мету, сферу застосування, ролі, обов'язки, відповідальність керівництва, координацію між організаційними підрозділами та системою контролю (complaince);\\ 
(b) відповідає чинному законодавству, виконавчим наказам, директивам, нормам, політикам, стандартам та керівним принципам; процедури, що полегшують здійснення планування політики безпеки та приватності й пов'язані з ними заходи.\\ 
b. Призначити [Призначення: визначену організацією посадову особу] для управління політикою та процедурами планування політики безпеки та приватності.\\ 
c. Переглядати та оновлювати поточне планування:\\ 
1. політики планування безпеки та приватності [Призначення:з визначеною організацією частотою] та наступні [Призначення: визначені організацією події];\\ 
2. поточні процедури планування політики [Призначення: з визначеною організацією [Призначення: визначені організацією події]. безпеки та приватності частотою] та наступні

\vspace{10pt}
{\footnotesize\bfseries
No: 1\\
Name: pl\_1\_odp\_01\\
Type: list\\
Default: [\textquotedbl{}admin\textquotedbl{}, \textquotedbl{}security\_officer\textquotedbl{}]
}

{\footnotesize Визначено персонал або ролі, до яких має бути доведена політика планування безпеки}


{\footnotesize\bfseries
No: 2\\
Name: pl\_1\_odp\_02\\
Type: list\\
Default: [\textquotedbl{}admin\textquotedbl{}, \textquotedbl{}security\_officer\textquotedbl{}]
}

{\footnotesize Визначено персонал або ролі, на які поширюватимуться процедури планування безпеки}


{\footnotesize\bfseries
No: 3\\
Name: pl\_1\_odp\_03\\
Type: string\\
Default: nil
}

{\footnotesize Вибрано одне або декілька з наступних ЗНАЧЕНЬ ПАРАМЕТРІВ: \{рівень організації; рівень місії/бізнеспроцесу; рівень системи\}}


{\footnotesize\bfseries
No: 4\\
Name: pl\_1\_odp\_04\\
Type: list\\
Default: [\textquotedbl{}admin\textquotedbl{}, \textquotedbl{}security\_officer\textquotedbl{}]
}

{\footnotesize Визначено посадову особу, яка керуватиме політикою та процедурами планування безпеки}


{\footnotesize\bfseries
No: 5\\
Name: pl\_1\_odp\_05\\
Type: list\\
Default: [\textquotedbl{}default\_deny\_rule\textquotedbl{}, \textquotedbl{}abac\_rule\_1\textquotedbl{}]
}

{\footnotesize Визначено періодичність перегляду та оновлення поточної політики планування безпеки}


{\footnotesize\bfseries
No: 6\\
Name: pl\_1\_odp\_06\\
Type: list\\
Default: [\textquotedbl{}default\_deny\_rule\textquotedbl{}, \textquotedbl{}abac\_rule\_1\textquotedbl{}]
}

{\footnotesize Є події, які потребують перегляду та оновлення поточної політики планування безпеки}


{\footnotesize\bfseries
No: 7\\
Name: pl\_1\_odp\_07\\
Type: string\\
Default: \textquotedbl{}щорічно\textquotedbl{}
}

{\footnotesize Визначена частота, з якою переглядаються та оновлюються поточні процедури планування безпеки}




\subsection{ПЛАНИ ЗАХИСТУ ІНФОРМАЦІЇ ТА ПЕРСОНАЛЬНИХ ДАНИХ (PL-2)}
a. Розробити план захисту інформації та персональних даних для інформаційної системи, який:\\ 
1. узгоджується з архітектурою підприємства організації;\\ 
2. чітко визначає складові компоненти системи;\\ 
3. описує оперативний контекст інформаційної системи з точки зору завдань та процесів;\\ 
4. визначає осіб, які виконують системні ролі та обов'язки;\\ 
5. визначає тип інформації, яка обробляється, зберігається та передається системою;\\ 
6. надає огляд вимог безпеки та приватності інформаційної системи;\\ 
7. описує будь які конкретні загрози системи, які викликають стурбованість організації;\\ 
8. надає результати оцінки ризику конфіденційності для систем, в яких обробляються персональні дані;\\ 
9. описує робоче середовище інформаційної системи та будь які залежності від систем або компонентів систем або підключень до таких систем та їх компонентів;\\ 
10. надає огляд вимог безпеки та конфіденційності системи;\\ 
11. визначає будь які відповідні контрольні базові рівні або накладання, якщо вони застосовуються;\\ 
12. описує чинні або заплановані заходи щодо забезпечення безпеки та приватності, включно з обґрунтуванням рішень щодо налаштування\\ 
13. включає виявлення ризиків для архітектури безпеки і приватності, а також проєктних рішень;\\ 
14. включає дії, пов'язані з безпекою та конфіденційністю, які впливають на систему, виконання яких вимагає планування та координацію з [Призначення: визначені організацією окремі особи або групи];\\ 
15. розглядається та затверджується уповноваженою посадовою особою або призначеним представником до початку реалізації плану.\\ 
b. Поширити копії планів захисту інформації та персональних даних і повідомляти про подальші зміни планів серед [Призначення:визначеного організацією персоналу або ролей].\\ 
c. Переглядати плани захисту інформації та персональних даних [Призначення: з визначеною організацією частотою].\\ 
d. Оновлювати плани захисту інформації та персональних даних для врахування змін в інформаційній системі й робочому середовищі або проблем, виявлених у ході реалізації або оцінювання заходів безпеки та приватності.\\ 
e. забезпечити захист планів захисту інформації та персональних даних від несанкціонованого розголошення та змін.

\vspace{10pt}
{\footnotesize\bfseries
No: 1\\
Name: pl\_2\_odp\_01\\
Type: list\\
Default: [\textquotedbl{}admin\textquotedbl{}, \textquotedbl{}security\_officer\textquotedbl{}]
}

{\footnotesize Призначені особи або групи, з якими пов'язана діяльність з безпекою та конфіденційністю, що впливає на систему, яка потребує планування та координації}


{\footnotesize\bfseries
No: 2\\
Name: pl\_2\_odp\_02\\
Type: list\\
Default: [\textquotedbl{}admin\textquotedbl{}, \textquotedbl{}security\_officer\textquotedbl{}]
}

{\footnotesize Призначено персонал або ролі для отримання розповсюджуваних копій планів захисту інформації та конфіденційності системи}


{\footnotesize\bfseries
No: 3\\
Name: pl\_2\_odp\_03\\
Type: list\\
Default: [\textquotedbl{}admin\textquotedbl{}, \textquotedbl{}security\_officer\textquotedbl{}]
}

{\footnotesize Визначено періодичність перегляду інформації та конфіденційності системи; планів захисту PL-02a.01[01] розроблено план захисту інформації, архітектурі підприємства організації; який відповідає PL-02a.01[02] розроблено план конфіденційності для системи, відповідає архітектурі підприємства організації; який PL-02a.02[01] розроблено план захисту інформації, який чітко визначає складові компоненти системи; PL-02a.02[02] розроблено для системи план забезпечення конфіденційності, який чітко визначає складові компоненти системи; PL-02a.03[01] розроблено план захисту інформації системи, який описує операційний контекст системи з точки зору місії та бізнеспроцесів; PL-02a.03[02] розроблено для системи план забезпечення конфіденційності, який описує операційний контекст системи з точки зору місії та бізнес-процесів; PL-02a.04[01] розроблено план захисту інформації, який визначає осіб, що виконують системні ролі та обов'язки; PL-02a.04[02] розроблено для системи план забезпечення конфіденційності, який визначає осіб, що виконують ролі та обов'язки в системі; PL-02a.05[01] розроблено план захисту інформації, який визначає типи інформації, що обробляється, зберігається та передається системою; PL-02a.05[02] розроблено план конфіденційності для системи, який визначає типи інформації, що обробляється, зберігається та передається системою; PL-02a.06[01] розроблено план захисту інформації, який передбачає категоризацію безпеки системи, включаючи відповідне обґрунтування; PL-02a.06[02] розроблено для системи план забезпечення конфіденційності, який передбачає категоризацію системи за рівнем безпеки, включаючи обґрунтування; PL-02a.07[01] розроблено план захисту інформації, який описує будь-які конкретні загрози для системи, що викликають потенційні ризики для рганізації; PL-02a.07[02] розроблено план конфіденційності для системи, який описує будь-які конкретні загрози для системи, що викликають потенційні ризики для організації; PL-02a.08[01] розроблено план захисту інформації, який містить результати оцінки ризиків конфіденційності для систем, що обробляють інформацію, яка ідентифікує особу; PL-02a.08[02] розроблено для системи план забезпечення конфіденційності, який містить результати оцінки ризиків конфіденційності для систем, що обробляють інформацію, яка ідентифікує особу; PL-02a.09[01] розроблено план захисту інформації, який описує операційне середовище системи та будь-які залежності або зв'язки з іншими системами чи компонентами системи; PL-02a.09[02] розроблено для системи план забезпечення конфіденційності, який описує робоче середовище системи та будь-які залежності або зв'язки з іншими системами чи компонентами системи; PL-02a.10[01] розроблено план захисту інформації, який містить огляд вимог до безпеки системи; PL-02a.10[02] розроблено для системи план забезпечення конфіденційності, який містить огляд вимог до конфіденційності системи; PL-02a.11[01] розроблено план захисту інформації, який визначає будь-які відповідні базові рівні контролю або обмеження, якщо такі є; PL-02a.11[02] розроблено для системи план забезпечення конфіденційності, який визначає будь-які відповідні базові рівні контролю або обмеження, якщо такі є; PL-02a.12[01] розроблено план захисту інформації, який описує наявні або заплановані засоби контролю для виконання вимог безпеки, включаючи обґрунтування будь-яких рішень щодо адаптації; PL-02a.12[02] розроблено для системи план забезпечення конфіденційності, який описує наявні або заплановані засоби контролю для виконання вимог щодо конфіденційності, включаючи обґрунтування будь-яких рішень, пов'язаних з адаптацією; PL-02a.13[01] розроблено план захисту інформації, який включає визначення ризиків для архітектури безпеки та проектних рішень; PL-02a.13[02] розроблено план забезпечення конфіденційності для системи, який включає визначення ризиків для архітектури конфіденційності та проектних рішень; PL-02a.14[01] розроблено захисту інформації, який включає діяльність, пов'язану з безпекою, що впливає на систему і потребує планування та координації з окремими особами або групами; PL-02a.14[02] розроблено для системи план забезпечення конфіденційності, який включає діяльність, пов'язану з конфіденційністю, що впливає на систему і потребує планування та координації з окремими особами або групами; PL-02a.15[01] розроблено план захисту інформації, який розглядається та затверджується уповноваженою посадовою особою або призначеним представником до початку реалізації плану; PL-02a.15[02] розроблено план забезпечення конфіденційності для системи, який перевіряється та затверджується уповноваженою посадовою особою або призначеним представником перед впровадженням плану. PL-02b.[01] розповсюджуються копії персонал або ролі>; PL-02b.[02] повідомляються наступні зміни до планів персонал або ролі; PL-02c. переглядаються плани частота; планів серед <PL-02\_ODP[02] PL-02d.[01] оновлюються плани відповідно до змін у системі та середовищі діяльності; PL-02d.[02] оновлюються плани для вирішення проблем, виявлених під час реалізації плану; PL-02d.[03] оновлюються плани для вирішення проблем, виявлених під час контрольних оцінок; PL-02e.[01] захищені плани від несанкціонованого розголошення; PL-02e.[02] захищені плани від несанкціонованої модифікації}




\subsection{ПРАВИЛА ПОВЕДІНКИ (PL-4)}
a. Створити та надати особам, що потребують доступу до інформаційної системи, правила, які описують їхні обов'язки й очікувану поведінку щодо інформації та використання інформаційної системи, безпеки та приватності.\\ 
b. Отримати документальне підтвердження від таких осіб про те, що вони прочитали, зрозуміли та погодилися дотримуватися правил поведінки, перш ніж дозволяти доступ до інформації та інформаційної системи.\\ 
c. Переглядати й оновлювати правила поведінки [Призначення: з визначеною організацією частотою].\\ 
d. Вимагати від осіб, які підписали попередню версію правил поведінки, перечитати та повторно підписати правила [Вибір (один або декілька): [Призначення: з визначеною організацією частотою]; коли правила переглядаються чи оновлюються].

\vspace{10pt}
{\footnotesize\bfseries
No: 1\\
Name: pl\_4\_odp\_01\\
Type: string\\
Default: \textquotedbl{}щорічно\textquotedbl{}
}

{\footnotesize Визначено періодичність перегляду та оновлення правил поведінки}


{\footnotesize\bfseries
No: 2\\
Name: pl\_4\_odp\_02\\
Type: list\\
Default: [\textquotedbl{}default\_deny\_rule\textquotedbl{}, \textquotedbl{}abac\_rule\_1\textquotedbl{}]
}

{\footnotesize Вибрано одне або декілька з наступних ЗНАЧЕНЬ ПАРАМЕТРА:\{частота; коли правила переглядаються або оновлюються\}}


{\footnotesize\bfseries
No: 3\\
Name: pl\_4\_odp\_03\\
Type: list\\
Default: [\textquotedbl{}admin\textquotedbl{}, \textquotedbl{}security\_officer\textquotedbl{}]
}

{\footnotesize Визначена періодичність перегляду та повторного підтвердження правил поведінки (якщо вибрано); PL-04a.[01] встановлені правила, які описують обов'язки та очікувану поведінку щодо використання інформації та системи, безпеки та конфіденційності для осіб, яким потрібен доступ до системи; PL-04a.[02] надаються правила, які описують обов'язки та очікувану поведінку щодо використання інформації та системи, безпеки та конфіденційності, особам, які отримують доступ до системи; PL-04b. отримано перед наданням доступу до інформації та системи задокументоване підтвердження від таких осіб про те, що вони прочитали, зрозуміли та згодні дотримуватися правил поведінки; PL-04c. переглядаються та оновлюються правила поведінки частота; PL-04d. потрібно особам, які визнали попередню версію правил поведінки, прочитати та повторно визнати ЗНАЧЕННЯ ВИБРАНОГО ПАРАМЕТРА(ів)}




\section{PS}
Клас заходів захисту PS --- КАДРОВА БЕЗПЕКА

Цей клас встановлює вимоги до перевірки, надання повноважень та звільнення персоналу з метою мінімізації ризиків інсайдерських загроз.

\textbf{Перелік заходів захисту:} Політика та процедури кадрової безпеки (PS-1); Перевірка персоналу (PS-3); Звільнення персоналу (PS-4); Переведення персоналу (PS-5).

\subsection{Політика та процедури кадрової безпеки (PS-1)}


\vspace{10pt}
{\footnotesize\bfseries
No: 1\\
Name: ps\_1\_odp\_01\\
Type: list\\
Default: [\textquotedbl{}admin\textquotedbl{}, \textquotedbl{}security\_officer\textquotedbl{}]
}

{\footnotesize Визначено персонал або ролі, на які поширюється політика кадрової безпеки}


{\footnotesize\bfseries
No: 2\\
Name: ps\_1\_odp\_02\\
Type: list\\
Default: [\textquotedbl{}admin\textquotedbl{}, \textquotedbl{}security\_officer\textquotedbl{}]
}

{\footnotesize Визначено персонал або ролі, на які поширюються процедури кадрової безпеки}


{\footnotesize\bfseries
No: 3\\
Name: ps\_1\_odp\_03\\
Type: string\\
Default: nil
}

{\footnotesize Вибрано одне або декілька з наступних ЗНАЧЕНЬ ПАРАМЕТРІВ: \{рівень організації; рівень місії/бізнеспроцесу; рівень системи\}}


{\footnotesize\bfseries
No: 4\\
Name: ps\_1\_odp\_04\\
Type: list\\
Default: [\textquotedbl{}admin\textquotedbl{}, \textquotedbl{}security\_officer\textquotedbl{}]
}

{\footnotesize Визначено посадову особу, яка керуватиме політикою та процедурами кадрової безпеки}


{\footnotesize\bfseries
No: 5\\
Name: ps\_1\_odp\_05\\
Type: list\\
Default: [\textquotedbl{}default\_deny\_rule\textquotedbl{}, \textquotedbl{}abac\_rule\_1\textquotedbl{}]
}

{\footnotesize Визначена періодичність перегляду та оновлення поточної політики кадрової безпеки}


{\footnotesize\bfseries
No: 6\\
Name: ps\_1\_odp\_06\\
Type: list\\
Default: [\textquotedbl{}default\_deny\_rule\textquotedbl{}, \textquotedbl{}abac\_rule\_1\textquotedbl{}]
}

{\footnotesize Є події, які вимагають перегляду та оновлення поточної політики кадрової безпеки}


{\footnotesize\bfseries
No: 7\\
Name: ps\_1\_odp\_07\\
Type: string\\
Default: \textquotedbl{}щорічно\textquotedbl{}
}

{\footnotesize Визначено періодичність перегляду та оновлення поточних процедур кадрової безпеки}




\subsection{Перевірка персоналу (PS-3)}


\vspace{10pt}
{\footnotesize\bfseries
No: 1\\
Name: ps\_3\_odp\_01\\
Type: list\\
Default: []
}

{\footnotesize Визначені умови, що вимагають повторної перевірки осіб}


{\footnotesize\bfseries
No: 2\\
Name: ps\_3\_odp\_02\\
Type: list\\
Default: [\textquotedbl{}admin\textquotedbl{}, \textquotedbl{}security\_officer\textquotedbl{}]
}

{\footnotesize Визначена частота повторної перевірки осіб, для яких це показано; PS-03a. проходять особи перевірку перед тим, як надати їм доступ до системи; PS-03b.[01] проходять особи повторну перевірку відповідно до умови, що вимагають повторної перевірки; PS-03b.[02] проводиться повторна перевірка у випадках, коли це зазначено, частота}




\subsection{Звільнення персоналу (PS-4)}


\vspace{10pt}
{\footnotesize\bfseries
No: 1\\
Name: ps\_4\_odp\_01\\
Type: integer\\
Default: 30
}

{\footnotesize Визначено період часу, протягом якого забороняється доступ до системи}


{\footnotesize\bfseries
No: 2\\
Name: ps\_4\_odp\_02\\
Type: list\\
Default: [\textquotedbl{}admin\textquotedbl{}, \textquotedbl{}security\_officer\textquotedbl{}]
}

{\footnotesize Визначені теми інформаційної безпеки для обговорення під час проведення співбесід; PS-04a. при звільненні працівника доступ до системи вимикається протягом часового періоду; PS-04b. припиняють дію або анулюють будь-які автентифікатори та облікові дані після припинення трудових відносин з окремими особами; PS-04c. проводяться при звільненні окремих працівників співбесіди, які включають обговорення питань інформаційної безпеки; PS-04d. отримується після звільнення особи все майно, пов'язане з безпекою організаційної системи; PS-04e. зберігається доступ до організаційної інформації та систем, які раніше перебували під контролем особи, що звільняється, після припинення нею трудових відносин}




\subsection{Переведення персоналу (PS-5)}


\vspace{10pt}
{\footnotesize\bfseries
No: 1\\
Name: ps\_5\_odp\_01\\
Type: list\\
Default: [\textquotedbl{}login\textquotedbl{}, \textquotedbl{}logout\textquotedbl{}, \textquotedbl{}failed\_attempt\textquotedbl{}]
}

{\footnotesize Визначені дії, які мають бути ініційовані після переведення або перепризначення}


{\footnotesize\bfseries
No: 2\\
Name: ps\_5\_odp\_02\\
Type: list\\
Default: [\textquotedbl{}login\textquotedbl{}, \textquotedbl{}logout\textquotedbl{}, \textquotedbl{}failed\_attempt\textquotedbl{}]
}

{\footnotesize Визначено період часу, протягом якого мають бути здійснені дії з переведення або перепризначення після переведення або перепризначення}


{\footnotesize\bfseries
No: 3\\
Name: ps\_5\_odp\_03\\
Type: list\\
Default: [\textquotedbl{}admin\textquotedbl{}, \textquotedbl{}security\_officer\textquotedbl{}]
}

{\footnotesize Визначено персонал або ролі, про які необхідно повідомляти, коли осіб призначають на інші посади або переводять на інші посади в організації}


{\footnotesize\bfseries
No: 4\\
Name: ps\_5\_odp\_04\\
Type: list\\
Default: [\textquotedbl{}admin\textquotedbl{}, \textquotedbl{}security\_officer\textquotedbl{}]
}

{\footnotesize Визначено період часу, протягом якого необхідно повідомляти визначений організацією персонал або ролі, коли осіб перепризначають або переводять на інші посади в межах організації; PS-05a. переглядаються та підтверджуються поточні потреби в логічних та фізичних дозволах на доступ до систем та об'єктів при перепризначенні або переведенні осіб на інші посади в організації; PS-05b. були ініційовані дії з переведення або перепризначення протягом періоду часу після формальної дії з переведення; PS-05c. змінюється авторизація доступу за необхідності, щоб відповідати будь-яким змінам в оперативних потребах у зв'язку з перепризначенням або переведенням; PS-05d. було повідомлено персонал або ролі протягом часового періоду}




\section{RA}
Клас заходів захисту RA --- ОЦІНЮВАННЯ РИЗИКУ

Цей клас забезпечує систематичний підхід до виявлення, аналізу та реагування на ризики, пов'язані з обробкою інформації та функціонуванням системи.

\textbf{Перелік заходів захисту:} Політика та процедури оцінювання ризику (RA-1); Оцінювання ризику (RA-3); Оцінювання постачання (RA-3(1)); Сканування вразливостей (RA-5); Оновлення за частотою, перед новим скануванням або при ідентифікації (RA-5(2)); Реагування на ризик (RA-7).

\subsection{ПОЛІТИКА ТА ПРОЦЕДУРИ ОЦІНЮВАННЯ РИЗИКУ (RA-1)}
a. Розробити, задокументувати та поширити серед [Призначення: визначеного організацією персоналу або посадових осіб]:\\ 
1. 2. Політику оцінювання ризику, яка:\\ 
(a) містить мету, сферу застосування, ролі, обов'язки, відповідальність керівництва, координацію між організаційними підрозділами та систему контролю відповідності (complaince);\\ 
(b) відповідає чинному законодавству, виконавчим наказам, директивам, нормам, політикам, стандартам і керівним принципам. Процедури, що сприяють здійсненню політики оцінювання ризику та пов'язаних з ними заходів оцінювання ризику.\\ 
b. Призначити [Призначення: визначена організацією посадову особу] для управління політикою та процедурами оцінювання ризику.\\ 
c. Переглядати й оновлювати:\\ 
1. Поточну політику оцінювання організацією частотою]. ризику [Призначення: з визначеною\\ 
2. Поточні процедури оцінювання ризику [Призначення: з визначеною організацією частотою].

\vspace{10pt}
{\footnotesize\bfseries
No: 1\\
Name: ra\_1\_odp\_01\\
Type: list\\
Default: [\textquotedbl{}admin\textquotedbl{}, \textquotedbl{}security\_officer\textquotedbl{}]
}

{\footnotesize Визначено персонал або ролі, на які поширюється політика оцінювання ризику}


{\footnotesize\bfseries
No: 2\\
Name: ra\_1\_odp\_02\\
Type: list\\
Default: [\textquotedbl{}admin\textquotedbl{}, \textquotedbl{}security\_officer\textquotedbl{}]
}

{\footnotesize Визначено персонал або ролі, на які поширюються процедури, що сприяють здійсненню політики оцінювання ризику}


{\footnotesize\bfseries
No: 3\\
Name: ra\_1\_odp\_03\\
Type: string\\
Default: nil
}

{\footnotesize Вибрано одне або декілька з наступних ЗНАЧЕНЬ ПАРАМЕТРІВ: \{рівень організації; рівень місії/бізнеспроцесу; рівень системи\}}


{\footnotesize\bfseries
No: 4\\
Name: ra\_1\_odp\_04\\
Type: list\\
Default: [\textquotedbl{}admin\textquotedbl{}, \textquotedbl{}security\_officer\textquotedbl{}]
}

{\footnotesize Визначена посадова особа, відповідальна за управління політикою та процедурами оцінювання ризику}


{\footnotesize\bfseries
No: 5\\
Name: ra\_1\_odp\_05\\
Type: list\\
Default: [\textquotedbl{}default\_deny\_rule\textquotedbl{}, \textquotedbl{}abac\_rule\_1\textquotedbl{}]
}

{\footnotesize Визначена періодичність перегляду та оновлення поточної політики оцінювання ризику}


{\footnotesize\bfseries
No: 6\\
Name: ra\_1\_odp\_06\\
Type: list\\
Default: [\textquotedbl{}default\_deny\_rule\textquotedbl{}, \textquotedbl{}abac\_rule\_1\textquotedbl{}]
}

{\footnotesize Є події, які вимагають перегляду та оновлення поточної політики оцінювання ризику}


{\footnotesize\bfseries
No: 7\\
Name: ra\_1\_odp\_07\\
Type: string\\
Default: \textquotedbl{}щорічно\textquotedbl{}
}

{\footnotesize Визначено періодичність перегляду поточних процедур оцінювання ризику}




\subsection{ОЦІНЮВАННЯ РИЗИКУ (RA-3)}


\vspace{10pt}
{\footnotesize\bfseries
No: 1\\
Name: ra\_3\_odp\_01\\
Type: string\\
Default: nil
}

{\footnotesize Вибрано одне з наступних ЗНАЧЕНЬ ПАРАМЕТРІВ: \{плани безпеки та приватності; звіт про оцінювання ризику; документ\}}


{\footnotesize\bfseries
No: 2\\
Name: ra\_3\_odp\_02\\
Type: string\\
Default: nil
}

{\footnotesize Визначено документ, в якому мають бути задокументовані результати оцінювання ризику (якщо вони не задокументовані в планах безпеки та приватності або в звіті про оцінювання ризику) (якщо вибрано)}


{\footnotesize\bfseries
No: 3\\
Name: ra\_3\_odp\_03\\
Type: string\\
Default: \textquotedbl{}щорічно\textquotedbl{}
}

{\footnotesize Визначено періодичність оцінювання ризику}


{\footnotesize\bfseries
No: 4\\
Name: ra\_3\_odp\_04\\
Type: list\\
Default: [\textquotedbl{}admin\textquotedbl{}, \textquotedbl{}security\_officer\textquotedbl{}]
}

{\footnotesize Визначено персонал або ролі, до яких мають бути доведені результати оцінювання ризику}


{\footnotesize\bfseries
No: 5\\
Name: ra\_3\_odp\_05\\
Type: list\\
Default: [\textquotedbl{}admin\textquotedbl{}, \textquotedbl{}security\_officer\textquotedbl{}]
}

{\footnotesize Визначена періодичність оновлення оцінювання ризику; перегляду результатів RA-03a.01 проводиться оцінювання ризику для виявлення загроз і вразливостей у системі; RA-03a.02 проводиться оцінювання ризику для визначення ймовірності та розміру шкоди від несанкціонованого доступу, використання, розкриття, порушення, модифікації або знищення системи; інформації, яку вона обробляє, зберігає або передає; а також будь-якої пов'язаної з нею інформації; RA-03a.03 проводиться оцінювання ризику для визначення ймовірності та впливу несприятливих наслідків для фізичних осіб, що виникають у зв'язку з обробкою інформації, яка ідентифікує особу; RA-03b. інтегровані результати оцінювання ризику та рішення з управління ризиками з точки зору організації та місії або бізнес-процесів з оцінкою ризиків на системному рівні; RA-03c. результати оцінювання ризику задокументовані в ЗНАЧЕННЯ ВИБРАНОГО ПАРАМЕТРА; RA-03d. переглядаються результати 03\_ODP[03] частота>; RA-03e. поширюються результати оцінювання ризику серед персоналу або ролей; RA-03f. оновлюється оцінювання ризику частота або коли відбуваються значні зміни в системі, середовищі її функціонування або інших умовах, які можуть вплинути на стан безпеки або приватності системи оцінювання ризику <RA-}




\subsubsection{ОЦІНЮВАННЯ ПОСТАЧАННЯ (RA-3(1))}


\vspace{10pt}
{\footnotesize\bfseries
No: 1\\
Name: ra\_3\_1\_a\\
Type: string\\
Default: nil
}

{\footnotesize Оцінювання ризику ланцюга постачання, пов'язані з системами, системними компонентами та системними послугами}


{\footnotesize\bfseries
No: 2\\
Name: ra\_3\_1\_b\\
Type: string\\
Default: nil
}

{\footnotesize Потрібно оновлювати оцінювання ризику ланцюга постачання, коли відбуваються значні зміни у відповідному ланцюгу постачання, або коли зміни в системі, середовищі функціонування чи інших умовах можуть вимагати змін у ланцюгу постачання}




\subsection{СКАНУВАННЯ ВРАЗЛИВОСТЕЙ (RA-5)}
Використовувати заходи ПДТР за [Призначення: визначені організацією місця] [Вибір (один або кілька): [Призначення: з визначеною організацією частотою]; [Призначення: за визначеними організацією подіями або показниками]].

\vspace{10pt}
{\footnotesize\bfseries
No: 1\\
Name: ra\_5\_odp\_01\\
Type: string\\
Default: nil
}

{\footnotesize Визначена необхідність моніторингу систем та розміщених застосунків на наявність вразливостей}


{\footnotesize\bfseries
No: 2\\
Name: ra\_5\_odp\_02\\
Type: string\\
Default: \textquotedbl{}щорічно\textquotedbl{}
}

{\footnotesize Визначена періодичність перевірки систем та розміщених на них застосунків на наявність вразливостей}


{\footnotesize\bfseries
No: 3\\
Name: ra\_5\_odp\_03\\
Type: integer\\
Default: 30
}

{\footnotesize Визначено час реагування на усунення законних вразливостей відповідно до організаційної оцінення ризику}


{\footnotesize\bfseries
No: 4\\
Name: ra\_5\_odp\_04\\
Type: list\\
Default: [\textquotedbl{}admin\textquotedbl{}, \textquotedbl{}security\_officer\textquotedbl{}]
}

{\footnotesize Потрібно ділитися інформацією, отриманою в процесі сканування вразливостей та оцінок контролю, з персоналом або ролями, з якими потрібно ділитися; RA-05a.[01] здійснюється моніторинг систем та розміщених застосунків на наявність вразливостей частота та/або випадковість відповідно до визначеного організацією процесу, а також коли виявляються та повідомляються нові вразливості, що потенційно можуть вплинути на систему; RA-05a.[02] перевіряються системи та розміщені застосунки на наявність вразливостей частота та/або випадковим чином відповідно до визначеного організацією процесу, а також коли виявляються та повідомляються нові вразливості, що потенційно можуть вплинути на систему; RA-05b. застосовуються інструменти та методи моніторингу вразливостей для забезпечення сумісності між інструментами; RA-05b.01 застосовуються інструменти та методи моніторингу вразливостей для автоматизації частини процесу управління вразливостями, використовуючи стандарти для переліку платформ, недоліків програмного забезпечення та неправильних конфігурацій; RA-05b.02 застосовуються інструменти та методи моніторингу вразливостей для полегшення взаємодії між інструментами та автоматизації частини процесу управління вразливостями шляхом використання стандартів для формування контрольних списків та процедур тестування; RA-05b.03 застосовуються інструменти та методи моніторингу вразливостей для полегшення взаємодії між інструментами та автоматизації частин процесу управління вразливостями шляхом використання стандартів для вимірювання впливу вразливостей; RA-05c. аналізуються звіти про сканування вразливостей та результати моніторингу вразливостей; RA-05d. усуваються легітимні вразливості час реагування відповідно до організаційної оцінки ризиків; RA-05e. надається інформація, отримана в процесі моніторингу вразливостей та оцінки контролю, персонал або ролі, щоб допомогти усунути подібні вразливості в інших системах; RA-05f. використовуються інструменти моніторингу вразливостей, які передбачають можливість швидкого оновлення вразливостей, що підлягають скануванню}




\subsubsection{ОНОВЛЕННЯ ЗА ЧАСТОТОЮ, ПЕРЕД НОВИМ СКАНУВАННЯМ АБО ПРИ ІДЕНТИФІКАЦІЇ (RA-5(2))}


\vspace{10pt}
{\footnotesize\bfseries
No: 1\\
Name: ra\_5\_2\_01\\
Type: integer\\
Default: 30
}

{\footnotesize Визначено час реагування на усунення законних вразливостей відповідно до організаційної оцінення ризику}




\subsection{РЕАГУВАННЯ НА РИЗИК (RA-7)}
Реагувати на результати оцінювання, моніторингу й аудиту безпеки та приватності.

\vspace{10pt}
{\footnotesize\bfseries
No: 1\\
Name: ra\_7\_01\\
Type: string\\
Default: nil
}

{\footnotesize Вживаються заходи реагування на результати оцінок безпеки відповідно до організаційної толерантності до ризиків}


{\footnotesize\bfseries
No: 2\\
Name: ra\_7\_02\\
Type: string\\
Default: nil
}

{\footnotesize Вживаються заходи реагування на результати оцінювання приватності відповідно до організаційної толерантності до ризиків}


{\footnotesize\bfseries
No: 3\\
Name: ra\_7\_03\\
Type: string\\
Default: nil
}

{\footnotesize Вживаються заходи реагування на результати моніторингу відповідно до організаційної толерантності до ризиків}


{\footnotesize\bfseries
No: 4\\
Name: ra\_7\_04\\
Type: string\\
Default: nil
}

{\footnotesize Вживаються заходи реагування на висновки аудиту відповідно до організаційної толерантності до ризиків}




\section{SA}
Клас заходів захисту SA --- ПРИДБАННЯ СИСТЕМИ ТА ПОСЛУГ

Цей клас зосереджується на інтеграції вимог безпеки на всіх етапах життєвого циклу розробки та закупівлі ІТ-продуктів чи послуг.

\textbf{Перелік заходів захисту:} Політики та процедури придбання систем та послуг (SA-1); Безпека та приватність принципів інжинірингу (SA-8); Зовнішні послуги для системи (SA-9); Компоненти системи, що не підтримуються (SA-22).

\subsection{ПОЛІТИКИ ТА ПРОЦЕДУРИ ПРИДБАННЯ СИСТЕМ ТА ПОСЛУГ (SA-1)}
a. Розробити, задокументувати та поширити серед [Призначення: визначеного організацією персоналу або ролей]:\\ 
1. 2.\\ 
b. [Вибір (один або декілька): рівень організації; рівень місії/бізнес-процесу; рівень системи] політики придбання систем і послуг, яка:\\ 
(a) містить мету, сферу застосування, ролі, обов'язки, відповідальність керівництва, координацію між організаційними підрозділами та систему контролю (complaince);\\ 
(b) відповідає чинному законодавству, виконавчим наказам, директивам, нормам, політикам, стандартам і рекомендаціям. Процедури, що полегшують впровадження політики та заходів придбання систем і послуг. Призначити [Призначення: визначену організацією посадову особу] для управління політикою та процедурами придбання системи та послуг.\\ 
c. Переглядати й оновлювати поточні політику та процедури придбання систем та послуг:\\ 
1. Поточну політику придбання системи та послуг [Призначення: з визначеною організацією частотою].\\ 
2. Поточні процедури придбання системи та послуг [Призначення: з визначеною організацією частотою] та наступні [Призначення: події, визначені організацією].

\vspace{10pt}
{\footnotesize\bfseries
No: 1\\
Name: sa\_1\_odp\_01\\
Type: list\\
Default: [\textquotedbl{}admin\textquotedbl{}, \textquotedbl{}security\_officer\textquotedbl{}]
}

{\footnotesize Визначено персонал або ролі, на які поширюватиметься політика придбання систем і послуг}


{\footnotesize\bfseries
No: 2\\
Name: sa\_1\_odp\_02\\
Type: list\\
Default: [\textquotedbl{}admin\textquotedbl{}, \textquotedbl{}security\_officer\textquotedbl{}]
}

{\footnotesize Визначено персонал або ролі, на які поширюються процедури придбання систем і послуг}


{\footnotesize\bfseries
No: 3\\
Name: sa\_1\_odp\_03\\
Type: string\\
Default: nil
}

{\footnotesize Вибрано одне або декілька з наступних ЗНАЧЕНЬ ПАРАМЕТРІВ: \{рівень організації; рівень місії/бізнеспроцесу; рівень системи\}}


{\footnotesize\bfseries
No: 4\\
Name: sa\_1\_odp\_04\\
Type: list\\
Default: [\textquotedbl{}admin\textquotedbl{}, \textquotedbl{}security\_officer\textquotedbl{}]
}

{\footnotesize Визначено посадову особу, яка керуватиме політикою та процедурами придбання систем і послуг}


{\footnotesize\bfseries
No: 5\\
Name: sa\_1\_odp\_05\\
Type: list\\
Default: [\textquotedbl{}default\_deny\_rule\textquotedbl{}, \textquotedbl{}abac\_rule\_1\textquotedbl{}]
}

{\footnotesize Визначено періодичність перегляду та оновлення поточної політики придбання систем і послуг}


{\footnotesize\bfseries
No: 6\\
Name: sa\_1\_odp\_06\\
Type: list\\
Default: [\textquotedbl{}default\_deny\_rule\textquotedbl{}, \textquotedbl{}abac\_rule\_1\textquotedbl{}]
}

{\footnotesize Є події, які вимагають перегляду та оновлення поточної політики придбання систем і послуг}


{\footnotesize\bfseries
No: 7\\
Name: sa\_1\_odp\_07\\
Type: integer\\
Default: 30
}

{\footnotesize Визначено частоту, з якою переглядаються та оновлюються поточні процедури придбання систем і послуг}




\subsection{БЕЗПЕКА ТА ПРИВАТНІСТЬ ПРИНЦИПІВ ІНЖИНІРИНГУ (SA-8)}
Застосовувати [Призначення: визначені організацією принципи інжинірингу безпеки та конфіденційності системи] в специфікації, проєктуванні, розробці, впровадженні та зміні системи й компонентів системи.

\vspace{10pt}
{\footnotesize\bfseries
No: 1\\
Name: sa\_8\_01\\
Type: string\\
Default: nil
}

{\footnotesize Застосовуються принципи інжинірингу безпеки систем у специфікації системи та компонентів системи}


{\footnotesize\bfseries
No: 2\\
Name: sa\_8\_02\\
Type: string\\
Default: nil
}

{\footnotesize Застосовуються принципи інжинірингу безпеки систем при розробці системи та її компонентів}


{\footnotesize\bfseries
No: 3\\
Name: sa\_8\_03\\
Type: string\\
Default: nil
}

{\footnotesize Були застосовані принципи інжинірингу безпеки систем при розробці системи та компонентів систем}


{\footnotesize\bfseries
No: 4\\
Name: sa\_8\_04\\
Type: string\\
Default: nil
}

{\footnotesize Застосовуються принципи інжинірингу безпеки систем при реалізації системи та компонентів систем}


{\footnotesize\bfseries
No: 5\\
Name: sa\_8\_05\\
Type: string\\
Default: nil
}

{\footnotesize Застосовуються принципи інжинірингу безпеки систем при модифікації системи та компонентів систем}


{\footnotesize\bfseries
No: 6\\
Name: sa\_8\_06\\
Type: string\\
Default: nil
}

{\footnotesize Застосовуються принципи інжинірингу конфіденційності у специфікації системи та компонентів систем; принципи інжинірингу конфіденційності SA-08[07] застосовуються принципи інжинірингу конфіденційності при розробці системи та компонентів систем}


{\footnotesize\bfseries
No: 7\\
Name: sa\_8\_08\\
Type: string\\
Default: nil
}

{\footnotesize Застосовуються принципи інжинірингу конфіденційності при розробці системи та компонентів систем}


{\footnotesize\bfseries
No: 8\\
Name: sa\_8\_09\\
Type: string\\
Default: nil
}

{\footnotesize Застосовуються принципи інжинірингу конфіденційності при реалізації системи та компонентів систем}


{\footnotesize\bfseries
No: 9\\
Name: sa\_8\_10\\
Type: string\\
Default: nil
}

{\footnotesize Застосовуються принципи інжинірингу конфіденційності при модифікації системи та компонентів систем}


{\footnotesize\bfseries
No: 10\\
Name: sa\_8\_odp\_01\\
Type: string\\
Default: nil
}

{\footnotesize Визначені принципи інжинірингу безпеки систем}


{\footnotesize\bfseries
No: 11\\
Name: sa\_8\_odp\_02\\
Type: string\\
Default: nil
}

{\footnotesize Визначені системи}




\subsection{ЗОВНІШНІ ПОСЛУГИ ДЛЯ СИСТЕМИ (SA-9)}
a. Вимагати, щоб постачальники зовнішніх послуг для системи відповідали вимогам безпеки та приватності в організації та застосовували такі заходи захисту [Призначення: встановлені організацією заходи безпеки та приватності].\\ 
b. Визначити та задокументувати нагляд організаціїповноваження та обов'язки користувачів щодо зовнішніх послуг для системи.\\ 
c. Використовувати наступні процеси, методи та техніки для постійного моніторингу дотримання контролю зовнішніми постачальниками послуг: [Призначення: визначені організацією процеси, методи та техніки].

\vspace{10pt}
{\footnotesize\bfseries
No: 1\\
Name: sa\_9\_odp\_01\\
Type: string\\
Default: \textquotedbl{}автоматизований засіб моніторингу\textquotedbl{}
}

{\footnotesize Визначені засоби контролю, які будуть застосовуватися зовнішніми постачальниками системних послуг}


{\footnotesize\bfseries
No: 2\\
Name: sa\_9\_odp\_02\\
Type: list\\
Default: [\textquotedbl{}admin\textquotedbl{}, \textquotedbl{}security\_officer\textquotedbl{}]
}

{\footnotesize Визначені процеси, методи та техніки, що застосовуються для моніторингу дотримання вимог контролю зовнішніми постачальниками послуг; SA-09a.[01] дотримуються постачальники зовнішніх системних послуг вимог організаційної безпеки; SA-09a.[02] дотримуються постачальники зовнішніх системних послуг вимог приватності організації; SA-09a.[03] використовують постачальники зовнішніх системних послуг елементи керування; SA-09b.[01] визначено та задокументовано організаційний нагляд за зовнішніми системними послугами; SA-09b.[02] визначені та задокументовані ролі та обов'язки користувачів щодо зовнішніх системних сервісів; SA-09c. визначені та задокументовані ролі та обов'язки користувачів щодо зовнішніх системних послуг}




\subsection{КОМПОНЕНТИ СИСТЕМИ, ЩО НЕ ПІДТРИМУЮТЬСЯ (SA-22)}
a. Замінювати компоненти системи, якщо підтримка компонентів більше не доступна розробнику, постачальнику або виробнику.\\ 
b. Надавати такі варіанти альтернативних джерел для подальшої підтримки непідтримуваних компонентів [Вибір (один або більше): внутрішня підтримка; [Призначення: підтримка, визначена організацією від зовнішніх постачальників]].

\vspace{10pt}
{\footnotesize\bfseries
No: 1\\
Name: sa\_22\_odp\_01\\
Type: string\\
Default: nil
}

{\footnotesize Вибрано одне або декілька з наступних ЗНАЧЕНЬ ПАРАМЕТРІВ: \{внутрішня підтримка; підтримка від зовнішніх постачальників\}}




\section{SC}
Клас заходів захисту SC --- ЗАХИСТ ІНФОРМАЦІЙНОЇ СИСТЕМИ ТА

Цей клас охоплює механізми захисту інформації під час її передачі мережами зв'язку, криптографічний захист та ізоляцію критичних компонентів.

\textbf{Перелік заходів захисту:} Політика та процедури захисту системи та комунікацій (SC-1); Інформація в загальних системних ресурсах (SC-4); Доступність ресурсів (SC-7); Відмова за замовчуванням - дозвіл за винятком (SC-7(5)); Конфіденційність та цілісність передачі (SC-8); Конфіденційність та криптографічний захист (SC-8(1)); Відключення мережі (SC-10); Встановлення ключами (SC-12); Криптографічний захист (SC-13); Спільні обчислювальні пристрої та застосунки (SC-15); Мобільний код (SC-18); Автентифікація сесії (SC-23); Захист інформації в стані спокою (SC-28); ЗАХИСТ ІНФОРМАЦІЇ В СТАНІ СПОКОЮ | КРИПТОГРАФІЧНИЙ ЗАХИСТ.

\subsection{ПОЛІТИКА ТА ПРОЦЕДУРИ ЗАХИСТУ СИСТЕМИ ТА КОМУНІКАЦІЙ (SC-1)}
a. Розробити, задокументувати та поширити серед [Призначення: визначеного організацією персоналу або посадових осіб]:\\ 
1. 2. Політику захисту системи та комунікацій, яка:\\ 
(a) містить мету, сферу застосування, ролі, обов'язки, відповідальність керівництва, координацію між організаційними підрозділами та систему контролю відповідності (complaince);\\ 
(b) відповідає чинному законодавству, виконавчим наказам, директивам, нормам, політикам, стандартам та керівним принципам. Процедури для сприяння впровадженню політики в області захисту систем і комунікацій, а також пов'язаних з ними систем і засобів захисту зв'язку.\\ 
b. Призначити [Призначення: визначена організацією посадову особу] для управління політикою та процедурами захисту системи та комунікацій.\\ 
c. Переглядати та оновлювати:\\ 
1. поточну політику захисту системи та комунікацій [Призначення: визначена організацією частота];\\ 
2. поточні процедури захисту системи та комунікацій [Призначення: визначена організацією частота].

\vspace{10pt}
{\footnotesize\bfseries
No: 1\\
Name: sc\_1\_odp\_01\\
Type: list\\
Default: [\textquotedbl{}admin\textquotedbl{}, \textquotedbl{}security\_officer\textquotedbl{}]
}

{\footnotesize Визначено персонал або ролі, до яких має бути доведена політика захисту системи та комунікацій}


{\footnotesize\bfseries
No: 2\\
Name: sc\_1\_odp\_02\\
Type: list\\
Default: [\textquotedbl{}admin\textquotedbl{}, \textquotedbl{}security\_officer\textquotedbl{}]
}

{\footnotesize Визначено персонал або ролі, на які поширюються процедури захисту системи та комунікацій}


{\footnotesize\bfseries
No: 3\\
Name: sc\_1\_odp\_03\\
Type: string\\
Default: nil
}

{\footnotesize Вибрано жодне або декілька з наступних ЗНАЧЕНЬ ПАРАМЕТРІВ: \{рівень організації; рівень місії/бізнеспроцесу; рівень системи\}}


{\footnotesize\bfseries
No: 4\\
Name: sc\_1\_odp\_04\\
Type: list\\
Default: [\textquotedbl{}admin\textquotedbl{}, \textquotedbl{}security\_officer\textquotedbl{}]
}

{\footnotesize Визначено посадову особу, яка керуватиме політикою та процедурами захисту системи та комунікацій}


{\footnotesize\bfseries
No: 5\\
Name: sc\_1\_odp\_05\\
Type: list\\
Default: [\textquotedbl{}default\_deny\_rule\textquotedbl{}, \textquotedbl{}abac\_rule\_1\textquotedbl{}]
}

{\footnotesize Визначена періодичність перегляду та оновлення поточної політики захисту системи та комунікацій}


{\footnotesize\bfseries
No: 6\\
Name: sc\_1\_odp\_06\\
Type: list\\
Default: [\textquotedbl{}default\_deny\_rule\textquotedbl{}, \textquotedbl{}abac\_rule\_1\textquotedbl{}]
}

{\footnotesize Є події, які вимагають перегляду та оновлення поточної політики захисту системи та комунікацій}


{\footnotesize\bfseries
No: 7\\
Name: sc\_1\_odp\_07\\
Type: string\\
Default: \textquotedbl{}щорічно\textquotedbl{}
}

{\footnotesize Визначена періодичність перегляду та оновлення поточних процедур захисту системи та засобів зв'язку}




\subsection{ІНФОРМАЦІЯ В ЗАГАЛЬНИХ СИСТЕМНИХ РЕСУРСАХ (SC-4)}
Запобігати несанкціонованій та ненавмисній передачі інформації через спільні системні ресурси.

\vspace{10pt}
{\footnotesize\bfseries
No: 1\\
Name: sc\_4\_01\\
Type: string\\
Default: nil
}

{\footnotesize Запобігається несанкціонована передача інформації через спільні системні ресурси}


{\footnotesize\bfseries
No: 2\\
Name: sc\_4\_02\\
Type: string\\
Default: nil
}

{\footnotesize Запобігається ненавмисна передача інформації через спільні системні ресурси. ПОТЕНЦІЙНІ МЕТОДИ ТА ОБ'ЄКТИ ОЦІНКИ:}




\subsection{ДОСТУПНІСТЬ РЕСУРСІВ (SC-7)}
a. Контролювати та управляти зв'язком на зовнішньому периметрі системи та на ключових внутрішніх периметрах всередині системи.\\ 
b. Реалізувати підмережі для загальнодоступних компонентів системи, які є [Вибір: фізично; логічно] відділені від внутрішніх мереж організації.\\ 
c. Підключатися до зовнішніх мереж або систем тільки через керовані інтерфейси, що складаються з пристроїв захисту периметру, і розташованих відповідно до архітектури безпеки та приватності організації.

\vspace{10pt}
\textit{Немає параметрів для цього контролю.}
\vspace{10pt}


\subsubsection{ВІДМОВА ЗА ЗАМОВЧУВАННЯМ - ДОЗВІЛ ЗА ВИНЯТКОМ (SC-7(5))}


\vspace{10pt}
\textit{Немає параметрів для цього контролю.}
\vspace{10pt}


\subsection{КОНФІДЕНЦІЙНІСТЬ ТА ЦІЛІСНІСТЬ ПЕРЕДАЧІ (SC-8)}
Забезпечити [Вибір (один або кілька): конфіденційність; цілісність] інформації, що передається.

\vspace{10pt}
\textit{Немає параметрів для цього контролю.}
\vspace{10pt}


\subsubsection{КОНФІДЕНЦІЙНІСТЬ ТА КРИПТОГРАФІЧНИЙ ЗАХИСТ (SC-8(1))}


\vspace{10pt}
{\footnotesize\bfseries
No: 1\\
Name: sc\_8\_1\_odp\\
Type: string\\
Default: nil
}

{\footnotesize Вибрано одне або декілька з наступних ЗНАЧЕНЬ ПАРАМЕТРІВ: \{запобігти несанкціонованому розголошенню інформації; виявити зміни в інформації\}}




\subsection{ВІДКЛЮЧЕННЯ МЕРЕЖІ (SC-10)}
Завершити з'єднання з мережею, яке пов'язане із сеансом зв'язку в кінці сеансу або після [Призначення: визначений організацією період часу] бездіяльності.

\vspace{10pt}
{\footnotesize\bfseries
No: 1\\
Name: sc\_10\_odp\\
Type: integer\\
Default: 30
}

{\footnotesize Визначено період бездіяльності, після якого система розриває мережеве з'єднання, пов'язане з сеансом зв'язку; SC08(05)\_ODP[02] мережеве з'єднання, пов'язане з сеансом зв'язку, розірвано в кінці сеансу або після період часу бездіяльності}




\subsection{ВСТАНОВЛЕННЯ КЛЮЧАМИ (SC-12)}
Встановити та управляти криптографічними ключами для криптографічних засобів, які використовуються в системі відповідно до [Призначення: визначені організацією вимоги до генерації, поширення, зберігання, доступу та знищення ключів].

\vspace{10pt}
{\footnotesize\bfseries
No: 1\\
Name: sc\_12\_01\\
Type: string\\
Default: \textquotedbl{}AES-256-GCM\textquotedbl{}
}

{\footnotesize Встановлюються криптографічні ключі, коли в системі використовується криптографія відповідно до < SC-12\_ODP вимог >}


{\footnotesize\bfseries
No: 2\\
Name: sc\_12\_02\\
Type: string\\
Default: \textquotedbl{}AES-256-GCM\textquotedbl{}
}

{\footnotesize Здійснюється управління криптографічними ключами, коли в системі використовується криптографія, відповідно до вимог }


{\footnotesize\bfseries
No: 3\\
Name: sc\_12\_odp\\
Type: string\\
Default: nil
}

{\footnotesize Визначені вимоги до генерації, розповсюдження, зберігання, доступу та знищення ключів}




\subsection{КРИПТОГРАФІЧНИЙ ЗАХИСТ (SC-13)}
a. Визначити [Призначення: використання криптографічних засобів, визначених організацією];\\ 
b. Впровадити [Завдання: визначені організацією види криптографії для кожного визначеного криптографічного використання].

\vspace{10pt}
{\footnotesize\bfseries
No: 1\\
Name: sc\_13\_01\\
Type: string\\
Default: \textquotedbl{}AES-256-GCM\textquotedbl{}
}

{\footnotesize КРИПТОГРАФІЧНИЙ ЗАХИСТ МЕТА ОЦІНКИ: Визначити, чи:}


{\footnotesize\bfseries
No: 2\\
Name: sc\_13\_odp\_01\\
Type: string\\
Default: \textquotedbl{}AES-256-GCM\textquotedbl{}
}

{\footnotesize Визначено використання криптографічних засобів}


{\footnotesize\bfseries
No: 3\\
Name: sc\_13\_odp\_02\\
Type: string\\
Default: \textquotedbl{}AES-256-GCM\textquotedbl{}
}

{\footnotesize Визначено типи криптографії для кожного вказаного криптографічного використання; SC-13a. ідентифіковано використання>; SC-13b. реалізовано типи криптографії для кожного вказаного криптографічного використання (визначеного в SC-13\_ODP[01]). криптографічне <SC-13\_ODP[01]}




\subsection{СПІЛЬНІ ОБЧИСЛЮВАЛЬНІ ПРИСТРОЇ ТА ЗАСТОСУНКИ (SC-15)}
a. Заборонити віддалену активацію спільних обчислювальних пристроїв (хмар) та застосунків з такими виключеннями: [Призначення: визначені організацією виключення, у яких дозволена віддалена активація].\\ 
b. Надати явну вказівку щодо використання користувачами фізично присутніми пристроями.

\vspace{10pt}
\textit{Немає параметрів для цього контролю.}
\vspace{10pt}


\subsection{МОБІЛЬНИЙ КОД (SC-18)}
a. Визначати прийнятні та неприйнятні мобільні коди та технології мобільних кодів.\\ 
b. Проводити авторизацію, відстежувати та контролювати використання мобільного коду всередині системи.

\vspace{10pt}
\textit{Немає параметрів для цього контролю.}
\vspace{10pt}


\subsection{АВТЕНТИФІКАЦІЯ СЕСІЇ (SC-23)}
Забезпечити автентифікацію сеансів зв'язку.

\vspace{10pt}
{\footnotesize\bfseries
No: 1\\
Name: sc\_23\_01\\
Type: string\\
Default: nil
}

{\footnotesize Захищено автентифікацію сеансів зв'язку}




\subsection{ЗАХИСТ ІНФОРМАЦІЇ В СТАНІ СПОКОЮ (SC-28)}
Забезпечити [Вибір (один або кілька): конфіденційність; цілісність] [Призначення: визначена організацією інформація] в стані спокою.

\vspace{10pt}
{\footnotesize\bfseries
No: 1\\
Name: sc\_28\_odp\\
Type: string\\
Default: nil
}

{\footnotesize Вибрано одне або декілька з наступних ЗНАЧЕНЬ ПАРАМЕТРІВ: \{конфіденційність; цілісність\}}


{\footnotesize\bfseries
No: 2\\
Name: sc\_28\_odp\_02\\
Type: string\\
Default: nil
}

{\footnotesize Є інформація в стані спокою, яка потребує захисту}




\subsubsection{ЗАХИСТ ІНФОРМАЦІЇ В СТАНІ СПОКОЮ | КРИПТОГРАФІЧНИЙ ЗАХИСТ}


\vspace{10pt}
{\footnotesize\bfseries
No: 1\\
Name: sc\_28\_1\_01\\
Type: string\\
Default: \textquotedbl{}AES-256-GCM\textquotedbl{}
}

{\footnotesize Реалізовані криптографічні механізми для запобігання несанкціонованому розкриттю інформації, що знаходиться в стані спокою на системних компонентах або носіях}


{\footnotesize\bfseries
No: 2\\
Name: sc\_28\_1\_02\\
Type: string\\
Default: \textquotedbl{}AES-256-GCM\textquotedbl{}
}

{\footnotesize Реалізовані криптографічні механізми для запобігання несанкціонованій модифікації інформації, що знаходиться в стані спокою на системних компонентах або носіях}




\section{SI}
Клас заходів захисту SI --- ЦІЛІСНІСТЬ СИСТЕМИ ТА ІНФОРМАЦІЇ

Цей клас спрямований на захист системи від шкідливого коду, виявлення вразливостей та запобігання несанкціонованим змінам інформації.

\textbf{Перелік заходів захисту:} Політика і процедури цілісності інформації (SI-1); Виправлення дефектів (SI-2); Захист від шкідливого коду (SI-3); Моніторинг системи (SI-4); Моніторинг системи комунікацій (SI-4(4)); Попередження, рекомендації та директиви з безпеки (SI-5); Управління та збереження інформації (SI-12).

\subsection{ПОЛІТИКА І ПРОЦЕДУРИ ЦІЛІСНОСТІ ІНФОРМАЦІЇ (SI-1)}


\vspace{10pt}
{\footnotesize\bfseries
No: 1\\
Name: si\_1\_odp\_01\\
Type: list\\
Default: [\textquotedbl{}admin\textquotedbl{}, \textquotedbl{}security\_officer\textquotedbl{}]
}

{\footnotesize Визначено персонал або ролі, до яких має бути доведена політика цілісності системи та інформації}


{\footnotesize\bfseries
No: 2\\
Name: si\_1\_odp\_02\\
Type: list\\
Default: [\textquotedbl{}admin\textquotedbl{}, \textquotedbl{}security\_officer\textquotedbl{}]
}

{\footnotesize Визначено персонал або ролі, на які поширюються процедури цілісності системи та інформації }


{\footnotesize\bfseries
No: 3\\
Name: si\_1\_odp\_03\\
Type: string\\
Default: nil
}

{\footnotesize Вибрано одне або декілька з наступних ЗНАЧЕНЬ ПАРАМЕТРІВ: \{рівень організації; рівень місії/бізнеспроцесу; рівень системи\}}


{\footnotesize\bfseries
No: 4\\
Name: si\_1\_odp\_04\\
Type: list\\
Default: [\textquotedbl{}admin\textquotedbl{}, \textquotedbl{}security\_officer\textquotedbl{}]
}

{\footnotesize Визначено посадову особу, відповідальну за управління системою та політикою і процедурами цілісності інформації}


{\footnotesize\bfseries
No: 5\\
Name: si\_1\_odp\_05\\
Type: list\\
Default: [\textquotedbl{}default\_deny\_rule\textquotedbl{}, \textquotedbl{}abac\_rule\_1\textquotedbl{}]
}

{\footnotesize Визначено періодичність перегляду та оновлення поточної політики цілісності системи та інформації}


{\footnotesize\bfseries
No: 6\\
Name: si\_1\_odp\_06\\
Type: list\\
Default: [\textquotedbl{}default\_deny\_rule\textquotedbl{}, \textquotedbl{}abac\_rule\_1\textquotedbl{}]
}

{\footnotesize Є події, які вимагають перегляду та оновлення поточної політики цілісності системи та інформації}


{\footnotesize\bfseries
No: 7\\
Name: si\_1\_odp\_07\\
Type: integer\\
Default: 30
}

{\footnotesize Визначено частоту, з якою переглядаються та оновлюються поточні цілісності системи та інформації}




\subsection{ВИПРАВЛЕННЯ ДЕФЕКТІВ (SI-2)}


\vspace{10pt}
{\footnotesize\bfseries
No: 1\\
Name: si\_2\_odp\\
Type: integer\\
Default: 30
}

{\footnotesize Визначено період часу, протягом якого необхідно встановити оновлення програмного забезпечення, пов'язані з безпекою, після виходу оновлень; SI-02a.[01] виявлено недоліки системи; SI-02a.[02] повідомляється про недоліки системи; SI-02a.[03] виправлені недоліки системи; SI-02b.[01] перевіряються оновлення програмного забезпечення, пов'язані з усуненням недоліків, на ефективність перед встановленням; SI-02b.[02] перевіряються оновлення програмного забезпечення, пов'язані з виправленням дефектів, на наявність потенційних побічних ефектів перед встановленням; SI-02b.[03] перевіряються оновлення прошивки, пов'язані з усуненням недоліків, на ефективність перед встановленням; SI-02b.[04] перевіряються оновлення прошивки, пов'язані з усуненням недоліків, на наявність потенційних побічних ефектів перед встановленням; SI-02c.[01] встановлено оновлення програмного забезпечення, що стосуються безпеки, протягом часовий проміжок з моменту випуску оновлень; SI-02c.[02] встановлено оновлення мікропрограми, що стосуються безпеки, протягом часового періоду з моменту випуску оновлень; SI-02d. включено відновлення порушених прав у процес управління організаційною конфігурацією}




\subsection{ЗАХИСТ ВІД ШКІДЛИВОГО КОДУ (SI-3)}


\vspace{10pt}
{\footnotesize\bfseries
No: 1\\
Name: si\_3\_odp\_01\\
Type: string\\
Default: nil
}

{\footnotesize Вибрано одне або декілька з наступних ЗНАЧЕНЬ ПАРАМЕТРІВ: \{підписаний; непідписаний\}}


{\footnotesize\bfseries
No: 2\\
Name: si\_3\_odp\_02\\
Type: integer\\
Default: 30
}

{\footnotesize Визначено частоту, з якою механізми шкідливого коду виконують сканування}


{\footnotesize\bfseries
No: 3\\
Name: si\_3\_odp\_03\\
Type: string\\
Default: nil
}

{\footnotesize Вибрано одне або декілька з наступних ЗНАЧЕНЬ ПАРАМЕТРІВ: \{кінцева точка; точки входу та виходу з мережі\}}


{\footnotesize\bfseries
No: 4\\
Name: si\_3\_odp\_04\\
Type: string\\
Default: nil
}

{\footnotesize Вибрано одне або декілька з наступних ЗНАЧЕНЬ ПАРАМЕТРІВ: \{block malicious code; quarantitne malicious code; take action\}}


{\footnotesize\bfseries
No: 5\\
Name: si\_3\_odp\_05\\
Type: list\\
Default: [\textquotedbl{}login\textquotedbl{}, \textquotedbl{}logout\textquotedbl{}, \textquotedbl{}failed\_attempt\textquotedbl{}]
}

{\footnotesize Визначено дії, яких слід вжити у відповідь на виявлення шкідливого коду (якщо вибрано)}




\subsection{МОНІТОРИНГ СИСТЕМИ (SI-4)}


\vspace{10pt}
{\footnotesize\bfseries
No: 1\\
Name: si\_4\_odp\_02\\
Type: list\\
Default: [\textquotedbl{}admin\textquotedbl{}, \textquotedbl{}security\_officer\textquotedbl{}]
}

{\footnotesize Визначені методи та способи, що використовуються для виявлення несанкціонованого використання системи}


{\footnotesize\bfseries
No: 2\\
Name: si\_4\_odp\_03\\
Type: list\\
Default: [\textquotedbl{}admin\textquotedbl{}, \textquotedbl{}security\_officer\textquotedbl{}]
}

{\footnotesize Визначена інформація про моніторинг системи, яка повинна надаватися персоналу або ролям}


{\footnotesize\bfseries
No: 3\\
Name: si\_4\_odp\_04\\
Type: list\\
Default: [\textquotedbl{}admin\textquotedbl{}, \textquotedbl{}security\_officer\textquotedbl{}]
}

{\footnotesize Визначено персонал або ролі, яким має надаватися інформація про моніторинг системи}


{\footnotesize\bfseries
No: 4\\
Name: si\_4\_odp\_05\\
Type: string\\
Default: \textquotedbl{}щорічно\textquotedbl{}
}

{\footnotesize Вибрано одне або декілька з наступних ЗНАЧЕНЬ ПАРАМЕТРІВ: \{за потребою; частота\}}




\subsubsection{МОНІТОРИНГ СИСТЕМИ КОМУНІКАЦІЙ (SI-4(4))}


\vspace{10pt}
{\footnotesize\bfseries
No: 1\\
Name: si\_4\_4\_a\_01\\
Type: list\\
Default: []
}

{\footnotesize Визначені критерії незвичної або несанкціонованої діяльності або умови для вхідного трафіку зв'язку}


{\footnotesize\bfseries
No: 2\\
Name: si\_4\_4\_a\_02\\
Type: list\\
Default: []
}

{\footnotesize Визначені критерії незвичайної або несанкціонованої діяльності або умови для вихідного трафіку зв'язку}


{\footnotesize\bfseries
No: 3\\
Name: si\_4\_4\_b\_01\\
Type: string\\
Default: \textquotedbl{}щорічно\textquotedbl{}
}

{\footnotesize Здійснюється моніторинг вхідного комунікаційного трафіку частота на предмет незвичних або несанкціонованих дій або умов}


{\footnotesize\bfseries
No: 4\\
Name: si\_4\_4\_b\_02\\
Type: string\\
Default: \textquotedbl{}щорічно\textquotedbl{}
}

{\footnotesize Контролюється вихідний трафік зв'язку частота на предмет незвичних або несанкціонованих дій або умов. вихідного виявлення}




\subsection{ПОПЕРЕДЖЕННЯ, РЕКОМЕНДАЦІЇ ТА ДИРЕКТИВИ З БЕЗПЕКИ (SI-5)}


\vspace{10pt}
{\footnotesize\bfseries
No: 1\\
Name: si\_5\_odp\_01\\
Type: string\\
Default: nil
}

{\footnotesize Визначені зовнішні організації, від яких необхідно постійно отримувати оповіщення, поради та директиви щодо безпеки системи}


{\footnotesize\bfseries
No: 2\\
Name: si\_5\_odp\_02\\
Type: list\\
Default: [\textquotedbl{}admin\textquotedbl{}, \textquotedbl{}security\_officer\textquotedbl{}]
}

{\footnotesize Вибрано одне або декілька з наступних ЗНАЧЕНЬ ПАРАМЕТРІВ: \{персонал або ролі; елементи; зовнішні організації\}}


{\footnotesize\bfseries
No: 3\\
Name: si\_5\_odp\_03\\
Type: list\\
Default: [\textquotedbl{}admin\textquotedbl{}, \textquotedbl{}security\_officer\textquotedbl{}]
}

{\footnotesize Визначено персонал або ролі, до яких мають бути доведені попередження, поради та директиви з безпеки (якщо визначено)}


{\footnotesize\bfseries
No: 4\\
Name: si\_5\_odp\_04\\
Type: string\\
Default: nil
}

{\footnotesize Визначені елементи в організації, до яких мають надсилатися оповіщення, поради та директиви з безпеки (якщо вони були обрані)}




\subsection{УПРАВЛІННЯ ТА ЗБЕРЕЖЕННЯ ІНФОРМАЦІЇ (SI-12)}
Управляти та зберігати інформацію всередині системи та виводити інформацію із системи відповідно до чинного законодавства, виконавчих наказів, директив, правил, політик, стандартів, керівних принципів та експлуатаційних вимог.

\vspace{10pt}
{\footnotesize\bfseries
No: 1\\
Name: si\_12\_01\\
Type: list\\
Default: [\textquotedbl{}default\_deny\_rule\textquotedbl{}, \textquotedbl{}abac\_rule\_1\textquotedbl{}]
}

{\footnotesize Здійснюється управління інформацією в системі відповідно до чинних законів, наказів, директив, положень, політик, стандартів, інструкцій та експлуатаційних вимог}


{\footnotesize\bfseries
No: 2\\
Name: si\_12\_02\\
Type: list\\
Default: [\textquotedbl{}default\_deny\_rule\textquotedbl{}, \textquotedbl{}abac\_rule\_1\textquotedbl{}]
}

{\footnotesize Зберігається інформація в системі відповідно до чинних законів, указів Президента, директив, положень, політик, стандартів, інструкцій та експлуатаційних вимог}


{\footnotesize\bfseries
No: 3\\
Name: si\_12\_03\\
Type: list\\
Default: [\textquotedbl{}default\_deny\_rule\textquotedbl{}, \textquotedbl{}abac\_rule\_1\textquotedbl{}]
}

{\footnotesize Управління інформацією, що виводиться з системи, здійснюється відповідно до чинних законів, указів Президента, директив, положень, політик, стандартів, інструкцій та операційних вимог}


{\footnotesize\bfseries
No: 4\\
Name: si\_12\_04\\
Type: list\\
Default: [\textquotedbl{}default\_deny\_rule\textquotedbl{}, \textquotedbl{}abac\_rule\_1\textquotedbl{}]
}

{\footnotesize Зберігається інформація, що виводиться з системи, відповідно до чинних законів, наказів, директив, положень, політик, стандартів, інструкцій та експлуатаційних вимог}




\section{SR}
Клас заходів захисту SR --- УПРАВЛІННЯ РИЗИКАМИ ЛАНЦЮГА

Цей клас встановлює вимоги для виявлення та мінімізації загроз, що виникають через зовнішніх постачальників продуктів та послуг.

\textbf{Перелік заходів захисту:} Політика та процедури управління ризиками ланцюга постачання (SR-1); План управління ризиками ланцюга постачання (SR-2); Контроль ланцюга постачання і процесів (SR-3); Стратегії придбання, інструменти і методи (SR-5); Оцінка постачальників (SR-6).

\subsection{ПОЛІТИКА ТА ПРОЦЕДУРИ УПРАВЛІННЯ РИЗИКАМИ ЛАНЦЮГА ПОСТАЧАННЯ (SR-1)}
a. Розробіть, задокументуйте та поширте [Призначення: персонал або ролі, визначені організацією]:\\ 
1. [Вибір (один або декілька): Рівень організації; Рівень місії/бізнес-процесу; рівень системи] політика управління ризиками ланцюга постачання, яка: a) Розглядає мету, сферу діяльності, ролі, відповідальність, зобов'язання керівництва, координацію між організаційними підрозділами та відповідність; b) Відповідає чинним законам, виконавчим наказам, директивам, положенням, політикам, стандартам і вказівкам;\\ 
2. Процедури для сприяння впровадженню політики управління ризиками ланцюга постачання та відповідних засобів контролю управління ризиками ланцюга постачання;\\ 
b. Призначити [Призначення: посадова особа, визначена організацією] для управління розробкою, документуванням і розповсюдженням політики та процедур управління ризиками ланцюга постачання;\\ 
c. Перегляньте та оновіть поточне управління ризиками ланцюга постачання:\\ 
1. Політика [Призначення: частота, визначена організацією] та наступне [Призначення: події, визначені організацією];\\ 
2. Процедури [Призначення: частота, визначена організацією] та наступні [Призначення: події, визначені організацією].

\vspace{10pt}
{\footnotesize\bfseries
No: 1\\
Name: sr\_1\_odp\_01\\
Type: list\\
Default: [\textquotedbl{}admin\textquotedbl{}, \textquotedbl{}security\_officer\textquotedbl{}]
}

{\footnotesize Визначено персонал або ролі, на які поширюється політика управління ризиками ланцюга постачання}


{\footnotesize\bfseries
No: 2\\
Name: sr\_1\_odp\_02\\
Type: list\\
Default: [\textquotedbl{}admin\textquotedbl{}, \textquotedbl{}security\_officer\textquotedbl{}]
}

{\footnotesize Визначено персонал або ролі, на які поширюються процедури управління ризиками ланцюга постачання}


{\footnotesize\bfseries
No: 3\\
Name: sr\_1\_odp\_03\\
Type: string\\
Default: nil
}

{\footnotesize Вибрано одне або декілька з наступних ЗНАЧЕНЬ ПАРАМЕТРІВ: \{рівень організації; рівень завдань/бізнеспроцесу; рівень системи\}}


{\footnotesize\bfseries
No: 4\\
Name: sr\_1\_odp\_04\\
Type: list\\
Default: [\textquotedbl{}admin\textquotedbl{}, \textquotedbl{}security\_officer\textquotedbl{}]
}

{\footnotesize Визначена посадова особа, відповідальна за розробку, документування та розповсюдження політики та процедур управління ризиками ланцюга постачання}


{\footnotesize\bfseries
No: 5\\
Name: sr\_1\_odp\_05\\
Type: list\\
Default: [\textquotedbl{}default\_deny\_rule\textquotedbl{}, \textquotedbl{}abac\_rule\_1\textquotedbl{}]
}

{\footnotesize Визначена періодичність перегляду та оновлення поточної політики управління ризиками ланцюга постачання}


{\footnotesize\bfseries
No: 6\\
Name: sr\_1\_odp\_06\\
Type: list\\
Default: [\textquotedbl{}default\_deny\_rule\textquotedbl{}, \textquotedbl{}abac\_rule\_1\textquotedbl{}]
}

{\footnotesize Є події, які вимагають перегляду та оновлення поточної політики управління ризиками ланцюга постачання}


{\footnotesize\bfseries
No: 7\\
Name: sr\_1\_odp\_07\\
Type: string\\
Default: \textquotedbl{}щорічно\textquotedbl{}
}

{\footnotesize Визначена періодичність перегляду та оновлення поточної процедури управління ризиками ланцюга постачання}




\subsection{ПЛАН УПРАВЛІННЯ РИЗИКАМИ ЛАНЦЮГА ПОСТАЧАННЯ (SR-2)}
a. Розробіть план управління ризиками ланцюга постачання, пов'язаними з дослідженнями та розробкою, проектуванням, виробництвом, придбанням, доставкою, інтеграцією, експлуатацією та обслуговуванням, а також утилізацією таких систем, компонентів системи або послуг для системи: [Призначення: системи, визначені організацією, системні компоненти або системні служби];\\ 
b. Перегляньте та оновіть план управління ризиками ланцюга постачання [Призначення: частота, визначена організацією] або за потреби для усунення загроз;\\ 
c. Захистіть план управління ризиками ланцюга постачання від несанкціонованого розголошення та модифікації.

\vspace{10pt}
{\footnotesize\bfseries
No: 1\\
Name: sr\_2\_odp\_01\\
Type: string\\
Default: nil
}

{\footnotesize Визначені системи, компоненти системи або системні послуги, для яких розробляється план управління ризиками ланцюга постачання}




\subsection{КОНТРОЛЬ ЛАНЦЮГА ПОСТАЧАННЯ І ПРОЦЕСІВ (SR-3)}
a. Встановлення процесу або процесів для виявлення та усунення слабких місць або недоліків в елементах і процесах ланцюга постачання [Призначення: визначена організацією система або компонент системи] у координації з [Завдання: персонал ланцюга постачання, визначений організацією];\\ 
b. Використовуйте такі заходи захисту, щоб захистити систему, компонент системи або системну службу від ризиків ланцюга постачання та обмежити шкоду чи наслідки від подій, пов'язаних із ланцюгом постачання: [Призначення: заходи захисту ланцюга постачання, визначені організацією];\\ 
c. Задокументуйте обрані та впроваджені процеси та заходи захисту ланцюгом постачання у [Вибір: плани безпеки та приватності; план управління ризиками ланцюга постачання; [Призначення: документ, визначений організацією]].

\vspace{10pt}
{\footnotesize\bfseries
No: 1\\
Name: sr\_3\_odp\_01\\
Type: string\\
Default: nil
}

{\footnotesize Визначено систему або компонент системи, який потребує процесу або процесів для виявлення та усунення слабких місць або недоліків}


{\footnotesize\bfseries
No: 2\\
Name: sr\_3\_odp\_02\\
Type: list\\
Default: [\textquotedbl{}admin\textquotedbl{}, \textquotedbl{}security\_officer\textquotedbl{}]
}

{\footnotesize Визначено персонал ланцюга поставок, з яким необхідно координувати процес або процеси виявлення та усунення слабких місць або недоліків в елементах і процесах ланцюга постачання}


{\footnotesize\bfseries
No: 3\\
Name: sr\_3\_odp\_03\\
Type: string\\
Default: \textquotedbl{}автоматизований засіб моніторингу\textquotedbl{}
}

{\footnotesize Визначені засоби контролю ланцюга постачання, що застосовуються для захисту від ризиків ланцюга постачання для системи, системного компонента або системної послуги, а також для обмеження шкоди або наслідків від подій, пов'язаних з ланцюгом постачання}


{\footnotesize\bfseries
No: 4\\
Name: sr\_3\_odp\_04\\
Type: list\\
Default: [\textquotedbl{}admin\textquotedbl{}, \textquotedbl{}security\_officer\textquotedbl{}]
}

{\footnotesize Вибрано одне або декілька з наступних ЗНАЧЕНЬ ПАРАМЕТРІВ: \{плани безпеки та конфіденційності; план управління ризиками ланцюга постачання; документ\}; SR-03\_ODP[05] визначено документ, що ідентифікує обрані та впроваджені процеси та засоби контролю ланцюга постачання (якщо обрано); SR-03a.[01] запроваджено процес або процеси для виявлення та усунення слабких місць або недоліків в елементах та процесах ланцюга постачання; SR-03a.[02] процес або процеси виявлення та усунення слабких місць або недоліків в елементах та процесах ланцюга постачання системи або компонента системи координується/координуються з персоналом ланцюга постачання; SR-03b. застосовуються засоби контролю ланцюга постачання для захисту від ризиків ланцюга постачання для системи, системного компонента або системної послуги, а також для обмеження шкоди або наслідків від подій, пов'язаних з ланцюгом постачання; SR-03c. задокументовані обрані та впроваджені процеси та засоби контролю ланцюга постачання в ЗНАЧЕННЯ ВИБІРКОВОГО ПАРАМЕТРА(ів)}


{\footnotesize\bfseries
No: 5\\
Name: sr\_3\_odp\_05\\
Type: list\\
Default: [\textquotedbl{}admin\textquotedbl{}, \textquotedbl{}security\_officer\textquotedbl{}]
}

{\footnotesize Визначено документ, що ідентифікує обрані та впроваджені процеси та засоби контролю ланцюга постачання (якщо обрано); SR-03a.[01] запроваджено процес або процеси для виявлення та усунення слабких місць або недоліків в елементах та процесах ланцюга постачання; SR-03a.[02] процес або процеси виявлення та усунення слабких місць або недоліків в елементах та процесах ланцюга постачання системи або компонента системи координується/координуються з персоналом ланцюга постачання; SR-03b. застосовуються засоби контролю ланцюга постачання для захисту від ризиків ланцюга постачання для системи, системного компонента або системної послуги, а також для обмеження шкоди або наслідків від подій, пов'язаних з ланцюгом постачання; SR-03c. задокументовані обрані та впроваджені процеси та засоби контролю ланцюга постачання в ЗНАЧЕННЯ ВИБІРКОВОГО ПАРАМЕТРА(ів)}




\subsection{СТРАТЕГІЇ ПРИДБАННЯ, ІНСТРУМЕНТИ І МЕТОДИ (SR-5)}
Використовуйте наступні стратегії придбання, контрактні інструменти та методи закупівель, щоб захистити від ризиків ланцюга постачання, визначити та пом'якшити їх: [Призначення: визначені організацією стратегії придбання, контрактні інструменти та методи закупівель].

\vspace{10pt}
{\footnotesize\bfseries
No: 1\\
Name: sr\_5\_01\\
Type: string\\
Default: nil
}

{\footnotesize Застосовуються стратегії, інструменти та методи для захисту від ризиків ланцюга постачання}


{\footnotesize\bfseries
No: 2\\
Name: sr\_5\_02\\
Type: string\\
Default: nil
}

{\footnotesize Застосовуються стратегії, інструменти та методи для виявлення ризиків ланцюга постачання}


{\footnotesize\bfseries
No: 3\\
Name: sr\_5\_03\\
Type: string\\
Default: nil
}

{\footnotesize Застосовуються стратегії, інструменти та методи для зменшення ризиків ланцюга постачання}


{\footnotesize\bfseries
No: 4\\
Name: sr\_5\_odp\\
Type: string\\
Default: nil
}

{\footnotesize Визначені стратегії закупівель, контрактні інструменти та методи закупівель для захисту, виявлення та пом'якшення ризиків ланцюга постачання}




\subsection{ОЦІНКА ПОСТАЧАЛЬНИКІВ (SR-6)}
Оцініть і перегляньте ризики ланцюга постачання, пов'язані з постачальниками або підрядниками, системою, системним компонентом або системною послугою, яку вони надають [Призначення: частота, визначена організацією].

\vspace{10pt}
{\footnotesize\bfseries
No: 1\\
Name: sr\_6\_01\\
Type: string\\
Default: \textquotedbl{}щорічно\textquotedbl{}
}

{\footnotesize Оцінюються та аналізуються ризики, пов'язані з ланцюгом постачання, які стосуються постачальників або підрядників та систем, компонентів системи або системних послуг, які вони надають < SR-06\_ODP частота >}


{\footnotesize\bfseries
No: 2\\
Name: sr\_6\_odp\\
Type: string\\
Default: \textquotedbl{}щорічно\textquotedbl{}
}

{\footnotesize Визначена періодичність оцінки та аналізу ризиків, пов'язаних з ланцюгом постачання, що стосуються постачальників або підрядників, а також систем, компонентів системи або системних послуг, які вони надають}




\end{document}
